% Consolidation des acquis de la méthodologie de l’exercice sur texte philosophique — Philosophie Tle Tech — Cours signé OnBuch+
% Programme officiel MINESEC. Compiler avec : tectonic N17-consolidation-acquis-methodologie-exercice-texte-philos.tex
\newcommand{\DOCMATIERE}{Philosophie}
\newcommand{\DOCNIVEAU}{Tle Tech}
\newcommand{\DOCMODULE}{Philosophie – classe de Terminale technique}
\newcommand{\DOCLECON}{N17}
\newcommand{\DOCTITRE}{Consolidation des acquis de la méthodologie de l’exercice sur texte philosophique}
% OnBuch+ — préambule « cours signés » ENRICHI (compilé avec tectonic / XeLaTeX)
% Système visuel « Pop » d'OnBuch+ : fond crème (jamais blanc), bordures encre
% épaisses, ombres « dures » décalées, titres Archivo Black en capitales.
% Version MATHÉMATIQUES : graphes (pgfplots/tikz), boîtes pédagogiques
% (définition, propriété, méthode, attention, le savais-tu, exercice résolu,
% exercice, TP, bilan), page de garde et signature OnBuch+ ancrée en bas.
%
% Macros attendues dans main.tex AVANT \input{preamble} :
%   \DOCMATIERE  \DOCNIVEAU  \DOCMODULE  \DOCLECON (numéro)  \DOCTITRE
\documentclass[11pt]{article}

\usepackage{fontspec}
\usepackage[french]{babel}
\usepackage[a4paper,margin=2cm,top=3.4cm,bottom=2.8cm,headheight=1.4cm,footskip=1.3cm]{geometry}
\usepackage{amsmath,amssymb,amsfonts}
\usepackage{mathtools} % \xmapsto, \coloneqq... souvent utilisés par les modèles
\usepackage{cancel} % simplifications barrées (bilans de forces)
% \abs{z} (module, valeur absolue) et \norm{u} (norme), utilisés par les modèles
\providecommand{\abs}{}\let\abs\relax\DeclarePairedDelimiter{\abs}{\lvert}{\rvert}
\providecommand{\norm}{}\let\norm\relax\DeclarePairedDelimiter{\norm}{\lVert}{\rVert}
\usepackage[table]{xcolor} % [table] : fournit aussi \rowcolors (alternance de lignes)
\usepackage{pagecolor}
\usepackage{tikz}
\usetikzlibrary{arrows.meta,positioning,calc,shapes.geometric,shapes.misc,shapes.arrows,decorations.pathmorphing,decorations.pathreplacing,decorations.markings,patterns,shadows,mindmap,trees,fit,backgrounds,matrix,intersections,angles,quotes}
\usepackage{pgfplots}
\usepackage[version=4]{mhchem} % équations nucléaires \ce{^{235}_{92}U}
\usepackage[european,siunitx]{circuitikz} % schémas électriques
\pgfplotsset{compat=1.18}
\usepgfplotslibrary{fillbetween,groupplots} % aires entre courbes (calcul intégral)
\usepackage{enumitem}
\usepackage{fancyhdr}
% [most] charge toutes les bibliothèques tcolorbox (listings, minted, raster,
% documentation...) et gonfle inutilement la mémoire TeX sur un document long
% (~300 boîtes) : on ne charge que ce qui est réellement utilisé.
\usepackage[skins,breakable]{tcolorbox}
\usepackage{graphicx}
\usepackage{colortbl}
\usepackage{booktabs}
\usepackage{tabularx}
\usepackage{multirow}
\usepackage{diagbox} % case d'en-tête diagonale des tableaux à double entrée
% Colonnes extensibles centrée / gauche / droite (C, L, R), souvent utilisées par les modèles
\newcolumntype{C}{>{\centering\arraybackslash}X}
\newcolumntype{L}{>{\raggedright\arraybackslash}X}
\newcolumntype{R}{>{\raggedleft\arraybackslash}X}
\usepackage{array}
\usepackage{multicol}
\usepackage{siunitx}
% parse-numbers=false : accepte aussi \SI{2,5 \times 10^{-3}}{...}
\sisetup{locale=FR,output-decimal-marker={,},group-minimum-digits=4,parse-numbers=false}
\DeclareSIUnit\ohm{\ensuremath{\symup{\Omega}}} % U+2126 absent de la police
\DeclareSIUnit\division{div} % réglages d'oscilloscope (ms/div, V/div)
\DeclareSIUnit\tr{tr} % tours (vitesse de rotation, ex. \SI{1500}{\tr\per\minute})
\DeclareSIUnit\molar{M} % concentration molaire (ex. \SI{3}{\molar}, \SI{1}{\micro\molar})
\DeclareSIUnit\an{an} % année (le modèle confond parfois avec \year, un compteur TeX) — ex. \SI{5}{\kilo\meter\per\an}
\DeclareSIUnit\FCFA{FCFA} % franc CFA (ex. \SI{500}{\FCFA\per\kilo\gram})
\DeclareSIUnit\degreeBrix{\textdegree Brix} % teneur en sucre d'un sirop/jus (réfractométrie)
\DeclareSIUnit\month{mois} % le modèle utilise parfois \month, un compteur TeX, comme unité de durée
\DeclareSIUnit\microliter{\micro\liter} % le modèle écrit parfois \microliter en un seul token
\DeclareSIUnit\million{millions} % dénombrement (ex. \SI{15}{\million\per\mL} spermatozoïdes)
\usepackage{qrcode}
% \degree : commande fréquemment utilisée par les modèles pour un angle ou une
% température en dehors de \SI{}{} (non fournie par les paquets chargés).
\providecommand{\degree}{^{\circ}}
% \qed (fin de démonstration), utilisé par les modèles sans amsthm
\providecommand{\qed}{\hfill$\square$}
% \barre : double barre des valeurs interdites dans un tableau de variations (tkz-tab)
\providecommand{\barre}{\ensuremath{\|}}
% environnement proof (amsthm non chargé) utilisé par les modèles
\ifdefined\proof\else\newenvironment{proof}[1][Démonstration]{\par\noindent\textit{#1.}\ }{\qed\par}\fi
\usepackage{lastpage}
\usepackage{hyperref}
\hypersetup{hidelinks}

% ============================================================
% Palette « Pop » OnBuch+ — mêmes hex que la classe P (plus_ui.dart)
% ============================================================
\definecolor{poporange}{HTML}{F59321}
\definecolor{popdark}{HTML}{E07A0C}
\definecolor{popcream}{HTML}{FFF6E8}
\definecolor{popink}{HTML}{1C1714}
\definecolor{poppurple}{HTML}{7A5AE0}
\definecolor{popgreen}{HTML}{2E9E6A}
\definecolor{poppink}{HTML}{E0457B}
\definecolor{popblue}{HTML}{2F7DE1}
\definecolor{popgold}{HTML}{C98A12}
\definecolor{popmuted}{HTML}{8A7F76}
\definecolor{popline}{HTML}{EADFD2}
\definecolor{popgray}{HTML}{8A7F76}
\colorlet{popgrey}{popgray}
% Alias tolérants (le modèle nomme parfois les couleurs différemment)
\colorlet{popwhite}{white}
\colorlet{popblack}{popink}
\colorlet{popred}{poppink}
\colorlet{popyellow}{popgold}
% Teintes claires (fonds de tableaux/encadrés) — le modèle nomme parfois
% les couleurs avec un suffixe L (ex. popblueL) sans les avoir définies.
\colorlet{poporangeL}{poporange!20}
\colorlet{popdarkL}{popdark!20}
\colorlet{poppurpleL}{poppurple!20}
\colorlet{popgreenL}{popgreen!20}
\colorlet{poppinkL}{poppink!20}
\colorlet{popblueL}{popblue!20}
\colorlet{popgoldL}{popgold!20}
\colorlet{popredL}{popred!20}
\colorlet{popgrayL}{popgray!20}
% Teintes claires pour fonds de tableaux / zones
\colorlet{popblueL}{popblue!10!white}
\colorlet{poporangeL}{poporange!14!white}
\colorlet{popgreenL}{popgreen!12!white}
\colorlet{poppurpleL}{poppurple!12!white}
\colorlet{poppinkL}{poppink!12!white}
\colorlet{popgoldL}{popgold!14!white}

\pagecolor{popcream}

% ============================================================
% Typographie
% ============================================================
\defaultfontfeatures{Ligatures=TeX}
\setmainfont{PlusJakartaSans-Medium.ttf}[Path=../../../fonts/,BoldFont=PlusJakartaSans-Bold.ttf]
% Maths en sans-serif (Fira Math) avec chiffres et lettres droites de Plus
% Jakarta, pour que les formules se fondent dans le texte.
\usepackage[math-style=ISO,bold-style=ISO]{unicode-math}
\setmathfont{FiraMath-Regular.otf}[Path=../../../fonts/]
\setmathfont{PlusJakartaSans-Medium.ttf}[Path=../../../fonts/,range={up/num,up/{Latin,latin},\mathup}]
\usepackage{icomma} % virgule décimale sans espace en mode math
% Caractères mathématiques Unicode absents des polices de texte (Plus Jakarta,
% Archivo Black) : rendus par leur équivalent mathématique, en texte comme en
% formule — sinon ils s'affichent en carré vide (ex. « Arithmétique dans ℤ »).
\usepackage{newunicodechar}
\newunicodechar{ℝ}{\ensuremath{\mathbb{R}}}
\newunicodechar{ℤ}{\ensuremath{\mathbb{Z}}}
\newunicodechar{ℂ}{\ensuremath{\mathbb{C}}}
\newunicodechar{ℕ}{\ensuremath{\mathbb{N}}}
\newunicodechar{ℚ}{\ensuremath{\mathbb{Q}}}
\newunicodechar{𝔻}{\ensuremath{\mathbb{D}}}
\newunicodechar{α}{\ensuremath{\alpha}}
\newunicodechar{β}{\ensuremath{\beta}}
\newunicodechar{γ}{\ensuremath{\gamma}}
\newunicodechar{δ}{\ensuremath{\delta}}
\newunicodechar{ε}{\ensuremath{\varepsilon}}
\newunicodechar{θ}{\ensuremath{\theta}}
\newunicodechar{λ}{\ensuremath{\lambda}}
\newunicodechar{μ}{\ensuremath{\mu}}
\newunicodechar{σ}{\ensuremath{\sigma}}
\newunicodechar{τ}{\ensuremath{\tau}}
\newunicodechar{φ}{\ensuremath{\varphi}}
\newunicodechar{ψ}{\ensuremath{\psi}}
\newunicodechar{ω}{\ensuremath{\omega}}
\newunicodechar{Σ}{\ensuremath{\Sigma}}
% majuscules produites par \MakeUppercase dans les titres (α-aminés -> Α)
\newunicodechar{Α}{\ensuremath{\alpha}}
\newunicodechar{Β}{\ensuremath{\beta}}
\newunicodechar{Γ}{\ensuremath{\Gamma}}
\newunicodechar{Δ}{\ensuremath{\Delta}}
\newunicodechar{Ε}{\ensuremath{\varepsilon}}
\newunicodechar{Θ}{\ensuremath{\Theta}}
\newunicodechar{Λ}{\ensuremath{\Lambda}}
\newunicodechar{Μ}{\ensuremath{\mu}}
\newunicodechar{Τ}{\ensuremath{\tau}}
\newunicodechar{Φ}{\ensuremath{\Phi}}
\newunicodechar{Ψ}{\ensuremath{\Psi}}
\newunicodechar{Ω}{\ensuremath{\Omega}}
\newunicodechar{↦}{\ensuremath{\mapsto}}
\newunicodechar{∈}{\ensuremath{\in}}
\newunicodechar{∉}{\ensuremath{\notin}}
\newunicodechar{∧}{\ensuremath{\wedge}}
\newunicodechar{⇔}{\ensuremath{\Leftrightarrow}}
\newunicodechar{⟺}{\ensuremath{\Longleftrightarrow}}
\newunicodechar{⇒}{\ensuremath{\Rightarrow}}
\newunicodechar{⟹}{\ensuremath{\Longrightarrow}}
\newunicodechar{∘}{\ensuremath{\circ}}
\newunicodechar{∪}{\ensuremath{\cup}}
\newunicodechar{∩}{\ensuremath{\cap}}
\newunicodechar{≡}{\ensuremath{\equiv}}
\newunicodechar{⊂}{\ensuremath{\subset}}
\newunicodechar{⊥}{\ensuremath{\perp}}
\newunicodechar{∃}{\ensuremath{\exists}}
\newunicodechar{∀}{\ensuremath{\forall}}
\newunicodechar{∛}{\ensuremath{\sqrt[3]{\,}}}
\newunicodechar{∜}{\ensuremath{\sqrt[4]{\,}}}
\newunicodechar{⃗}{\ensuremath{{}^{\to}}}
\newunicodechar{ⁿ}{\ifmmode^{n}\else\textsuperscript{\ensuremath{n}}\fi}
\newunicodechar{ˣ}{\ifmmode^{x}\else\textsuperscript{\ensuremath{x}}\fi}
\newunicodechar{ᵇ}{\ifmmode^{b}\else\textsuperscript{\ensuremath{b}}\fi}
\newunicodechar{ʳ}{\ifmmode^{r}\else\textsuperscript{\ensuremath{r}}\fi}
\newunicodechar{ᵘ}{\ifmmode^{u}\else\textsuperscript{\ensuremath{u}}\fi}
\newunicodechar{ʸ}{\ifmmode^{y}\else\textsuperscript{\ensuremath{y}}\fi}
\newunicodechar{ᵃ}{\ifmmode^{a}\else\textsuperscript{\ensuremath{a}}\fi}
\newunicodechar{ᵉ}{\ifmmode^{e}\else\textsuperscript{\ensuremath{e}}\fi}
\newunicodechar{ᵏ}{\ifmmode^{k}\else\textsuperscript{\ensuremath{k}}\fi}
\newunicodechar{ᵖ}{\ifmmode^{p}\else\textsuperscript{\ensuremath{p}}\fi}
\newunicodechar{ᴺ}{\ifmmode^{N}\else\textsuperscript{\ensuremath{N}}\fi}
\newunicodechar{ᶜ}{\ifmmode^{c}\else\textsuperscript{\ensuremath{c}}\fi}
\newunicodechar{ʰ}{\ifmmode^{h}\else\textsuperscript{\ensuremath{h}}\fi}
\newunicodechar{ᵠ}{\ifmmode^{q}\else\textsuperscript{\ensuremath{q}}\fi}
\newunicodechar{ₙ}{\ifmmode_{n}\else\textsubscript{\ensuremath{n}}\fi}
\newunicodechar{ₐ}{\ifmmode_{a}\else\textsubscript{\ensuremath{a}}\fi}
\newunicodechar{ᵢ}{\ifmmode_{i}\else\textsubscript{\ensuremath{i}}\fi}
\newunicodechar{ᵣ}{\ifmmode_{r}\else\textsubscript{\ensuremath{r}}\fi}
\newunicodechar{ₖ}{\ifmmode_{k}\else\textsubscript{\ensuremath{k}}\fi}
\newunicodechar{ᵦ}{\ifmmode_{\beta}\else\textsubscript{\ensuremath{\beta}}\fi}
\newunicodechar{ₓ}{\ifmmode_{x}\else\textsubscript{\ensuremath{x}}\fi}
\newfontfamily\popdisplay{ArchivoBlack-Regular.ttf}[Path=../../../fonts/]
\newfontfamily\poplogo{SpaceGrotesk-Bold.ttf}[Path=../../../fonts/]
\newfontfamily\popsemi{PlusJakartaSans-SemiBold.ttf}[Path=../../../fonts/]
\newfontfamily\popxbold{PlusJakartaSans-ExtraBold.ttf}[Path=../../../fonts/]
\newfontfamily\popxboldls{PlusJakartaSans-ExtraBold.ttf}[Path=../../../fonts/,LetterSpace=3.0]

\setlist[enumerate]{leftmargin=1.7em,itemsep=5pt,topsep=4pt}
\setlist[itemize]{leftmargin=1.7em,itemsep=4pt,topsep=4pt,label={\color{poporange}\textbullet}}
\setlength{\parindent}{0pt}
\setlength{\parskip}{5pt}
\renewcommand{\arraystretch}{1.25}

% pgfplots : style Pop par défaut
\pgfplotsset{
  every axis/.append style={axis line style={popink,line width=1pt},tick style={popink},
    grid style={popline,line width=0.5pt},label style={font=\small},tick label style={font=\footnotesize}},
  popcurve/.style={poporange,line width=1.8pt,no markers,smooth},
}
% Même style disponible dans un \draw TikZ simple (hors pgfplots)
\tikzset{popcurve/.style={poporange,line width=1.8pt}}
\tikzset{popgrid/.style={popline,very thin}}
\colorlet{popcurve}{poporange} % le nom du style est parfois utilisé comme couleur

% ============================================================
% Pill / squiggle
% ============================================================
\newcommand{\poppill}[3]{%
  \tikz[baseline=-0.55ex]\node[draw=popink,line width=1.2pt,rounded corners=2.6pt,%
    fill=#2,inner xsep=6.5pt,inner ysep=3pt]{%
    \popxboldls\fontsize{7}{7}\selectfont\color{#3}\MakeUppercase{#1}};%
}
\newcommand{\popsquiggle}[1]{%
  \tikz[baseline=-0.4ex]{
    \draw[#1,line width=2.6pt,line cap=round]
      plot [smooth] coordinates {(0,0.09) (0.34,0.01) (0.68,0.21) (1.02,0.01) (1.36,0.10)};
  }%
}
% Mot-clé surligné (marqueur orange sous le texte)
\newcommand{\cle}[1]{{\popxbold\color{popdark}#1}}

% ============================================================
% En-tête / pied
% ============================================================
\pagestyle{fancy}
\fancyhf{}
\renewcommand{\headrulewidth}{0pt}
\renewcommand{\footrulewidth}{0pt}
\lhead{\raisebox{-0.15ex}{\poplogo\fontsize{13}{13}\selectfont\color{popink}OnBuch\color{poporange}+}}
\rhead{\poppill{\DOCMATIERE\ \textbullet\ \DOCNIVEAU\ \textbullet\ Leçon \DOCLECON}{white}{popink}}
\lfoot{\popxbold\fontsize{7.6}{7.6}\selectfont\color{popmuted}\MakeUppercase{Cours signé OnBuch+}\ \textbullet\ {\popsemi\DOCTITRE}}
\rfoot{\tikz[baseline=-0.6ex]\node[rounded corners=8.5pt,fill=popink,minimum height=17pt,minimum width=17pt,inner xsep=5pt,inner sep=0]{\popxbold\fontsize{8}{8}\selectfont\color{popcream}\thepage\,/\,\pageref*{LastPage}};}

% ============================================================
% Titres
% ============================================================
\newcommand{\chaptitre}[2]{%
  \begin{tcolorbox}[enhanced,colback=poporange,colframe=popink,boxrule=2.6pt,arc=11pt,
      drop shadow=popink,left=15pt,right=15pt,top=12pt,bottom=11pt,before skip=3pt,after skip=18pt]
    \poppill{#2}{popink}{popcream}\par\vspace{9pt}
    {\raggedright\hyphenpenalty=10000\exhyphenpenalty=10000\popdisplay\fontsize{22}{24}\selectfont\color{popcream}\MakeUppercase{#1}\par}
  \end{tcolorbox}
}
% Section numérotée : pastille orange avec le numéro + titre Archivo
\newcounter{popsec}
\newcommand{\coursec}[1]{%
  \stepcounter{popsec}\setcounter{popsub}{0}%
  \par\vspace{16pt}\noindent
  \begin{minipage}{\linewidth}
  \tikz[baseline=-0.7ex]\node[draw=popink,line width=1.6pt,fill=poporange,rounded corners=4pt,minimum size=22pt,inner sep=0]{\popdisplay\fontsize{12}{12}\selectfont\color{popcream}\thepopsec};%
  \hspace{8pt}{\popdisplay\fontsize{15}{17}\selectfont\color{popink}\MakeUppercase{#1}}\par
  \vspace{4pt}\noindent{\color{poporange}\rule{36pt}{3.6pt}}
  \end{minipage}\par\nopagebreak\vspace{8pt}
}
\newcounter{popsub}
\newcommand{\courssub}[1]{%
  \stepcounter{popsub}%
  \par\vspace{9pt}\noindent{\popxbold\fontsize{11.5}{13}\selectfont\color{popink}{\color{poporange}\thepopsec.\thepopsub}\hspace{6pt}#1}\par\nopagebreak\vspace{4pt}
}

% ============================================================
% Boîtes pédagogiques — même recette : fond blanc, bordure épaisse colorée,
% ombre dure encre, badge plein. Titre optionnel [#1].
% ============================================================
\tcbset{popbase/.style={breakable,enhanced,colback=white,boxrule=2.3pt,arc=9pt,
  shadow={3pt}{-3pt}{0pt}{popink},left=14pt,right=14pt,top=11pt,bottom=11pt,before skip=11pt,after skip=15pt,
  fontupper=\popsemi\fontsize{10.6}{14.5}\selectfont\color{popink},
  fontlower=\popsemi\fontsize{10.6}{14.5}\selectfont\color{popink}}}
% Coche dessinée (le glyphe ✓ n'existe pas dans les polices du document)
\newcommand{\popcheck}[1]{\tikz[baseline=-0.5ex]\draw[#1,line width=1.6pt,line cap=round,line join=round](-0.13,0.0)--(-0.04,-0.09)--(0.13,0.1);}
\providecommand{\coche}{\popcheck{popgreen}}
\providecommand{\resultat}[1]{\par\smallskip\noindent\fcolorbox{popgreen}{popgreenL}{\parbox{\dimexpr\linewidth-2\fboxsep-2\fboxrule\relax}{\textbf{Résultat.} #1}}\par\smallskip}
\newcommand{\popboxhead}[3]{\poppill{#1}{#2}{white}\ifx\relax#3\relax\else\hspace{6pt}{\popxbold\fontsize{10.6}{12}\selectfont\color{popink}#3}\fi\par\vspace{7pt}}

\newtcolorbox{aretenir}[1][]{popbase,colframe=popgreen,before upper={\popboxhead{À retenir}{popgreen}{#1}}}
\newtcolorbox{exemplebox}[1][]{popbase,colframe=poporange,before upper={\popboxhead{Exemple}{poporange}{#1}}}
\newtcolorbox{definition}[1][]{popbase,colframe=popblue,before upper={\popboxhead{Définition}{popblue}{#1}}}
\newtcolorbox{propriete}[1][]{popbase,colframe=poppurple,before upper={\popboxhead{Propriété}{poppurple}{#1}}}
\newtcolorbox{methode}[1][]{popbase,colframe=popgold,before upper={\popboxhead{Méthode}{popgold}{#1}}}
\newtcolorbox{attention}[1][]{popbase,colframe=poppink,before upper={\popboxhead{Attention}{poppink}{#1}}}
\newtcolorbox{experience}[1][]{popbase,colframe=popblue,colback=popblueL,before upper={\popboxhead{Expérience}{popblue}{#1}}}
% « Le savais-tu ? » : carte violette pleine (contexte, culture, Cameroun)
\newtcolorbox{savaistu}[1][]{popbase,colback=poppurple,colframe=popink,
  fontupper=\popsemi\fontsize{10.4}{14.2}\selectfont\color{white},
  before upper={\poppill{Le savais-tu\,?}{white}{poppurple}\ifx\relax#1\relax\else\hspace{6pt}{\popxbold\color{white}#1}\fi\par\vspace{7pt}}}
% Exercice résolu : énoncé en haut, \tcblower puis la solution sur fond crème
\newcounter{popexo}\newcounter{popexr}
\newtcolorbox{exoresolu}[1][]{popbase,colframe=poporange,colbacklower=popcream,
  segmentation style={popink,line width=1.2pt,dashed},
  before upper={\stepcounter{popexr}\popboxhead{Exercice résolu \thepopexr}{poporange}{#1}},
  before lower={\poppill{Solution}{popink}{popcream}\par\vspace{6pt}}}
% Exercice (non résolu en place) — la correction est dans la partie Corrigés
\newtcolorbox{exercice}[1][]{popbase,colframe=popink,
  before upper={\stepcounter{popexo}\popboxhead{Exercice \thepopexo}{popink}{#1}}}
% Corrigé
\newtcolorbox{corrige}[1][]{popbase,colframe=popgreen,colback=popgreenL,
  before upper={\popboxhead{Corrigé}{popgreen}{#1}}}
% Objectifs / prérequis / bilan
\newtcolorbox{objectifs}{popbase,colframe=popink,colback=poporangeL,
  before upper={\popboxhead{Objectifs de la leçon}{popink}{}}}
\newtcolorbox{prerequis}{popbase,colframe=popink,colback=popblueL,
  before upper={\popboxhead{Prérequis}{popblue}{}}}
\newtcolorbox{bilan}{popbase,colframe=popink,colback=white,boxrule=2.8pt,
  before upper={\popboxhead{Fiche bilan}{poporange}{L'essentiel à connaître pour l'examen}}}
% Figure encadrée avec légende
\newtcolorbox{popfigure}[1][]{enhanced,breakable=false,colback=white,colframe=popline,boxrule=1.4pt,arc=8pt,
  left=10pt,right=10pt,top=10pt,bottom=8pt,before skip=10pt,after skip=12pt,halign=center,
  lower separated=false,halign lower=center,
  fontlower=\popsemi\fontsize{9}{11.5}\selectfont\color{popmuted},
  #1}
\newcommand{\legende}[1]{\par\vspace{6pt}{\popsemi\fontsize{9}{11.5}\selectfont\color{popmuted}#1}}

% ============================================================
% Signature OnBuch+
% ============================================================
\newcommand{\poplogoinline}[1]{{\poplogo\fontsize{#1}{#1}\selectfont\color{popink}OnBuch\color{poporange}+}}
\newcommand{\signatureonbuch}{%
  \begin{tcolorbox}[enhanced,colback=popink,colframe=popink,boxrule=2.2pt,arc=10pt,
      drop shadow=poporange,left=14pt,right=14pt,top=10pt,bottom=10pt,before skip=0pt,after skip=0pt]
    \tikz[baseline=-0.7ex]\node[circle,fill=popgreen,draw=popcream,line width=1.3pt,minimum size=22pt,inner sep=0]{\popcheck{white}};%
    \hspace{9pt}\begin{minipage}[c]{0.62\linewidth}
      {\popxbold\fontsize{12}{13}\selectfont\color{popcream}Signé\ \textbullet\ {\poplogo OnBuch\color{poporange}+}}\par\vspace{2pt}
      {\raggedright\popsemi\fontsize{8}{10.5}\selectfont\color{popline}\MakeUppercase{Équipe pédagogique OnBuch+}\\\MakeUppercase{Conforme au programme officiel MINESEC}\par}
    \end{minipage}\hfill
    {\poplogo\fontsize{22}{22}\selectfont\color{popcream}OnBuch\color{poporange}+}
  \end{tcolorbox}%
}

% ============================================================
% Page de garde
% #1 titre · #2 matière · #3 niveau · #4 module · #5 n° leçon · #6 durée ·
% #7 description / objectifs (texte ou itemize)
% ============================================================
\newcommand{\pagegarde}[7]{%
  \thispagestyle{empty}
  \newgeometry{a4paper,margin=1.7cm,top=1.6cm,bottom=1.4cm}%
  \begin{tikzpicture}[remember picture,overlay]
    \fill[poporange!18] ([xshift=-3.2cm,yshift=-3.6cm]current page.north east) circle (4.6cm);
    \fill[poppurple!14] ([xshift=2.4cm,yshift=7.6cm]current page.south west) circle (3.8cm);
  \end{tikzpicture}%
  {\poplogo\fontsize{30}{30}\selectfont\color{popink}OnBuch\color{poporange}+}\hfill
  \raisebox{0.6ex}{\poppill{Cours signé \textbullet\ Premium}{popink}{popcream}}\par
  \vspace{1pt}\popsquiggle{poppurple}\par
  \vspace{8pt}
  \begin{tcolorbox}[enhanced,colback=poporange,colframe=popink,boxrule=2.8pt,arc=13pt,
      drop shadow=popink,left=18pt,right=18pt,top=12pt,bottom=13pt,after skip=10pt]
    \poppill{#2 \textbullet\ #3}{popink}{popcream}\hspace{5pt}\poppill{#4}{white}{popink}\par\vspace{8pt}
    {\popdisplay\fontsize{14}{14}\selectfont\color{popink}LEÇON #5}\par\vspace{4pt}
    {\raggedright\hyphenpenalty=10000\exhyphenpenalty=10000\popdisplay\fontsize{25}{27}\selectfont\color{popcream}\MakeUppercase{#1}\par}
    \vspace{8pt}
    {\popxbold\fontsize{9.5}{9.5}\selectfont\color{popink}\MakeUppercase{Durée conseillée : #6}}
  \end{tcolorbox}
  \begin{tcolorbox}[enhanced,colback=white,colframe=popink,boxrule=2.2pt,arc=10pt,
      drop shadow=popink,left=16pt,right=16pt,top=10pt,bottom=10pt,after skip=8pt]
    \poppill{Dans cette leçon}{poppurple}{white}\par\vspace{5pt}
    \popsemi\fontsize{9.3}{12.6}\selectfont\color{popink}#7
  \end{tcolorbox}
  \noindent\begin{minipage}[t]{0.487\linewidth}
    \begin{tcolorbox}[enhanced,colback=popgreen,colframe=popink,boxrule=2.2pt,arc=10pt,
        drop shadow=popink,left=11pt,right=11pt,top=10pt,bottom=10pt,height=3.1cm,valign=center]
      \tcbox[colback=white,colframe=popink,boxrule=1.3pt,arc=5pt,boxsep=0pt,nobeforeafter,
        left=3pt,right=3pt,top=3pt,bottom=3pt,valign=center]{\qrcode[height=1.9cm]{https://play.google.com/store/apps/details?id=com.luvvix.onbuch&hl=fr}}%
      \hspace{8pt}\begin{minipage}[c]{0.5\linewidth}
        {\popdisplay\fontsize{11}{12.5}\selectfont\color{white}\MakeUppercase{Télécharge OnBuch}}\par\vspace{4pt}
        {\raggedright\popsemi\fontsize{8.4}{11}\selectfont\color{white}Gratuit \textbullet\ Google Play\\Tuteur IA, annales, concours, résultats\par}
      \end{minipage}
    \end{tcolorbox}
  \end{minipage}\hfill
  \begin{minipage}[t]{0.487\linewidth}
    \begin{tcolorbox}[enhanced,colback=poppurple,colframe=popink,boxrule=2.2pt,arc=10pt,
        drop shadow=popink,left=11pt,right=11pt,top=10pt,bottom=10pt,height=3.1cm,valign=center]
      \tikz[baseline=-0.7ex]\node[circle,fill=white,draw=popink,line width=1.3pt,minimum size=1.6cm,inner sep=0]{\popdisplay\fontsize{24}{24}\selectfont\color{poppurple}+};%
      \hspace{8pt}\begin{minipage}[c]{0.6\linewidth}
        {\popdisplay\fontsize{11}{12.5}\selectfont\color{white}\MakeUppercase{Passe à OnBuch+}}\par\vspace{4pt}
        {\raggedright\popsemi\fontsize{8.4}{11}\selectfont\color{white}Tous les cours signés, Tuteur IA illimité, fiches PDF — onglet OnBuch+ de l'app\par}
      \end{minipage}
    \end{tcolorbox}
  \end{minipage}
  \vspace{8pt}
  \popcontenu
  \vfill
  \signatureonbuch
  \restoregeometry
  \newpage
}

% Rangée « Ce cours contient » (page de garde)
\newcommand{\poptile}[3]{% #1 couleur #2 gros texte #3 légende
  \begin{tcolorbox}[enhanced,colback=#1,colframe=popink,boxrule=2pt,arc=9pt,shadow={2.5pt}{-2.5pt}{0pt}{popink},
      left=6pt,right=6pt,top=9pt,bottom=9pt,halign=center,valign=center,height=1.75cm,width=\linewidth,nobeforeafter]
    {\popdisplay\fontsize{12.5}{13}\selectfont\color{white}\MakeUppercase{#2}}\par\vspace{3pt}
    {\centering\popsemi\fontsize{7.8}{9.5}\selectfont\color{white}#3\par}
  \end{tcolorbox}}
\newcommand{\popcontenu}{%
  \par\noindent{\popxboldls\fontsize{7.5}{7.5}\selectfont\color{popmuted}CE COURS SIGNÉ CONTIENT}\par\vspace{7pt}
  \noindent\begin{minipage}[t]{0.235\linewidth}\poptile{popblue}{Cours}{complet, illustré, pas à pas}\end{minipage}\hfill
  \begin{minipage}[t]{0.235\linewidth}\poptile{poporange}{Méthodes}{et exercices résolus}\end{minipage}\hfill
  \begin{minipage}[t]{0.235\linewidth}\poptile{poppink}{Exercices}{type examen + corrigés}\end{minipage}\hfill
  \begin{minipage}[t]{0.235\linewidth}\poptile{popgreen}{Fiche bilan}{l'essentiel à retenir}\end{minipage}\par}

% Sommaire
\newtcolorbox{sommaire}{enhanced,colback=white,colframe=popink,boxrule=2.2pt,arc=10pt,
  drop shadow=popink,left=18pt,right=18pt,top=13pt,bottom=13pt,before skip=6pt,after skip=20pt,
  before upper={\poppill{Sommaire}{poppurple}{white}\par\vspace{9pt}\popsemi\fontsize{10.6}{14.5}\selectfont\color{popink}}}

% Signature de fin de cours — poussée tout en bas de la dernière page
\newcommand{\findecours}{%
  \par\vfill\nopagebreak
  \begin{center}\popsquiggle{poporange}\end{center}\vspace{4pt}
  \signatureonbuch
}
\AtBeginDocument{\def\llbracket{\mathopen{[\mkern-3.5mu[}}\def\rrbracket{\mathclose{]\mkern-3.5mu]}}} % intervalles d'entiers (glyphe ⟦ absent de la police)
% Glyphes absents de Fira Math : redéfinis à partir de glyphes disponibles
\AtBeginDocument{%
  \def\setminus{\mathbin{\backslash}}\let\smallsetminus\setminus
  \def\vdots{\mathord{\vbox{\baselineskip3.2pt\lineskiplimit0pt\kern4pt\hbox{.}\hbox{.}\hbox{.}}}}%
  \def\ddots{\mathinner{\mkern1mu\raise6.5pt\hbox{.}\mkern2mu\raise3.6pt\hbox{.}\mkern2mu\raise0.7pt\hbox{.}\mkern1mu}}%
  \def\triangle{\mathord{\scalebox{0.9}{$\symup{\Delta}$}}}%
  \def\nabla{\mathord{\raisebox{\depth}{\scalebox{1}[-1]{$\symup{\Delta}$}}}}%
  \def\checkmark{\popcheck{popdark}}%
  \def\star{\mathbin{\ast}}\def\bigstar{\mathord{\ast}}%
  \def\diamond{\mathbin{\scalebox{0.6}{$\Diamond$}}}%
  \def\bigcup{\mathop{\mathchoice{\raisebox{-0.3em}{\scalebox{1.7}{$\cup$}}}{\scalebox{1.25}{$\cup$}}{\cup}{\cup}}\displaylimits}%
  \def\bigcap{\mathop{\mathchoice{\raisebox{-0.3em}{\scalebox{1.7}{$\cap$}}}{\scalebox{1.25}{$\cap$}}{\cap}{\cap}}\displaylimits}%
  \def\lesseqqgtr{\mathrel{\lessgtr}}\def\bigoplus{\mathop{\oplus}\displaylimits}%
}
\providecommand{\ssi}{si et seulement si} % abréviation fréquente
\AtBeginDocument{\providecommand{\wideparen}{\overparen}} % arc de cercle

%% FIGFIX-BEGIN v5 — correctifs globaux des figures (docs/onbuch-generation/tools/figfix)
\usepackage{adjustbox}\usepackage{etoolbox}\tcbuselibrary{raster}
% toute figure TikZ/pgfplots plus large que la ligne est réduite à la largeur disponible
\BeforeBeginEnvironment{tikzpicture}{\begin{adjustbox}{max width=\linewidth}}
\AfterEndEnvironment{tikzpicture}{\end{adjustbox}}
% légendes pgfplots lisibles par-dessus les courbes
\pgfplotsset{every axis/.append style={legend style={fill=white,fill opacity=0.92,text opacity=1,draw=black!15,inner sep=3pt}}}
% étiquettes d'arêtes (syntaxe edge["..."]) avec fond blanc
\tikzset{every edge quotes/.append style={fill=white,inner sep=1.2pt}}
%% FIGFIX-END
\begin{document}

\pagegarde{Consolidation des acquis de la méthodologie de l’exercice sur texte philosophique}{Philosophie}{Tle Tech}{Philosophie – classe de Terminale technique}{N17}{2 séances (fiche de progression harmonisée)}{Cette leçon de consolidation structure et automatise la méthodologie de l'exercice sur texte philosophique, de la lecture analytique à la rédaction finale, à travers des entraînements progressifs ancrés dans des thématiques camerounaises. Elle vise à transformer l'élève en un lecteur actif capable de problématiser, structurer et argumenter avec précision pour réussir l'épreuve du Baccalauréat.
\vspace{4pt}
\begin{multicols}{2}\small
\begin{itemize}
\item À la fin de cette leçon, je sais identifier la thèse, les concepts clés et la structure logique d'un texte philosophique donné au Bac.
\item Je sais construire une problématique pertinente qui guide l'explication et évite le hors-sujet ou la paraphrase.
\item Je sais rédiger une introduction complète (accroche, présentation, thèse, problématique, annonce du plan) en 15 minutes chrono.
\item Je sais organiser un développement en deux ou trois parties équilibrées, articulées par des transitions logiques, en m'appuyant sur l'analyse interne du texte.
\item Je sais distinguer et traiter les trois types de questions classiques : Explication, Analyse/Commentaire, Discussion/Critique.
\item Je sais mobiliser mes connaissances philosophiques (auteurs, notions, oppositions) pour éclairer le texte sans le noyer sous les citations.
\end{itemize}\end{multicols}}

\begin{sommaire}
\begin{multicols}{2}
\begin{enumerate}
\item 1. Diagnostic \& Rappel des Exigences Officielles
\item 2. L'Atelier Problématisation : Du Sujet à la Question Directrice
\item 3. Rédaction Chirurgicale : Introduction, Développement, Conclusion
\item 4. Maîtrise des Questions Types \& Pièges Spécifiques Bac Tech
\item 5. Simulation Bac Blanc \& Méta-cognition (Auto-évaluation)
\item Activité d'intégration
\item Méthodes et exercices résolus
\item Exercices
\item Corrigés des exercices
\item Fiche bilan
\end{enumerate}
\end{multicols}
\end{sommaire}

\begin{objectifs}
\begin{itemize}
\item À la fin de cette leçon, je sais identifier la thèse, les concepts clés et la structure logique d'un texte philosophique donné au Bac.
\item Je sais construire une problématique pertinente qui guide l'explication et évite le hors-sujet ou la paraphrase.
\item Je sais rédiger une introduction complète (accroche, présentation, thèse, problématique, annonce du plan) en 15 minutes chrono.
\item Je sais organiser un développement en deux ou trois parties équilibrées, articulées par des transitions logiques, en m'appuyant sur l'analyse interne du texte.
\item Je sais distinguer et traiter les trois types de questions classiques : Explication, Analyse/Commentaire, Discussion/Critique.
\item Je sais mobiliser mes connaissances philosophiques (auteurs, notions, oppositions) pour éclairer le texte sans le noyer sous les citations.
\item Je sais conclure en synthétisant la démarche, en répondant explicitement à la problématique et en ouvrant sur une perspective concrète.
\item Je sais m'auto-évaluer à l'aide de la grille officielle de notation du Baccalauréat technique pour corriger mes faiblesses récurrentes.
\end{itemize}
\end{objectifs}

\begin{prerequis}
\begin{itemize}
\item Maîtrise des grands courants philosophiques au programme (Rationalisme, Empirisme, Criticisme, Existentialisme, Philosophie africaine).
\item Connaissance de la structure type d'une dissertation philosophique (Introduction, Développement, Conclusion) acquise en Première.
\item Maîtrise des connecteurs logiques (donc, cependant, par conséquent, en outre) pour assurer la cohérence du discours.
\item Capacité à définir rigoureusement les concepts clés du programme (Liberté, État, Science, Culture, Tradition, Modernité).
\item Expérience de la lecture active (surlignage, annotation, repérage des connecteurs) lors des séances de TP précédents.
\end{itemize}
\end{prerequis}

\newpage

\coursec{Diagnostic \& Rappel des Exigences Officielles}

En plein cœur de Yaoundé, un élève de Terminale Technique fixe l'extrait de Paulin Hountondji au Bac : il connaît l'auteur, mais reste paralysé face à la consigne « Dégagez la thèse et expliquez-la ». Comme ce jeune commerçant du Marché Central qui doit décrypter un contrat écrit en français juridique pour protéger sa marchandise, l'élève réalise que comprendre un texte philosophique ne demande pas de la magie, mais une méthode rigoureuse, transposable à tout document complexe de la vie citoyenne. L'angoisse ne vient pas de la difficulté du texte : elle vient de l'absence de protocole. Or un technicien, qu'il soit mécanicien, électricien ou comptable, sait qu'on ne répare pas une machine sans en connaître le schéma de fonctionnement.

Cette première séance de consolidation pose donc la question fondatrice de toute la leçon : \textbf{qu'attend exactement le correcteur du Baccalauréat technique lorsqu'il lit une copie d'exercice sur texte, et pourquoi tant de candidats échouent-ils alors qu'ils « connaissent » leur cours ?} Pour y répondre, nous partirons du diagnostic dressé par les rapports de jury, nous rappellerons la définition officielle de l'exercice, nous distinguerons ses deux gestes fondamentaux (expliquer et discuter), puis nous détaillerons la grille d'évaluation et les trois types de consignes.

\courssub{Le diagnostic : ce que disent les rapports de jury (sessions 2023-2024)}

\begin{experience}[Observer avant d'agir : la démarche du diagnostic]
Avant de réparer un moteur, le mécanicien écoute, observe, mesure. Avant de réviser une méthode, le philosophe observe les copies réelles. Les rapports de jury du Baccalauréat technique (sessions 2023 et 2024) constituent notre « banc d'essai » : ils recensent les fautes les plus fréquentes commises par les candidats dans l'exercice sur texte. Retenons les quatre erreurs majeures :

\begin{enumerate}
\item \textbf{La paraphrase.} Le candidat recopie le texte en changeant quelques mots, sans jamais dégager la structure de la pensée. Il « répète » l'auteur au lieu de l'\emph{expliquer}. Le correcteur sanctionne immédiatement : paraphraser, c'est prouver qu'on n'a pas compris.
\item \textbf{Le hors-sujet.} Le candidat récite son cours sur « la technique » ou « l'État » en général, alors que le texte pose une question précise (par exemple : la technique libère-t-elle ou aliène-t-elle le travailleur ?). Tout ce qui n'est pas \emph{ancré dans le texte} est perdu.
\item \textbf{L'absence de problématique.} La copie enchaîne des remarques sans question directrice. Or un texte philosophique répond toujours à un problème : ne pas le formuler, c'est passer à côté de l'enjeu.
\item \textbf{Le plan invisible.} Le développement est un bloc compact, sans annonce des étapes, sans transitions, sans paragraphes identifiables. Le correcteur ne peut pas suivre la progression, donc ne peut pas la noter.
\end{enumerate}

\textbf{Interprétation :} ces quatre erreurs ont une cause commune — le candidat traite l'exercice sur texte comme un résumé ou comme une dissertation déguisée, alors qu'il s'agit d'un exercice spécifique, codifié, avec des attendus précis.
\end{experience}

\begin{attention}[Les quatre fautes qui plafonnent la note]
Les rapports de jury sont formels : une copie paraphrastique ne peut pas dépasser la moyenne, même si elle est bien écrite. De même, un hors-sujet, même brillant, est sanctionné car il ne répond pas au contrat de l'exercice. \textbf{Réflexe de survie :} à chaque paragraphe rédigé, demandez-vous : « Est-ce que je m'appuie sur une phrase précise du texte ? Est-ce que je fais avancer la réponse à ma problématique ? » Si la réponse est non deux fois de suite, vous dérivez.
\end{attention}

\courssub{La définition officielle : ni résumé, ni dissertation}

\begin{definition}[L'exercice sur texte philosophique]
L'\cle{exercice sur texte} est l'\textbf{explication ordonnée de la pensée d'un auteur à partir de son propre texte}, pour en vérifier la cohérence et la portée. Il comporte deux moments inséparables :
\begin{itemize}
\item l'\cle{explication}, qui restitue la thèse, les concepts et les arguments de l'auteur \emph{depuis l'intérieur} du texte ;
\item la \cle{discussion}, qui évalue la validité, la portée et les limites de cette pensée \emph{depuis l'extérieur}.
\end{itemize}
Il ne s'agit donc \textbf{ni d'un résumé} (qui abrège sans analyser), \textbf{ni d'une dissertation} (qui répond à une question personnelle en mobilisant sa culture).
\end{definition}

Prenons une analogie technique. Face à un appareil électrique en panne, le résumé consisterait à dire : « c'est un appareil avec des fils et un moteur ». L'explication, elle, consiste à \emph{ouvrir le capot} : identifier chaque composant, comprendre comment ils sont reliés, pourquoi le concepteur les a disposés ainsi. La discussion, enfin, consiste à se demander : ce montage est-il le meilleur possible ? Tient-il dans toutes les conditions d'usage ? L'exercice sur texte est exactement ce double geste : \textbf{démontage interne, puis expertise externe}.

\courssub{Explication et discussion : les deux gestes fondamentaux}

\begin{propriete}[La distinction explication / discussion]
\textbf{Expliquer} (moment 1) : se placer \emph{à l'intérieur} du système de pensée de l'auteur. On adopte temporairement ses définitions, on suit son raisonnement pas à pas, on montre comment la thèse s'impose à partir des arguments. Règle d'or : \textbf{fidélité}. On n'ajoute rien, on ne critique pas encore. Formules typiques : « L'auteur entend par là que... », « Cet argument repose sur l'idée que... », « Le texte procède en deux temps : d'abord..., ensuite... ».

\textbf{Discuter} (moment 2) : se placer \emph{à l'extérieur} pour évaluer. On interroge les présupposés (ce que l'auteur admet sans le démontrer), la portée (sa thèse vaut-elle partout et toujours ?) et les limites (quelles objections peut-on lui opposer ?). Formules typiques : « Cette thèse suppose que..., ce qui peut être contesté », « On objectera cependant que... », « La position de l'auteur éclaire..., mais elle peine à rendre compte de... ».
\end{propriete}

\begin{exemplebox}[Un même texte, deux gestes]
Soit un extrait où un auteur affirme : « Le travail est ce par quoi l'homme se humanise. »

\emph{Explication (dedans) :} L'auteur définit le travail non comme une simple activité de subsistance, mais comme une transformation de la nature qui transforme en retour l'homme lui-même : en fabriquant des outils, l'homme développe des capacités (prévision, patience, coopération) qui le distinguent de l'animal. La thèse repose donc sur l'idée d'un double mouvement : transformer le monde, c'est se transformer soi-même.

\emph{Discussion (dehors) :} Cette analyse, héritée de Hegel et de Marx, suppose que tout travail est créateur. Or le travail à la chaîne, répétitif et parcellisé, peut au contraire abrutir — c'est l'objection de l'aliénation. La thèse de l'auteur est donc forte pour penser le travail accompli, mais doit être complétée pour penser le travail subi, par exemple celui de certains ouvriers des zones industrielles de Douala.
\end{exemplebox}

\begin{attention}[Erreur fatale : confondre explication et dissertation]
L'erreur la plus grave consiste à confondre l'\emph{explication} (fidélité au texte) et la \emph{dissertation} (réponse personnelle à une question). \textbf{Ne faites pas votre propre dissertation en citant l'auteur.} Si la consigne demande d'expliquer le texte, votre culture personnelle (Platon, Descartes, Nietzsche) ne doit intervenir que dans la discussion, et encore, avec sobriété. Une copie qui aligne dix noms de philosophes sans jamais analyser les phrases du texte est une copie hors contrat. Le texte est votre pièce maîtresse : tout part de lui, tout y revient.
\end{attention}

\courssub{La méthode des trois lectures}

\begin{methode}[Les trois lectures successives du texte]
\begin{enumerate}
\item \textbf{Lecture 1 — Découverte.} Qui est l'auteur ? De quelle époque ? Quel est le sujet général du passage ? Quel est le ton (polémique, descriptif, ironique) ? Objectif : situer le texte sans encore l'analyser.
\item \textbf{Lecture 2 — Analyse.} Quelle est la \cle{thèse} (souvent une phrase clé, au début ou à la fin) ? Quels sont les \cle{concepts} employés et comment sont-ils définis ? Quels \cle{arguments} soutiennent la thèse ? Quelle est la \cle{structure} (étapes du raisonnement) ? Objectif : reconstituer la machine argumentative, crayon à la main, en soulignant et en numérotant les étapes.
\item \textbf{Lecture 3 — Critique.} Quels sont les \cle{présupposés} (ce qui est admis sans preuve) ? Quelle est la \cle{portée} de la thèse (universelle ou limitée) ? Quelles \cle{objections} peut-on formuler ? Objectif : préparer la discussion.
\end{enumerate}
\end{methode}

\courssub{La grille d'évaluation type MINESEC}

L'exercice sur texte est noté sur 20 points selon une répartition officielle qu'il faut connaître par cœur, car elle dicte la stratégie de rédaction :

\begin{center}
\begin{tabularx}{\linewidth}{|l|X|c|}\hline
\rowcolor{popblueL} \textbf{Partie} & \textbf{Attendus} & \textbf{Barème} \\ \hline
Introduction & Présentation du texte (auteur, sujet), dégagement de la problématique, annonce du plan & 4 pts \\ \hline
Développement & Explication fidèle et ordonnée (thèse, concepts, arguments, structure) puis discussion critique argumentée & 10 pts \\ \hline
Conclusion & Bilan de l'explication, réponse à la problématique, ouverture maîtrisée & 4 pts \\ \hline
Qualité de la langue & Orthographe, syntaxe, vocabulaire philosophique, clarté & 2 pts \\ \hline
\textbf{Total} & & \textbf{20 pts} \\ \hline
\end{tabularx}
\end{center}

\begin{exemplebox}[Le calcul qui doit vous réveiller]
Sur 20 points, l'Introduction (4 pts) et la Conclusion (4 pts) valent à elles seules \textbf{8 points, soit 40 \% de la note}. Conséquence implacable : une introduction ratée (pas de problématique) ou une conclusion absente (faute de temps) signifie une note plafonnée autour de 12/20, \emph{même avec un bon développement}. Inversement, un candidat moyen au développement mais rigoureux sur l'introduction et la conclusion sécurise déjà près de la moitié des points. Stratégie de gestion du temps : réservez impérativement les 20 dernières minutes à la conclusion, quoi qu'il arrive.
\end{exemplebox}

\courssub{Les trois types de consignes au Bac technique}

Toutes les consignes du Baccalauréat technique se ramènent à trois types. Savoir identifier le type dès la première minute est décisif, car chacun commande un plan différent.

\begin{itemize}
\item \textbf{Type A — Thèse + Explication.} Consigne typique : « Dégagez la thèse du texte et expliquez-la » ou « Dégagez la thèse et montrez comment l'auteur l'établit ». Plan attendu : I. La thèse et sa signification ; II. Les arguments qui la fondent ; III. Discussion de sa portée.
\item \textbf{Type B — Questions guidées.} La consigne se présente comme une série de questions (souvent trois) portant sur le texte : « a) Quelle est la thèse de l'auteur ? b) Sur quel argument repose-t-elle ? c) Cette position vous paraît-elle défendable ? » Plan attendu : suivre l'ordre des questions, qui correspondent presque toujours à thèse / argumentation / discussion. Chaque question = une partie.
\item \textbf{Type C — Commentaire libre.} Consigne typique : « Commentez ce texte » ou « Expliquez et discutez ce texte ». C'est le type le plus exigeant : c'est au candidat de construire lui-même la problématique et le plan complet (explication ordonnée puis discussion), sans rampe fournie.
\end{itemize}

\begin{popfigure}
\begin{center}
\begin{tikzpicture}[node distance=1.5cm, every node/.style={font=\small}]
\node[draw, rectangle, rounded corners, fill=popblueL] (start) {Texte Bac};
\node[draw, diamond, aspect=2, below=of start, fill=popgoldL] (q1) {Type de consigne ?};
\node[draw, rectangle, rounded corners, fill=popgreenL, right=of q1, xshift=1.5cm] (a) {Type A : Thèse + Explication};
\node[draw, rectangle, rounded corners, fill=popgreenL, below=of q1] (b) {Type B : Questions guidées};
\node[draw, rectangle, rounded corners, fill=popgreenL, left=of q1, xshift=-1.5cm] (c) {Type C : Commentaire libre};
\draw[->] (start) -- (q1);
\draw[->] (q1) -- node[above] {Identifie} (a);
\draw[->] (q1) -- node[right] {Identifie} (b);
\draw[->] (q1) -- node[left] {Identifie} (c);
\end{tikzpicture}
\end{center}
\legende{Arbre de décision : identifier le type de consigne avant toute rédaction. Cette identification, faite dans les cinq premières minutes, détermine le plan de toute la copie.}
\end{popfigure}

\begin{exoresolu}[Identifier le type de consigne]
\textbf{Énoncé.} Voici trois consignes inspirées de sujets récents. Pour chacune, indiquez le type (A, B ou C) et le plan qu'elle commande.

1. « Dégagez la thèse de ce texte et montrez comment l'auteur la justifie. »

2. « Expliquez et discutez ce texte. »

3. « a) Quel problème l'auteur pose-t-il ? b) Quelle solution propose-t-il ? c) Cette solution vous semble-t-elle acceptable ? »

\tcblower
\textbf{Solution détaillée.}

1. \textbf{Type A} (Thèse + Explication). La consigne nomme explicitement la thèse et sa justification. Plan : I. La thèse (ce que l'auteur affirme et ce que cela signifie) ; II. L'argumentation (les étapes et preuves du raisonnement) ; III. Discussion critique (portée et limites).

2. \textbf{Type C} (Commentaire libre). Aucune question n'est fournie : le candidat doit tout construire. Plan : introduction avec problématique dégagée par le candidat ; développement en explication ordonnée (thèse, concepts, arguments) puis discussion ; conclusion.

3. \textbf{Type B} (Questions guidées). Les trois questions tracent le plan : la question a) correspond à la problématique/thèse, la b) à l'argumentation, la c) à la discussion. Il suffit de suivre cet ordre, en rédigeant une partie par question, avec des transitions.
\end{exoresolu}

\begin{savaistu}[Le texte au Bac technique : plus court, mais plus dense]
Au Baccalauréat technique, les textes proposés sont souvent plus courts (15 à 25 lignes) qu'en série générale, mais ils exigent une précision accrue sur les concepts du programme technique : le \cle{Travail}, la \cle{Technique}, la \cle{Société}, l'État, la Justice, la Liberté. Les auteurs fréquemment proposés (Hountondji, Marx, Bergson, Kant) sont choisis parce que leurs analyses éclairent directement le monde du travail et de la production que le futur technicien va intégrer. Maîtriser l'exercice sur texte, c'est donc aussi se donner les moyens de lire, demain, un cahier des charges, un code du travail ou un contrat : la même rigueur d'analyse s'applique.
\end{savaistu}

\begin{aretenir}[L'essentiel de la section]
\begin{itemize}
\item Les rapports de jury sanctionnent quatre erreurs : paraphrase, hors-sujet, absence de problématique, plan invisible.
\item L'exercice sur texte = \textbf{explication} (fidélité, depuis l'intérieur) + \textbf{discussion} (évaluation, depuis l'extérieur). Ni résumé, ni dissertation.
\item Méthode des trois lectures : découverte, analyse, critique.
\item Barème : 4 + 10 + 4 + 2. Introduction et conclusion = 40 \% de la note.
\item Trois types de consignes : A (thèse + explication), B (questions guidées), C (commentaire libre). Identifier le type, c'est déjà construire le plan.
\end{itemize}
\end{aretenir}

\coursec{L'Atelier Problématisation : Du Sujet à la Question Directrice}

Dans la section précédente, nous avons diagnostiqué les erreurs fréquentes des copies de Terminale technique et rappelé les exigences officielles de l'exercice sur texte philosophique : fidélité au texte, rigueur de l'analyse, clarté de la rédaction. Mais toutes ces exigences reposent sur un acte fondateur, celui qui décide du destin de toute la copie : la \cle{problématisation}. Un élève peut connaître parfaitement Kant ou Foucault, s'il n'a pas dégagé la bonne question directrice, sa copie sera hors sujet. Cette section est donc un atelier pratique : nous allons apprendre, pas à pas, à transformer un texte en question, puis cette question en plan.

\courssub{Repérer la tension interne du texte}

\textbf{L'intuition d'abord.} Un texte philosophique n'est jamais un simple exposé d'informations comme une notice de machine ou un bulletin météo. Il est né d'un \emph{conflit} : deux idées qui s'affrontent, deux valeurs qui se contredisent, une évidence qui se fissure. Le philosophe écrit parce que quelque chose \emph{ne va pas de soi}. Trouver la problématique, c'est donc d'abord repérer cette faille, cette \cle{tension interne} qui fait avancer le texte.

\begin{definition}[La tension interne]
La \cle{tension interne} d'un texte philosophique est l'opposition de deux concepts, de deux valeurs ou de deux exigences que l'auteur cherche à articuler, à hiérarchiser ou à dépasser. Elle se repère par des marqueurs linguistiques : \emph{mais}, \emph{cependant}, \emph{pourtant}, \emph{si... alors}, \emph{d'un côté... de l'autre}, ou par l'alternance de concessions et de réfutations.
\end{definition}

\textbf{Exercice pratique sur trois extraits courts.} Appliquons ce repérage à trois textes d'auteurs étudiés dans l'unité.

\begin{exemplebox}[Trois tensions à repérer]
\textbf{Extrait 1 (Kant, morale) :} \og Agir par devoir, c'est agir par respect de la loi morale, et non par inclination ; car une action faite par intérêt, même conforme au devoir, n'a pas de valeur morale authentique. \fg{} \par
\emph{Tension repérée :} \textbf{Devoir vs Inclination}. Kant oppose la motivation pure (le respect de la loi) à la motivation intéressée (le désir, le plaisir, le calcul).\par\medskip
\textbf{Extrait 2 (Mbembe, politique africaine) :} \og Les sociétés africaines ne sauraient se penser en rupture totale avec leurs traditions, ni s'y enfermer ; il leur faut inventer une modernité qui leur soit propre. \fg{} \par
\emph{Tension repérée :} \textbf{Tradition vs Modernité}. Ni rejet ni enfermement : l'auteur cherche une voie d'appropriation.\par\medskip
\textbf{Extrait 3 (Foucault, pouvoir) :} \og La société disciplinaire surveille et punit au nom de la sécurité collective ; mais cette surveillance, en s'étendant, menace les libertés qu'elle prétend protéger. \fg{} \par
\emph{Tension repérée :} \textbf{Sécurité vs Liberté}. Le moyen (surveiller) risque de détruire la fin (protéger des êtres libres).
\end{exemplebox}

Remarque décisive : dans chaque cas, la tension n'est pas un prétexte pour détruire l'un des deux termes. Le philosophe ne dit jamais simplement \og il faut la liberté et pas la sécurité \fg{}. Il \emph{travaille} l'opposition. C'est exactement ce travail que votre problématique doit annoncer.

\courssub{La technique du \og Pourquoi ? / Comment ? \fg{} : transformer la thèse en question-problème}

\textbf{L'intuition.} Beaucoup d'élèves confondent la \cle{thèse} du texte (ce que l'auteur affirme) et la \cle{problématique} (la question à laquelle cette thèse répond). Or une thèse est une \emph{réponse} ; la problématique est la \emph{question}. Pour remonter de la réponse à la question, il suffit de se demander : \og Pourquoi l'auteur a-t-il besoin d'affirmer cela ? Contre quelle difficulté se bat-il ? \fg{}

\begin{methode}[Problématiser : trouver la question cachée]
Problématiser, c'est trouver la \emph{question cachée} que le texte résout. Démarche en quatre étapes :
\begin{enumerate}
\item \textbf{Dégager la thèse} : formuler en une phrase ce que l'auteur soutient (ex. : \og La technique libère l'homme \fg{}).
\item \textbf{Identifier la tension} : nommer les deux concepts en conflit (ex. : maîtrise de la nature vs asservissement de l'homme).
\item \textbf{Appliquer la formule type} : \og Comment [Concept A] peut-il [Verbe] [Concept B] sans [Risque/Contradiction] ? \fg{}
\item \textbf{Vérifier} : relire le texte et contrôler que chaque paragraphe apporte un élément de réponse à cette question. Si un paragraphe reste orphelin, la problématique est fausse ou incomplète.
\end{enumerate}
\end{methode}

\begin{exemplebox}[Application de la formule type]
\textbf{Thèse relevée :} \og La technique libère l'homme. \fg{}\par
\textbf{Tension :} libération (maîtrise de la nature, gain de temps) vs aliénation (dépendance aux machines, rythme imposé).\par
\textbf{Problématique construite :} \og Dans quelle mesure la technique, en maîtrisant la nature, ne risque-t-elle pas d'asservir l'homme qu'elle prétend libérer ? \fg{}\par\medskip
Autre exemple : thèse \og L'État garantit la paix sociale \fg{} $\rightarrow$ problématique : \og Comment l'État peut-il garantir la sécurité des citoyens sans porter atteinte à la liberté qui fonde sa légitimité ? \fg{}
\end{exemplebox}

\begin{attention}[Piège : la question trop large ou trop étroite]
Deux erreurs symétriques détruisent la problématique. \textbf{Question trop large} : \og Qu'est-ce que la technique ? \fg{} — elle ne rend pas compte de la tension spécifique du texte et ouvre un sujet infini. \textbf{Question trop étroite} : \og La machine à laver libère-t-elle la ménagère ? \fg{} — elle réduit la portée philosophique à un cas particulier. La bonne problématique épouse \emph{exactement} le texte : ni plus, ni moins.
\end{attention}

\courssub{Construire un plan déduit du texte}

\textbf{L'intuition.} Le plan n'est pas un corset que l'on impose au texte de l'extérieur ; c'est le \emph{squelette} du texte que l'on met à nu. L'auteur a raisonné en étapes : votre développement doit suivre ces étapes. Autrement dit, les grandes phases du raisonnement de l'auteur deviennent les parties de votre devoir.

\begin{definition}[Plan dialectique et plan analytique]
Le \cle{plan dialectique} suit le schéma Thèse / Antithèse / Synthèse : on expose d'abord une position, puis la position opposée, puis un dépassement. Le \cle{plan analytique} suit le fil du texte : Argument 1 / Argument 2 / Argument 3, en respectant l'ordre et la logique propres de l'auteur. Au Baccalauréat technique, le plan analytique \emph{rigoureux} est souvent plus sûr, car il garantit la fidélité au texte, première exigence des correcteurs.
\end{definition}

\begin{attention}[Le piège du \og plan plaqué \fg]
Imposer un plan Thèse/Antithèse/Synthèse à un texte qui suit une logique linéaire (déduction rigoureuse, narration, description progressive) déforme la pensée de l'auteur : c'est le \og plan plaqué \fg. De même, le \og plan catalogue \fg{} (I. L'auteur, II. Le contexte historique, III. Mon avis) est \textbf{interdit} : il ne répond à aucune problématique et ignore le texte. Règle d'or : \emph{respecter la logique de l'auteur}, qu'elle soit dialectique ou linéaire.
\end{attention}

Voici, pour notre exemple sur la technique, la visualisation d'un plan entièrement \emph{déduit du texte}, chaque partie renvoyant aux paragraphes précis :

\begin{popfigure}
\begin{center}
\begin{tikzpicture}[grow=right, level distance=3.5cm, sibling distance=1.5cm, every node/.style={draw, rectangle, rounded corners, fill=popcream, font=\small, align=center, text width=3.2cm}]
\node[fill=popblueL, text width=4cm] {Problématique : \textit{La technique libère-t-elle ou aliène-t-elle ?}}
child {node {I. La technique comme libération (maîtrise nature, gain temps)} child {node {\S 1--2 : outils, production}} child {node {\S 3 : condition d'émancipation}} }
child {node {II. La technique comme aliénation (dépendance, bureaucratie)} child {node {\S 4--5 : machine, rythme imposé}} child {node {\S 6 : perte d'autonomie}} }
child {node {III. Dépassement : maîtrise critique (éthique, appropriation)} child {node {\S 7 : éducation technique}} child {node {\S 8 : choix politique}} };
\end{tikzpicture}
\end{center}
\legende{Arbre de problématisation : la problématique centrale se ramifie en trois parties, chacune ancrée dans des paragraphes précis du texte (\S = numéro de paragraphe). Chaque branche du plan correspond à une étape du raisonnement de l'auteur.}
\end{popfigure}

\begin{savaistu}[Ce que cherchent les correcteurs MINESEC]
Les correcteurs MINESEC cherchent d'abord la \emph{trace du texte} dans la copie : citations précises entre guillemets, numéros de paragraphes ou de lignes, reprise des expressions de l'auteur. Une copie sans référence textuelle explicite est sanctionnée, même si elle est brillante. Citer le texte, ce n'est pas le réciter : c'est prouver qu'on l'a lu.
\end{savaistu}

\courssub{Atelier \og Plan en 5 min \fg{} : la gestion de l'eau au Cameroun}

Passons à l'entraînement chronométré, dans les conditions réelles de l'examen. Lisons le texte fictif suivant, de type philosophique.

\begin{exemplebox}[Texte d'atelier (texte philosophique fictif)]
\og L'eau, disent les anciens, appartient à tous et à personne. Pourtant, dans nos villes, de Yaoundé à Douala, l'accès à l'eau potable est devenu une marchandise que l'on achète, et que les plus pauvres ne peuvent plus se payer. [\S 1] Sans doute la gestion technique de l'eau — forages, châteaux d'eau, réseaux de distribution — exige-t-elle des investissements que seule une tarification peut financer ; traiter l'eau comme un bien gratuit, c'est la condamner au gaspillage et à la pénurie. [\S 2] Mais lorsque le prix devient la condition de l'accès, c'est le droit à la vie lui-même qui se trouve mis en vente, car nul ne peut vivre sans eau. [\S 3] La solution ne peut donc être ni la gratuité absolue, qui ruine l'entretien des infrastructures, ni le marché pur, qui exclut les démunis. [\S 4] Il appartient à la puissance publique de garantir à chacun un volume vital d'eau, tout en faisant payer la consommation superflue : ainsi la justice et l'efficacité cessent d'être ennemies. [\S 5] \fg{}
\end{exemplebox}

\begin{exoresolu}[Deux plans dialectiques distincts, et leur comparaison]
\textbf{Énoncé.} En 5 minutes, proposez deux plans dialectiques distincts pour ce texte, puis justifiez lequel est le meilleur.
\tcblower
\textbf{Étape 1 — Problématique.} Tension : \textbf{Justice (droit à l'eau pour tous) vs Efficacité (financement technique des infrastructures)}. Problématique : \og Comment garantir à tous l'accès à l'eau, bien vital, sans compromettre la gestion technique et financière qui rend cet accès possible ? \fg{}\par\medskip
\textbf{Plan A (dialectique classique) :}
\begin{itemize}
\item I. L'eau, bien commun, devrait être gratuite pour tous (\S 1, \S 3 : le droit à la vie).
\item II. Mais la gratuité absolue ruine les infrastructures et encourage le gaspillage (\S 2, \S 4).
\item III. Dépassement : la tarification différenciée concilie justice et efficacité (\S 4--5).
\end{itemize}
\textbf{Plan B (analytique, fil du texte) :}
\begin{itemize}
\item I. Le constat : la marchandisation de l'eau menace le droit vital (\S 1).
\item II. L'examen des deux impasses : le marché pur exclut les pauvres, la gratuité pure détruit le service (\S 2--4).
\item III. La solution de l'auteur : un volume vital garanti par l'État, le superflu payant (\S 5).
\end{itemize}
\textbf{Justification du meilleur plan.} Le Plan B est supérieur ici, pour trois raisons. (1) \emph{Fidélité} : il suit exactement l'ordre du texte (constat $\rightarrow$ impasses $\rightarrow$ solution), donc aucun paragraphe n'est déplacé. (2) \emph{Exhaustivité} : le \S 1, simple constat, trouve sa place dans la partie I, alors que le Plan A le dilue. (3) \emph{Honnêteté dialectique} : le texte ne présente pas d'abord une thèse \og gratuité \fg{} puis une antithèse \og marché \fg{} ; il les examine \emph{ensemble} comme deux impasses symétriques (\S 4 : \og ni... ni \fg). Le Plan A plaque un schéma Thèse/Antithèse sur un texte qui fonctionne par élimination conjointe. Conclusion : le Plan B, analytique et rigoureux, respecte la logique de l'auteur — c'est le plan à rédiger.
\end{exoresolu}

\courssub{Valider son plan : exhaustivité, progressivité, unité}

Un plan, même séduisant, doit passer trois contrôles avant la rédaction. C'est le \og contrôle qualité \fg{} de votre brouillon, comme un technicien vérifie une installation avant la mise en service.

\begin{methode}[Les trois critères de validation d'un plan]
\begin{enumerate}
\item \textbf{Exhaustivité} : tout le texte est-il couvert ? Numérotez les paragraphes et attribuez chacun à une partie. Si un paragraphe ne trouve sa place nulle part, le plan est incomplet ; si un paragraphe doit être tordu pour entrer dans une partie, le plan est faux.
\item \textbf{Progressivité} : va-t-on du simple au complexe, du constat à la solution, de la thèse la plus faible à la plus forte ? L'ordre des parties doit être nécessaire : on ne peut pas intervertir la partie I et la partie III sans détruire le raisonnement.
\item \textbf{Unité} : chaque partie répond-elle à la problématique ? Ajoutez mentalement à la fin de chaque intitulé : \og ...ce qui répond à notre question en montrant que... \fg. Si la phrase est impossible à compléter, la partie est hors sujet.
\end{enumerate}
\end{methode}

\begin{exemplebox}[Validation du Plan B de l'atelier]
\textbf{Exhaustivité :} \S 1 $\rightarrow$ partie I ; \S 2--4 $\rightarrow$ partie II ; \S 5 $\rightarrow$ partie III. Aucun paragraphe orphelin. Validé.\par
\textbf{Progressivité :} constat (I) $\rightarrow$ analyse des impasses (II) $\rightarrow$ solution (III). Ordre nécessaire : on ne peut exposer la solution avant d'avoir montré les impasses. Validé.\par
\textbf{Unité :} I montre \emph{pourquoi} la question se pose ; II montre \emph{pourquoi} les réponses simples échouent ; III montre \emph{comment} concilier justice et efficacité. Chaque partie répond à la problématique. Validé.
\end{exemplebox}

\begin{exemplebox}[Gestion du temps : le chronomètre de l'épreuve]
La problématisation et la construction du plan s'inscrivent dans un temps d'épreuve limité. Temps idéal au Bac technique (1 h 30) :
\begin{center}
\begin{tabularx}{\linewidth}{|l|X|c|}
\hline
\rowcolor{popblueL} \textbf{Phase} & \textbf{Travail réalisé} & \textbf{Durée} \\
\hline
Lecture / Analyse & Repérage de la tension, problématique, plan détaillé au brouillon & 25 min \\
\hline
Rédaction de l'introduction & Amorce, présentation du texte, problématique, annonce du plan & 10 min \\
\hline
Développement & Rédaction des trois parties, citations du texte, explications & 45 min \\
\hline
Conclusion & Bilan de la réponse, ouverture maîtrisée & 10 min \\
\hline
Relecture & Orthographe, citations, cohérence avec la problématique & 10 min \\
\hline
\end{tabularx}
\end{center}
Respecter ce chronomètre est un savoir-faire évalué : une problématique brillante trouvée à la 50$^{\text{ème}}$ minute ne sert à rien si le développement reste à l'état de brouillon.
\end{exemplebox}

\begin{aretenir}[L'essentiel de l'atelier]
\begin{itemize}
\item Tout texte philosophique naît d'une \cle{tension interne} (Devoir/Inclination, Tradition/Modernité, Sécurité/Liberté) : la repérer est le premier geste.
\item La problématique est la \emph{question cachée} dont la thèse est la réponse ; formule type : \og Comment [A] peut-il [verbe] [B] sans [risque] ? \fg
\item Le plan se \emph{déduit} du texte : les étapes du raisonnement de l'auteur deviennent vos parties. Ni plan catalogue, ni plan plaqué.
\item Avant de rédiger, valider le plan par trois contrôles : \cle{exhaustivité}, \cle{progressivité}, \cle{unité}.
\end{itemize}
\end{aretenir}

\begin{exercice}[À vous de jouer]
Soit l'extrait suivant : \og La tradition transmet les valeurs sans lesquelles une communauté se dissout ; mais une tradition qui interdit toute critique devient une prison, et une jeunesse qui rompt avec tout héritage devient orpheline d'elle-même. \fg
\begin{enumerate}
\item Repérez la tension interne et nommez les deux concepts en conflit.
\item Construisez la problématique avec la formule type.
\item Proposez un plan en deux ou trois parties déduit du texte, puis validez-le par les trois critères.
\end{enumerate}
\end{exercice}

\begin{corrige}[Exercice]
\begin{enumerate}
\item Tension : \textbf{Tradition vs Liberté critique} (ou Héritage vs Innovation). Le texte oppose la fonction protectrice de la tradition (transmission des valeurs) à son danger (la prison), puis équilibre par le danger symétrique (la rupture totale rend orphelin).
\item Problématique : \og Comment la tradition peut-elle transmettre à la jeunesse les valeurs qui fondent la communauté sans devenir une prison qui étouffe la critique et la liberté ? \fg
\item Plan possible (analytique) : I. La tradition est nécessaire : elle transmet les valeurs qui font vivre la communauté. II. Mais elle devient dangereuse lorsqu'elle interdit la critique, et la rupture totale est tout aussi dangereuse. III. (Synthèse implicite du texte) Il faut donc une tradition vivante, ouverte à la critique. Validation : exhaustivité (les trois moments de la phrase sont couverts), progressivité (nécessité $\rightarrow$ dangers $\rightarrow$ voie moyenne), unité (chaque partie éclaire la question de la transmission sans enfermement).
\end{enumerate}
\end{corrige}

Nous savons désormais dégager une problématique et construire un plan fidèle. Reste à transformer cette architecture en rédaction effective : c'est l'objet de la section suivante, consacrée à la rédaction chirurgicale de l'introduction, du développement et de la conclusion.

\coursec{Rédaction Chirurgicale : Introduction, Développement, Conclusion}

Dans la section précédente, nous avons appris à transformer un texte philosophique en une \cle{question directrice} précise : la problématique. Mais une bonne problématique gardée au brouillon ne rapporte aucun point. Il faut maintenant la \emph{mettre en mots}, c'est-à-dire passer de l'analyse silencieuse à la \cle{rédaction} proprement dite. Cette étape est décisive : au Baccalauréat technique, le correcteur lit des centaines de copies ; il juge en quelques minutes si l'élève maîtrise les codes de l'exercice. Nous allons donc construire, pièce par pièce, une copie « chirurgicale » : une introduction en entonnoir, un développement articulé en paragraphes d'analyse, et une conclusion qui boucle sans trahir.

\courssub{L'introduction en entonnoir : du large vers le précis}

L'introduction philosophique fonctionne comme un \cle{entonnoir} : on part d'un constat large (la vie humaine, l'actualité) et on resserre progressivement jusqu'à l'annonce du plan, point le plus étroit et le plus précis. Chaque étape est obligatoire et obéit à un ordre strict.

\begin{popfigure}
\begin{center}
\begin{tikzpicture}[node distance=0.9cm, every node/.style={align=center, font=\small}]
\node[draw, rectangle, fill=poppinkL, text width=8cm, rounded corners] (i1) {1. Accroche : fait social, actualité camerounaise ou mondiale};
\node[draw, rectangle, fill=poporangeL, text width=7.2cm, below of=i1, rounded corners] (i2) {2. Présentation : auteur, \oe uvre, date, contexte intellectuel};
\node[draw, rectangle, fill=popgoldL, text width=6.4cm, below of=i2, rounded corners] (i3) {3. Thèse : reformulée \emph{avec les mots de l'auteur}};
\node[draw, diamond, fill=popgreenL, below of=i3, aspect=2.4, inner sep=1pt, align=center] (i4){4. Problématique :\\ question directrice};
\node[draw, rectangle, fill=popblueL, text width=6.4cm, below of=i4, rounded corners] (i5) {5. Annonce du plan : « Nous verrons d'abord\ldots\\ Ensuite\ldots Enfin\ldots »};
\draw[->, thick, popdark] (i1) -- (i2);
\draw[->, thick, popdark] (i2) -- (i3);
\draw[->, thick, popdark] (i3) -- (i4);
\draw[->, thick, popdark] (i4) -- (i5);
\end{tikzpicture}
\legende{La structure en entonnoir de l'introduction : cinq étapes ordonnées, du plus large (l'accroche) au plus précis (l'annonce du plan). Aucune étape ne peut être sautée ni inversée.}
\end{center}
\end{popfigure}

Détaillons chaque marche de l'entonnoir.

\textbf{1. L'accroche contextuelle.} Elle part d'un fait que tout le monde connaît et qui touche au thème du texte. Pour un texte sur le travail, on peut partir de la réalité camerounaise : « À Douala comme à Yaoundé, des milliers de jeunes quittent chaque matin les quartiers populaires pour gagner leur pain dans le secteur informel : le travail semble être la condition première de la dignité humaine. » L'accroche doit être \emph{vraie}, \emph{brève} (deux ou trois phrases) et \emph{en lien réel} avec le texte. Évitez les formules creuses du type « Depuis la nuit des temps, l'homme\ldots » qui ne disent rien.

\textbf{2. La présentation précise du texte.} Nommez l'auteur, l'\oe uvre, la date et le contexte intellectuel : « C'est à cette interrogation que répond Hannah Arendt (1906-1975), philosophe allemande exilée aux États-Unis, dans \emph{Condition de l'homme moderne} (1958), ouvrage où elle repense les activités fondamentales de la vie humaine après les catastrophes du XX\textsuperscript{e} siècle. » Cette étape montre au correcteur que vous avez identifié le texte et que vous savez situer une pensée dans son époque.

\textbf{3. La thèse formulée avec ses mots.} Il ne suffit pas de dire « l'auteur parle du travail ». Il faut restituer la \cle{thèse}, c'est-à-dire la position défendue, en citant les termes mêmes du texte : « Arendt soutient que le \emph{travail} (\emph{labor}) doit être distingué de l'\emph{\oe uvre} (\emph{work}) : le premier ne fait que reproduire la vie biologique, tandis que la seconde fabrique un monde durable. »

\textbf{4. La problématique.} Reprenez la question directrice construite dans la section précédente, sous forme interrogative : « Dès lors, faut-il opposer, comme Arendt, le travail qui asservit l'homme à la nécessité et l'\oe uvre qui l'élève à la liberté ? Le travail est-il seulement une malédiction biologique ou peut-il aussi humaniser ? »

\textbf{5. L'annonce de plan signifiée.} Le verbe « signifier » est essentiel : il ne suffit pas de dire « nous verrons la première partie puis la deuxième ». Il faut \emph{donner le contenu} de chaque moment : « Nous verrons d'abord qu'Arendt distingue rigoureusement le travail de l'\oe uvre ; nous montrerons ensuite que cette distinction condamne le travail à la répétition servile ; nous demanderons enfin si le travail, même contraint, ne participe pas à la formation de l'homme. »

\begin{methode}[L'introduction est un contrat moral avec le correcteur]
L'annonce du plan vous \emph{engage} : si le plan annoncé n'est pas suivi dans le développement, la copie perd toute cohérence et la note chute lourdement. D'où la règle d'or pratique :
\begin{enumerate}
\item Analyser le texte et construire la problématique (section 2).
\item Rédiger le \textbf{plan détaillé au brouillon} (parties, sous-parties, citations choisies).
\item Rédiger l'introduction \textbf{en dernier}, une fois le plan fixé, pour garantir que l'annonce correspond exactement à ce qui sera réellement traité.
\end{enumerate}
Une introduction écrite « à l'aveugle » avant le plan promet presque toujours ce que la copie ne tiendra pas.
\end{methode}

\courssub{Le développement : la règle du paragraphe unique}

Le développement est le corps de la copie. Sa loi fondamentale est simple : \cle{une idée = un paragraphe}. Dès que vous changez d'idée, vous changez de paragraphe (alinéa visible). Un paragraphe qui contient deux idées brouille l'analyse ; deux paragraphes qui répètent la même idée signalent un vide.

\begin{propriete}[Structure interne du paragraphe d'analyse]
Un paragraphe d'explication de texte bien construit enchaîne cinq moments :
\begin{enumerate}
\item \textbf{Phrase-thèse du paragraphe} : l'idée directrice, annoncée en une phrase claire.
\item \textbf{Citation courte intégrée} : un extrait du texte entre guillemets, inséré dans votre propre phrase.
\item \textbf{Explication} : analyse du vocabulaire et des concepts (définir les mots-clés de la citation).
\item \textbf{Monstration de l'enchaînement logique} : montrer comment cette idée s'articule à la précédente et fait avancer l'argumentation de l'auteur.
\item \textbf{Lien avec la problématique et transition} : rappeler ce que cette idée apporte à la question directrice et annoncer la suite.
\end{enumerate}
Cette structure est exigible dans toute copie d'explication de texte au Baccalauréat technique.
\end{propriete}

\begin{exemplebox}[Un paragraphe d'analyse réussi (texte d'Arendt sur le travail)]
« Arendt établit d'abord une distinction fondatrice entre deux activités que le langage courant confond. Elle écrit : « le travail laisse derrière lui des objets d'usage » qui se consomment, tandis que l'\oe uvre produit des choses qui « demeurent ». Le terme « consommer » désigne ici la destruction par l'usage : le pain fabriqué le matin est mangé le soir, et le labeur du boulanger doit recommencer chaque jour. Le travail est donc enfermé dans le cycle de la vie biologique, à la différence de l'\oe uvre qui installe un monde stable. Cette distinction, première étape de l'argumentation, fonde la thèse selon laquelle le travail ne libère pas l'homme : c'est ce qu'il nous faut maintenant examiner. »
\end{exemplebox}

\begin{attention}[La paraphrase déguisée, faute éliminatoire]
L'erreur la plus fréquente consiste à \emph{répéter} le texte au lieu de l'\emph{analyser} : « L'auteur dit que le travail consomme. Ensuite il ajoute que l'\oe uvre demeure. Puis il dit que\ldots » C'est de la paraphrase déguisée : le correcteur y voit l'absence de pensée personnelle. Remplacez les verbes de parole par des \textbf{verbes d'argumentation} et montrez la démarche : « L'auteur \emph{argumente} que\ldots \emph{en s'appuyant sur}\ldots \emph{afin de montrer que}\ldots ». Vous ne racontez pas le texte, vous en démontez le mécanisme.
\end{attention}

\courssub{La gestion des citations : intégrer, ne pas parachuter}

Une citation n'a de valeur que si elle est \cle{intégrée syntaxiquement} à votre phrase et \emph{expliquée} juste après. Une citation « parachutée » (jetée seule entre deux paragraphes, sans introduction ni commentaire) est sanctionnée : elle prouve que vous avez recopié sans comprendre.

Règles techniques :
\begin{itemize}
\item \textbf{Intégration} : la citation s'insère dans votre phrase — Arendt affirme que « la vie est le bien suprême » du \emph{laborans} — et non l'inverse.
\item \textbf{Coupure} : utilisez les crochets avec points de suspension [\ldots] pour élaguer un passage trop long sans le déformer.
\item \textbf{Modification} : tout mot adapté (accord, temps) se met entre crochets : « le travail [est] enfermé dans le cycle vital ». L'italique signale une mise en valeur que vous ajoutez ; la mention \emph{[sic]} signale une particularité présente dans l'original.
\item \textbf{Dosage} : les citations ne doivent pas dépasser \textbf{10 à 15 \%} du volume de la copie. Une copie qui serait un collier de citations n'est plus une copie de philosophie mais un montage.
\end{itemize}

\courssub{La conclusion : boucler sans trahir}

La conclusion comporte deux temps et deux interdits.

\textbf{Premier temps : le bilan synthétique.} Répondez explicitement à la problématique posée dans l'introduction, en rappelant le chemin parcouru : « Nous avons vu qu'Arendt sépare le travail, esclave de la nécessité vitale, de l'\oe uvre, productrice d'un monde durable. Mais cette séparation ne condamne pas le travailleur : en reconnaissant la dignité de chaque activité, Arendt invite à replacer le travail à sa juste mesure. »

\textbf{Second temps : l'ouverture pertinente.} Elle doit prolonger la réflexion, non l'abandonner : un autre auteur du programme (Marx, pour qui le travail est au contraire la voie d'humanisation), une actualité camerounaise (la valorisation de l'artisanat local, l'importation de la main-d'\oe uvre qualifiée), ou une autre notion du programme (la liberté, la technique).

\textbf{Interdits absolus} : aucun \textbf{nouvel argument} (tout ce qui devait être démontré l'a été dans le développement), aucune \textbf{citation nouvelle}, et jamais la formule scolaire « En conclusion je dirais que\ldots », qui infantilise la copie.

\begin{savaistu}[Peut-on dire « je » en philosophie ?]
L'usage du « je » est \textbf{toléré} dans les moments de discussion personnelle (troisième partie, conclusion) pour marquer l'appropriation critique : « Nous pouvons nous demander si\ldots », « Il me semble toutefois que\ldots ». En revanche, il est \textbf{banni de l'explication} du texte : pendant que vous exposez la pensée d'Arendt, c'est elle qui parle, pas vous. Confondre les deux registres est une faute de méthode classique.
\end{savaistu}

\begin{exemplebox}[Les bonnes proportions d'une copie (repères chiffrés)]
\begin{center}
\begin{tabularx}{\linewidth}{|l|X|X|}\hline
\rowcolor{popblueL} \textbf{Partie} & \textbf{Longueur indicative} & \textbf{Fonction} \\ \hline
Introduction & 15 à 20 lignes & Poser le contrat (entonnoir complet) \\ \hline
Développement & 1,5 à 2 pages (deux ou trois grandes parties) & Analyser, citer, discuter \\ \hline
Conclusion & 10 à 15 lignes & Bilan + ouverture \\ \hline
\end{tabularx}
\end{center}
Retenez : la \textbf{qualité prime sur la quantité}. Une page et demie dense, structurée, où chaque citation est expliquée, vaut infiniment mieux que trois pages floues et paraphrasantes.
\end{exemplebox}

\begin{exoresolu}[Rédiger une introduction complète (sujet type Bac : texte d'Arendt sur le travail)]
Énoncé : à partir du texte de Hannah Arendt, \emph{Condition de l'homme moderne} (1958), distinguant travail et \oe uvre, rédigez l'introduction complète d'une explication de texte.
\tcblower
\textbf{Solution détaillée.} Suivons l'entonnoir en cinq étapes.
\emph{Accroche} : « Chaque matin, des milliers de Camerounais se lèvent avant l'aube pour vendre, transporter, cultiver : le travail apparaît comme la condition même de la survie et de la dignité. Pourtant, ce travail qui occupe l'essentiel de nos jours est-il ce qui nous rend pleinement humains ? »
\emph{Présentation} : « C'est le sens de la réflexion de Hannah Arendt, philosophe allemande du XX\textsuperscript{e} siècle, dans \emph{Condition de l'homme moderne} (1958), où elle repense les activités humaines fondamentales. »
\emph{Thèse} : « Elle y soutient que le \emph{travail}, enfermé dans la répétition du cycle vital, se distingue radicalement de l'\emph{\oe uvre}, seule capable d'édifier un monde durable. »
\emph{Problématique} : « Faut-il alors dévaloriser le travail au profit de l'\oe uvre, ou le travail possède-t-il lui aussi une dignité propre ? »
\emph{Annonce} : « Nous verrons d'abord comment Arendt établit la distinction entre travail et \oe uvre ; nous montrerons ensuite qu'elle enferme le travail dans la nécessité biologique ; nous discuterons enfin la possibilité d'une dignité propre du travail. »
Cette introduction respecte l'ordre de l'entonnoir, cite les mots de l'auteur et annonce un plan dont chaque étape est signifiée.
\end{exoresolu}

\begin{experience}[TP : rédaction collaborative en ateliers (séance 2)]
\textbf{Matériel} : tableau, craies de couleurs, texte du sujet de Bac 2022 (Hannah Arendt, \emph{Le Travail}), grille d'évaluation de la section 1.
\begin{enumerate}
\item Former quatre groupes. Les groupes A et B rédigent chacun une \textbf{introduction complète} du sujet au brouillon (20 minutes) ; les groupes C et D rédigent chacun une \textbf{conclusion} (bilan + ouverture).
\item Un rapporteur par groupe écrit la production au tableau.
\item Correction collective : la classe note chaque production avec la grille officielle (accroche ? présentation ? thèse avec les mots de l'auteur ? problématique ? plan signifié ? pour l'introduction ; bilan ? ouverture ? absence de nouvel argument ? pour la conclusion).
\item Réécriture commune d'une version « copie modèle » intégrant les corrections.
\end{enumerate}
\textbf{Observations attendues} : les premières versions oublient presque toujours la présentation de l'auteur ou annoncent un plan non signifié ; les conclusions introduisent souvent un argument nouveau. La correction par la grille rend ces défauts visibles et corrigibles.
\end{experience}

\begin{aretenir}[L'essentiel de la section]
\begin{itemize}
\item L'introduction suit l'\textbf{entonnoir} : accroche, présentation, thèse avec les mots de l'auteur, problématique, annonce de plan signifiée. Elle se rédige \emph{après} le plan détaillé.
\item Le développement obéit à la règle \textbf{une idée = un paragraphe}, chaque paragraphe combinant phrase-thèse, citation intégrée, explication, enchaînement logique et transition.
\item Les citations s'\textbf{intègrent} syntaxiquement, se coupent par [\ldots], se modifient entre crochets, et restent sous 15 \% de la copie.
\item La conclusion fait le \textbf{bilan} de la problématique puis \textbf{ouvre} (autre auteur, actualité camerounaise, autre notion), sans argument ni citation nouveaux.
\end{itemize}
\end{aretenir}

La rédaction étant désormais maîtrisée dans sa structure, il reste à affronter les formes précises que prendront les questions le jour de l'examen : c'est l'objet de la section suivante, consacrée aux questions types et aux pièges spécifiques du Baccalauréat technique.

\coursec{Maîtrise des Questions Types \& Pièges Spécifiques Bac Tech}

Après avoir maîtrisé la rédaction chirurgicale de l'introduction, du développement et de la conclusion (Section 3), nous devons maintenant nous armer pour affronter l'épreuve elle-même : les questions posées par le jury. Au Bac Tech, l'exercice sur texte ne se résume pas à une dissertation guidée ; il impose un \textbf{parcours obligé} en trois questions types (Q1, Q2, Q3) qui testent des compétences distinctes : la compréhension rigoureuse (Q1), l'exégèse précise (Q2) et l'esprit critique philosophique (Q3). Connaître le "format" de chaque question, c'est déjà gagner la moitié des points. Cette section décortique chaque type, révèle les pièges officiels et propose un cas pratique intégral sur un thème brûlant pour le Cameroun : la corruption.

\courssub{Typologie Officielle des Trois Questions \& Répartition des Points}

\begin{definition}[Archétype de l'épreuve sur texte (Bac Tech / Probatoire)]
L'épreuve se présente invariablement sous la forme d'un texte philosophique (extraits d'auteurs au programme ou prolongement de notions du programme : Technique, Travail, État, Société, Morale, Politique) suivi de \textbf{trois questions obligatoires} :
\begin{enumerate}
    \item \textbf{Q1 (Compréhension / Synthèse) :} « Dégagez la thèse de l'auteur et les arguments principaux qu'il utilise pour la soutenir. »
    \item \textbf{Q2 (Analyse / Exégèse) :} « Expliquez le passage suivant : "..." » (Citation d'une phrase ou d'un paragraphe clé du texte).
    \item \textbf{Q3 (Discussion / Évaluation critique) :} « Le texte vous semble-t-il convaincant ? » ou « Discutez la thèse de l'auteur. » ou « Quelles sont les limites de cette thèse ? »
\end{enumerate}
\end{definition}

\begin{exemplebox}[Répartition indicative des points (sur 20)]
Le barème officiel n'est pas publié dans le détail, mais la pratique corrigée et les rapports de jury convergent vers cette pondération qui doit dicter votre \textbf{gestion du temps} (3h à 4h d'épreuve) :
\begin{center}
\begin{tabularx}{\linewidth}{|l|X|c|}\hline
\textbf{Question} & \textbf{Compétence évaluée} & \textbf{Points estimés} \\ \hline
Q1 & Restitution structurée (Thèse + Args) & 4 à 5 pts \\ \hline
Q2 & Précision conceptuelle + Fonction argumentative & 5 à 6 pts \\ \hline
Q3 & Esprit critique + Culture philosophique + Argumentation personnelle & 8 à 10 pts \\ \hline
\end{tabularx}
\end{center}
\emph{Conséquence stratégique :} Ne passez pas 1h30 sur la Q1. La Q3 \emph{porte la note}. Si le temps manque, sacrifiez l'ouverture de la conclusion de la Q3, \textbf{jamais} la structure argumentée de son développement (thèses/antithèses/synthèse).
\end{exemplebox}

\courssub{Question 1 : La Synthèse Structurée (Thèse + Arguments)}

\begin{methode}[Réussir la Q1 : Le protocole "T.A.A." (Thèse - Arguments - Articulation)]
\begin{enumerate}
    \item \textbf{Identifiez la Thèse unique :} C'est la réponse directe de l'auteur au problème posé. Formulez-la en \textbf{une phrase affirmative complète} (sujet + verbe + complément). Évitez : « L'auteur parle de... » Préférez : « L'auteur soutient que la liberté n'est pas l'indépendance mais l'obéissance à une loi qu'on s'est prescrite (Rousseau). »
    \item \textbf{Extrayez 2 à 3 Arguments principaux :} Numérotez-les explicitement (\textbf{Argument 1}, \textbf{Argument 2}...). Un argument = une idée-force + son développement logique (pas une simple citation).
    \item \textbf{Montrez l'Articulation (Le "Fil rouge") :} Expliquez \emph{comment} les arguments s'enchaînent pour prouver la thèse. (Ex: « L'arg. 1 pose le principe ontologique, l'arg. 2 en tire la conséquence politique, l'arg. 3 répond à l'objection du relativisme. »)
\end{enumerate}
\end{methode}

\begin{attention}[Pièges Mortels Q1]
\begin{itemize}
    \item \textbf{Le "Catalogue de citations" :} Recopier des bouts de texte sans les reformuler ni les relier. \rightarrow \textbf{Zéro point} pour l'argumentation.
    \item \textbf{L'avis personnel :} « Je suis d'accord car... » \rightarrow \textbf{Hors sujet}. La Q1 est une \emph{restitution}, pas une discussion.
    \item \textbf{La thèse "parapluie" :} « L'auteur parle de la justice. » \rightarrow Trop vague. La thèse doit être \emph{la position spécifique} de l'auteur sur la justice.
    \item \textbf{Oublier l'articulation :} Lister des arguments sans montrer la logique du texte. Perte de 1 à 2 points.
\end{itemize}
\end{attention}

\courssub{Question 2 : L'Exégèse Locale (Explication de Passage)}

C'est l'exercice le plus technique. On vous demande d'éclairer une zone précise du texte. C'est ici que se joue la maîtrise des \cle{concepts}.

\begin{methode}[Protocole "S.D.R.F." pour l'Explication de Passage (Q2)]
\begin{enumerate}
    \item \textbf{Situer (Contexte) :} Début / Milieu / Fin du texte ? Rupture ou continuité ? Répond à quelle objection ? Prépare quelle conclusion ?
    \item \textbf{Définir (Concepts \emph{in situ}) :} Identifiez les termes techniques (ex: \emph{aliénation}, \emph{survaleur}, \emph{contractualisme}, \emph{graissage}). Donnez leur \textbf{sens philosophique précis dans la pensée de l'auteur}, pas le sens courant du dictionnaire.
    \item \textbf{Reformuler (Paraphrase) :} Réécrivez le passage \emph{avec vos mots}, en explicitant les implicites (références historiques, auteurs visés, concepts sous-jacents).
    \item \textbf{Fonction argumentative (Rôle) :} À \emph{quoi sert} ce passage dans la démonstration globale ? (Prouve l'argument 2 ? Réfute un contre-argument ? Illustre la thèse par un exemple ?)
\end{enumerate}
\end{methode}

\begin{definition}[Exégèse vs Paraphrase]
\textbf{Paraphrase} = Dire la même chose avec d'autres mots (niveau scolaire). \textbf{Exégèse} = Rendre intelligible le sens \emph{profond} en mobilisant le contexte de l'œuvre, l'histoire des concepts et la logique interne du texte (niveau Bac/Probatoire). L'exégèse \textbf{exige} l'étape 4 (Fonction argumentative).
\end{definition}

\begin{attention}[Piège Q2 : Le "Hors-Texte" Déguisé]
Ne faites pas une dissertation sur le concept cité (ex: faire un cours sur "La Justice chez Rawls" parce que le mot "justice" est dans le passage). \textbf{Ancrez tout dans le passage.} Chaque phrase de votre explication doit pouvoir se rattacher à un mot ou une structure de la citation.
\end{attention}

\courssub{Question 3 : La Discussion Critique (Le Cœur de l'Épreuve)}

C'est la seule question où votre \textbf{propre pensée philosophique} est attendue, mais elle doit être \emph{disciplinée} par la méthode.

\begin{definition}[Discussion Philosophique ≠ Opinion Personnelle]
\textbf{Opinion} : « Je pense que la corruption est mal car c'est illégal. » (Subjectif, dogmatique, sans portée universelle).
\textbf{Discussion} : « La thèse de l'auteur (la corruption comme régulateur) tient si l'on adopte une perspective \emph{utilitariste} (efficacité court terme), mais elle s'effondre sous l'angle \emph{déontologique} (Kant : l'homme comme fin) ou \emph{républicain} (Rousseau : corruption = destruction du lien social). Au Cameroun, l'opération \emph{Epervier} montre que le "graissage" paralyse in fine l'investissement public. »
$\rightarrow$ La discussion \textbf{confronte} la thèse à d'autres arguments (auteurs, concepts, faits sociaux avérés) pour en tester la solidité.
\end{definition}

\begin{popfigure}
\begin{tikzpicture}[node distance=1.2cm, every node/.style={font=\small, align=center}]
\node[draw, rectangle, fill=popblue!20, thick] (q) {Question 3 : Discutez};
\node[draw, rectangle, fill=popgreen!20, below left of=q, xshift=-1.2cm, align=center] (p1){1. Concessions\\ (Validité interne,\\ contexte historique,\\ part de vérité)};
\node[draw, rectangle, fill=poporange!20, below of=q, align=center] (p2){2. Limites / Objections\\ (Contre-exemples CM,\\ Auteurs opposés,\\ Concepts contraires)};
\node[draw, rectangle, fill=poppink!20, below right of=q, xshift=1.2cm, align=center] (p3){3. Nuance / Synthèse\\ (Condition de validité,\\ Distinction conceptuelle,\\ Dépassement)};
\draw[->, thick, popdark] (q.south west) -- (p1.north east);
\draw[->, thick, popdark] (q.south) -- (p2.north);
\draw[->, thick, popdark] (q.south east) -- (p3.north west);
\end{tikzpicture}
\legende{Architecture type d'une discussion philosophique réussie (Q3) : la dialectique en trois temps.}
\end{popfigure}

\begin{methode}[Structure Type de la Q3 : La Dialectique "Oui / Mais / Donc"]
\begin{enumerate}
    \item \textbf{Concessions (Le "Oui") :} Reconnaissez la force du texte. Dans quel contexte / sous quelle condition / au nom de quel principe la thèse est-elle \emph{justifiée} ? (Ex: « L'analyse de Bayart sur le "ventre" est pertinente pour comprendre la \emph{stabilité} apparente des régimes néopatrimoniaux au Cameroun des années 80. »)
    \item \textbf{Limites / Objections (Le "Mais") :} C'est le corps principal. Attaquez la thèse sur 3 fronts :
        \begin{itemize}
            \item \textbf{Conceptuel :} Concepts contraires (Justice vs Efficacité, Droit vs Fait, Universalisme vs Particularisme).
            \item \textbf{Auteur/École :} Mobilisez un auteur du programme qui contredit (Ex: Kant contre Machiavel / Bayart ; Rawls contre l'utilitarisme ; Marx contre la "ruse de la raison" hégélienne).
            \item \textbf{Factuel / Camerounais :} Contre-exemple réel, précis, daté (Ex: Coût de la corruption sur le budget santé/éducation ; Affaire Alucam ; Dysfonctionnement des guichets uniques).
        \end{itemize}
    \item \textbf{Nuance / Synthèse (Le "Donc") :} Ne laissez pas le lecteur sur une contradiction. Proposez une distinction qui sauve l'essentiel : « La thèse est valide \emph{si} on se place au niveau de l'analyse sociologique descriptive (comment ça marche), mais inacceptable \emph{si} on se place au niveau de la norme politique prescriptive (comment ça doit marcher). La distinction \emph{Fait / Valeur} (Hume) est ici cruciale. »
\end{enumerate}
\end{methode}

\begin{attention}[Piège Majeur Q3 : La Dissertation "Hors-Texte"]
\begin{itemize}
    \item \textbf{Symptôme :} Vous faites un plan de dissertation standard (Thèse/Antithèse/Synthèse) sur le thème général (ex: "La corruption") sans \emph{jamais} citer l'auteur du texte, sans reprendre ses arguments précis, sans discuter \emph{sa} thèse spécifique.
    \item \textbf{Règle d'or :} Chaque paragraphe de la Q3 doit commencer ou finir par une référence au texte : « \emph{L'auteur affirme que...}, or... » ; « \emph{Cet argument (le graissage) néglige que...} » ; « \emph{Contrairement à ce que suggère le texte...} »
    \item \textbf{Sanction :} Note plafonnée à 6-7/20 même si la dissertation est belle, car la consigne « Discutez \textbf{la thèse de l'auteur} » n'est pas respectée.
\end{itemize}
\end{attention}

\courssub{Cas Pratique Intégré : « La Corruption comme "Graissage" Social »}

Nous appliquons la méthode à un texte type Bac Tech, inspiré des analyses de Jean-François Bayart (\emph{L'État en Afrique : la politique du ventre}) et des débats sur l'économie informelle.

\begin{exemplebox}[Texte Support (Extrait fictif type Bac)]
« On condamne trop vite la corruption au nom d'une morale abstraite, importée de l'Occident. Dans nos sociétés camerounaises, où l'État est souvent absent ou défaillant, la corruption agit comme un \textbf{graissage} nécessaire des rouages administratifs. Le "pourboire" au fonctionnaire sous-payé n'est pas un détournement, c'est un \textbf{complément de salaire} négocié qui assure un service public que l'État ne finance plus. Le "piston" pour l'embauche n'est pas du népotisme, c'est une \textbf{régulation communautaire} du chômage : on place "les siens" pour assurer la survie du lignage. Ainsi, loin de bloquer la machine, la corruption \textbf{huile le moteur social}. Elle est l'expression pragmatique d'une solidarité mécanique (Durkheim) adaptée à la pénurie. La combattre frontalement sans augmenter les salaires ni créer d'emplois, c'est gripper le moteur et condamner les plus faibles à l'asphyxie. »
\end{exemplebox}

\textbf{Question 1 : Dégagez la thèse de l'auteur et les arguments principaux.}

\begin{exoresolu}[Corrigé Type Q1]
\textbf{Question 1 :} Dégagez la thèse de l'auteur et les arguments principaux qu'il utilise pour la soutenir. \tcblower
\textbf{Thèse :} L'auteur soutient que la corruption, loin d'être un dysfonctionnement moral, remplit une \textbf{fonction sociale positive et nécessaire} (graissage, régulation, survie) dans le contexte spécifique de défaillance de l'État camerounais, et que sa répression pure sans réformes structurelles est contre-productive.

\textbf{Arguments principaux :}
\begin{enumerate}
    \item \textbf{Argument anthropologique / Économique (Le "Graissage") :} La corruption compense l'insuffisance des salaires publics (fonctionnaires sous-payés) et garantit la continuité du service public (efficacité pragmatique).
    \item \textbf{Argument Sociologique / Culturel (La "Régulation communautaire") :} Le népotisme/piston n'est pas du favoritisme illégitime mais une redistribution solidaire des ressources rares (emplois) au profit du groupe de parenté (lignage), assurant la cohésion sociale (référence à Durkheim / solidarité mécanique).
    \item \textbf{Argument Politique / Conséquentialiste (Le risque de "Grippe") :} La lutte anti-corruption purement répressive (opération \emph{Epervier} style) sans mesures d'accompagnement (augmentation salaires, création emplois) détruit le dernier filet de sécurité des populations vulnérables.
\end{enumerate}
\textbf{Articulation :} Le texte progresse du niveau micro-individuel (fonctionnaire/salaire) au niveau méso-social (lignage/chômage) pour aboutir au niveau macro-politique (État/réforme), construisant une défense systémique de la corruption comme "mal nécessaire".
\end{exoresolu}

\textbf{Question 2 : Expliquez le passage suivant : « La corruption agit comme un graissage nécessaire des rouages administratifs... elle huile le moteur social. »}

\begin{exoresolu}[Corrigé Type Q2 - Application méthode S.D.R.F.]
\textbf{Question 2 :} Expliquez le passage suivant : « La corruption agit comme un graissage nécessaire des rouages administratifs... elle huile le moteur social. » \tcblower
\textbf{1. Situer :} Passage d'ouverture (thèse directrice + métaphore centrale). Donne le ton : provocateur, pragmatique, anti-moraliste.
\textbf{2. Définir \emph{in situ} :}
\begin{itemize}
    \item \textbf{Graissage / Huile :} Métaphore mécanique. Réduit les "frottements" (bureaucratie lente, formalismes excessifs, absence de ressources). Rend le système \emph{fluide} / \emph{efficient} à court terme.
    \item \textbf{Rouages administratifs / Moteur social :} L'appareil d'État et la société dans son ensemble vus comme une machine productive.
    \item \textbf{Nécessaire :} Au sens fort : \emph{condition de possibilité} du fonctionnement actuel. Sans lui, arrêt total (grippe).
\end{itemize}
\textbf{3. Reformuler :} Dans un État camerounais affaibli (sous-administration, sous-salaires), les pratiques corruptives (petits pots-de-vin, arrangements) ne sont pas des grains de sable mais l'huile qui permet à la machine administrative de tourner. Elles transforment une contrainte rigide (la règle officielle inapplicable) en transaction fluide (le service rendu contre rémunération réelle).
\textbf{4. Fonction argumentative :} Ce passage \textbf{pose la thèse centrale} (utilitarisme contextuel) et \textbf{installe la métaphore filée} (mécanique/thermique) qui structure tout le texte : corruption = maintenance / répression = sabotage. Il anticipe l'objection morale ("condamne trop vite") par l'argument de l'efficacité réelle ("ne bloque pas, huille").
\end{exoresolu}

\textbf{Question 3 : Le texte vous semble-t-il convaincant ? Discutez la thèse de l'auteur.}

\begin{exoresolu}[Corrigé Type Q3 - Structure Dialectique]
\textbf{Question 3 :} Le texte vous semble-t-il convaincant ? Discutez la thèse de l'auteur. \tcblower
\textbf{Introduction (Accroche + Rappel thèse + Problématique + Plan) :}
L'auteur nous invite à un renversement spectaculaire : la corruption, habituellement vue comme un \emph{frein} au développement (Banque Mondiale, Transparency International), serait ici un \emph{moteur} de survie sociale. \textbf{Problématique :} L'efficacité pragmatique du "graissage" à court terme peut-elle justifier moralement et politiquement l'instauration d'un système de prédation qui détruit l'État de droit à long terme ?

\textbf{Développement :}

\textbf{I. Concessions : La pertinence descriptive du "Graissage" (Validité contextuelle)}
L'auteur a raison sur le plan \textbf{sociologique descriptif} (Bayart, \emph{La politique du ventre}). Au Cameroun, le "petit corruption" (raket aux barrages, pourboire aux urgences, frais de dossier inventés) est effectivement le \emph{mode de fonctionnement réel} de l'administration depuis les crises structurelles des années 80-90 (dévaluation, gel embauche). Il assure une \textbf{redistribution informelle} de la rente étatique vers les agents de base et leurs réseaux. La métaphore du "graissage" rend compte de la \emph{fluidité transactionnelle} : sans "motivation", le fonctionnaire n'a aucun incitatif à traiter le dossier. La référence à Durkheim (solidarité mécanique) éclaire la logique du "piston" : dans une société à forte solidarité segmentaire (ethnie, clan, loge), placer un parent est un devoir moral, pas un délit. \textbf{Le texte a raison : réprimer sans comprendre ni remplacer le mécanisme, c'est produire de l'exclusion sociale immédiate.}

\textbf{II. Limites et Objections : Le prix caché du "Graissage" (Invalidité normative et structurelle)}
\textit{1. Objection Éthique / Déontologique (Kant / Droit) :} L'auteur confond \emph{Fait} (ça marche) et \emph{Droit} (c'est juste). Kant (\emph{Fondements de la métaphysique des mœurs}) : « Agis de telle sorte que tu traites l'humanité... toujours comme une fin, jamais simplement comme un moyen. » Le "graissage" instrumentalise le service public (droit du citoyen) en marchandise privée. Le fonctionnaire devient un \emph{racketier} légitimé. L'universalité de la loi s'effondre : le service n'est plus un droit égal pour tous, mais une faveur pour qui paie.
\textit{2. Objection Politique / Républicaine (Rousseau / Contrat social) :} La corruption détruit la \textbf{Volonté Générale}. Le "piston" ethnique/clanique (argument 2 du texte) nie le mérite républicain et fracture la nation en "ventres" concurrents (Bayart). Au Cameroun, la gestion ethnique des directions générales et des marchés publics (ex: secteur forestier, commande publique) a conduit à une \textbf{prédation systémique} où l'État n'est plus l'arbitre commun mais le butin des clans. Le "graissage" ne huile pas le moteur national, il \emph{corrode le bloc moteur} : l'impôt ne revient plus en service public, mais en rente privée.
\textit{3. Objection Économique Structurelle (Nouvelle Économie Institutionnelle / Acemoglu) :} L'argument de l'efficacité court terme ignore les \textbf{coûts de transaction systémiques}. L'incertitude juridique (le "graissage" est arbitraire, non contractuel) décourage l'investissement productif long terme (usines, R\&D) au profit de la rente rapide (import-export, marchés publics surfacturés). Le Cameroun perd des milliards en Investissements Directs Étrangers (IDE) à cause de l'imprévisibilité corruptive (Rapports Doing Business, INDICE de perception de la corruption). Le "moteur" tourne au ralenti, consomme tout son huile (recettes publiques) et ne produit plus de puissance (infrastructures, santé, éducation).

\textbf{III. Nuance / Synthèse : La distinction Fait / Valeur et l'urgence de la réforme structurelle}
La thèse de l'auteur est \textbf{valide comme diagnostic sociologique} (le "graissage" \emph{existe} et \emph{fonctionne} comme mode de survie), mais \textbf{inacceptable comme prescription politique} (le "graissage" ne doit pas \emph{devenir} la norme).
La solution n'est pas le choix binaire : « Corruption ou Asphyxie » (faux dilemme du texte).
La \textbf{synthèse} exige la \textbf{distinction conceptuelle} (Weber) entre :
\begin{itemize}
    \item \textbf{Légitimité rationnelle-légale :} L'État doit payer des salaires décents (supprimer le \emph{besoin} de graissage), numériser les procédures (supprimer l'\emph{occasion} du graissage : ex: portail unique douanier, concours anonymisés), sanctionner la grande corruption (récupérer l'huile volée).
    \item \textbf{Légitimité traditionnelle/charismatique :} La solidarité lignagère doit être canalisée vers la société civile (associations, mutuelles, coopératives) et non vers la prédation de la chose publique.
\end{itemize}
Le "graissage" est un \emph{symptôme} de la faillite de l'État-providence, pas son \emph{remède}. Le traiter comme remède, c'est condamner le Cameroun à l'économie de rente éternelle.

\textbf{Conclusion :}
Le texte a le mérite de briser l'hypocrisie moralisatrice pour décrire la réalité vécue des Camerounais. Mais en naturalisant la corruption ("nécessaire", "mécanique"), il désarme la citoyenneté. La philosophie politique (de Platon à Rawls) nous enseigne que la \textbf{Justice} n'est pas l'efficacité du plus fort, mais l'équité des règles du jeu. Le défi camerounais n'est pas de mieux "huiler" la machine corruptive, mais de \textbf{changer de moteur} : passer de la prédation personnalisée à l'institutionnalisation transparente.
\end{exoresolu}

\courssub{Gestion du Temps \& Méta-cognition : Votre Tableau de Bord Jour J}

\begin{methode}[Stratégie "Chrono-Bac" (4h d'épreuve)]
\begin{enumerate}
    \item \textbf{0h00 - 0h45 : Lecture active \& Brouillon Q1/Q2.} (Lecture 1: globale. Lecture 2: crayon en main, repérage thèse/args/concepts. Lecture 3: focus passage Q2). \textbf{Rédigez Q1 et Q2 sur brouillon propre.}
    \item \textbf{0h45 - 1h15 : Copie Propre Q1 \& Q2.} (Soignez la présentation, numérotez les arguments, structurez l'exégèse S.D.R.F.).
    \item \textbf{1h15 - 1h30 : Pause cérébrale / Relecture texte + Préparation Q3.} (Identifiez les 2-3 concepts clés à mobiliser, les 2 auteurs opposés, les 2 exemples Cameroun).
    \item \textbf{1h30 - 3h15 : Rédaction Q3 (Le temps fort).} \textbf{1h45 min minimum.} Appliquez la structure Concessions / Limites / Nuance. Citez le texte \emph{en permanence}.
    \item \textbf{3h15 - 3h30 : Relecture globale + Conclusion Q3 (Ouverture soignée si temps, sinon phrase de chute sèche).}
    \item \textbf{3h30 - 4h00 : Marge de sécurité / Relecture orthographe / Vérification consignes.}
\end{enumerate}
\end{methode}

\begin{aretenir}[Check-list "Zéro Regret" avant de rendre la copie]
\begin{itemize}
    \item [ ] \textbf{Q1 :} Thèse en \textbf{une phrase}. Arguments \textbf{numérotés 1, 2, 3}. Articulation explicitée (\emph{« Enchaînement... »}). \textbf{Zéro « Je »}.
    \item [ ] \textbf{Q2 :} Concepts définis \emph{in situ}. Passage reformulé. \textbf{Fonction argumentative} claire (prouve l'arg X, transition vers Y).
    \item [ ] \textbf{Q3 :} \textbf{Plan visible} (I, II, III ou A, B, C). \textbf{Ancre textuelle} dans chaque grand paragraphe (\emph{« L'auteur écrit... »}). Auteurs du programme cités \textbf{nommément} (Kant, Rousseau, Rawls, Marx, Bayart, Durkheim, Weber...). Exemples \textbf{Cameroun précis} (Opération Épervier, Portail unique, Code du travail, Prestation serment barreau...). \textbf{Nuance finale} (Distinction conceptuelle).
    \item [ ] \textbf{Forme :} Paragraphes aérés. Connecteurs logiques (\emph{Or, Toutefois, En effet, Par conséquent, En définitive}). Pas de SMS, pas de familier.
\end{itemize}
\end{aretenir}

\begin{savaistu}[Le "Secret" des Correcteurs Bac Tech]
Les correcteurs du Bac Tech (Philosophie) sont souvent des agrégés qui enseignent \emph{aussi} en Terminale Générale. Ils attendent de vous la \textbf{même rigueur conceptuelle} que les généraux, mais appliquée aux \textbf{réalités techniques et sociales} de vos spécialités (Génie Civil, Électrotechnique, Commerce, Secrétariat, Santé...).
\begin{itemize}
    \item Un élève en \textbf{Génie Civil} qui discute la corruption sur les marchés publics en citant \emph{le Cahier des Clauses Administratives Générales (CCAG)} ou la \emph{loi sur les marchés publics} marque des points majeurs.
    \item Un élève en \textbf{Santé} qui discute l'éthique médicale (secret professionnel, consentement) face à la corruption hospitalière (détournement médicaments, fausses factures) ancre sa philosophie dans \emph{sa} technique.
    \item Un élève en \textbf{Commerce/Gestion} qui mobilise la \emph{théorie de l'agence} (Jensen \& Meckling) ou la \emph{RSE (Responsabilité Sociétale des Entreprises)} pour discuter le "graissage" montre qu'il transfère ses acquis techniques en philosophie.
\end{itemize}
\textbf{Conseil :} Relisez vos cours de \textbf{Droit du Travail, Droit des Affaires, Éthique Professionnelle, Économie de l'Entreprise}. Ce sont des réservoirs d'exemples et de concepts \emph{autorisés} et \emph{valorisés} en Q3.
\end{savaistu}

\begin{exercice}[Entraînement Chronométré (Simu Bac Blanc - 1h30)]
\textbf{Texte :} Extrait de \emph{La Condition ouvrière} de Simone Weil (1951) sur la distinction entre \emph{besoins vitaux} (nourrir le corps) et \emph{besoins de l'âme} (liberté, vérité, propriété, sécurité, responsabilité) et l'aliénation du travail à la chaîne.
\textbf{Questions :}
\begin{enumerate}
    \item Dégagez la thèse et les arguments principaux. (Visé : 20 min)
    \item Expliquez : « L'obéissance forcée est la forme la plus courante de l'oppression. Mais elle n'est pas la plus cruelle. La plus cruelle, c'est celle où l'opprimé consent à son oppression. » (Visé : 25 min)
    \item « Le travail peut-il être source de libération ? » Discutez à partir du texte. (Visé : 45 min)
\end{enumerate}
\textit{Consigne :} Faites cet exercice en conditions réelles (chronomètre, pas de notes, copie blanche). Corrigez-vous avec la grille \textbf{Check-list "Zéro Regret"} ci-dessus.
\end{exercice}

\begin{corrige}[Exercice n°1 - Pistes de correction rapides]
\textbf{Q1 Thèse :} L'ouvrier à la chaîne est aliéné non par la fatigue physique mais par la \emph{privation des besoins de l'âme} (liberté, responsabilité, vérité) que l'organisation taylorienne du travail impose. \textbf{Args :} 1. Description phénoménologique de l'oppression invisible (consentement forcé). 2. Analyse des besoins de l'âme bafoués (propriété usage vs propriété légale, responsabilité réduite à zéro). 3. Critique du syndicalisme purement revendicatif (salaire) qui néglige la condition spirituelle.
\textbf{Q2 Passage :} Distinction \emph{obéissance forcée} (contrainte externe, visible, révoltante) vs \emph{consentement à l'oppression} (intériorisation, aliénation mentale, « servitude volontaire » à la La Boétie). Fonction : Montre que l'aliénation moderne est \emph{psychique} plus que physique, donc plus radicale.
\textbf{Q3 Discussion :} Concession : Weil a raison sur la dimension anthropologique (travail = condition de l'humanisation \emph{si} il porte sens/responsabilité). Limites : 1. Marx : l'aliénation est structurelle (propriété privée des moyens de production), pas seulement organisationnelle ; le "consentement" est idéologique (fausse conscience). 2. Arendt / Hannah Arendt : distinction \emph{Labor} (survie) vs \emph{Work} (fabrication) vs \emph{Action} (politique) ; le salaire peut financer la vie politique (libération \emph{hors} travail). 3. Contexte Cameroun : Informel / Chômage : le "travail décent" (OIT) est un luxe ; la priorité vitale (besoin du corps) précède le besoin de l'âme (Maslow / Sen). Nuance : La libération par le travail exige \textbf{double condition} : transformation des rapports de production (structure) ET réappropriation du sens par le sujet (subjectivité).
\end{corrige}

\coursec{Simulation Bac Blanc \& M\'eta-cognition (Auto-\'evaluation)}

Apr\`es avoir ma\^itris\'e, dans la section pr\'ec\'edente, les \cle{questions types} et d\'ejou\'e les \cle{pi\`eges sp\'ecifiques} au Bac Tech (reformulation, explicitation, discussion), il est temps de passer de l'entra\^inement fragment\'e \`a la \textbf{simulation int\'egrale}. C'est le moment de v\'erit\'e : mettre en musique l'ensemble de la cha\^ine op\'eratoire --- lecture, analyse, probl\'ematisation, planification, r\'edaction, relecture --- dans les conditions r\'eelles de l'examen (1h30). Cette s\'eance n'est pas un simple \og{}devoir surveill\'e\fg{} ; c'est un \textbf{dispositif m\'eta-cognitif} : vous allez devenir votre propre correcteur pour identifier, \emph{avant} le jour J, les derniers points de friction qui co\^utent des points.

\courssub{L'\'Epreuve Int\'egrale Chronom\'etr\'ee : Conditions R\'eelles}

\emph{Intuition :} On ne pr\'epare pas un marathon en courant 100 m\`etres. La gestion du temps, la fatigue cognitive, la tentation de l'improvisation ne se g\`erent qu'en situation.

\begin{methode}[Protocole de la Simulation Bac Blanc (1h30 exactes)]
\begin{enumerate}
    \item \textbf{Installation \& Consignes (0h00 - 0h05) :} Pas de t\'el\'ephone, pas de brouillon pr\'e-rempli. Uniquement : sujet, copies doubles, stylos (bleu/noir), correcteur (si autoris\'e), montre.
    \item \textbf{Phase 1 : Lecture Active \& Analyse (0h05 - 0h20 \textbf{[15 min]}) :}
    \begin{itemize}
        \item \textbf{Lecture 1 (3 min) :} \og{}Qu'est-ce que dit ce texte globalement ?\fg{} (Th\`ese brute).
        \item \textbf{Lecture 2 (7 min) :} Rep\'erage structurel (paragraphes, arguments, connecteurs), surlignage des \cle{concepts-cl\'es}, des \cle{citations} potentielles, des \cle{exemples}.
        \item \textbf{Lecture 3 (5 min) :} Formulation de la \cle{probl\'ematique} et esquisse du \cle{plan d\'etaill\'e} (2 ou 3 parties, 2 sous-parties chacune, avec arguments textuels et hors-texte pr\'ecis).
    \end{itemize}
    \item \textbf{Phase 2 : R\'edaction de l'Introduction au Propre (0h20 - 0h35 \textbf{[15 min]}) :} \'Ecrire l'intro \emph{au propre} directement (gain de temps, engagement). Structure canonique : Accroche $\rightarrow$ Pr\'esentation texte/auteur $\rightarrow$ Th\`ese $\rightarrow$ Probl\'ematique $\rightarrow$ Annonce de plan.
    \item \textbf{Phase 3 : D\'eveloppement (0h35 - 1h15 \textbf{[40 min]}) :} 20 min par grande partie. \textbf{R\`egle d'or :} 1 paragraphe = 1 id\'ee = 1 argument textuel (citation/paraphrase) + 1 analyse conceptuelle + 1 exemple/contre-exemple/auteur + 1 mini-transition. Sauter une ligne entre paragraphes. Indenter.
    \item \textbf{Phase 4 : Conclusion (1h15 - 1h25 \textbf{[10 min]}) :} Bilan synth\'etique (r\'eponse \`a la probl\'ematique) + Ouverture pertinente (autre auteur, actualit\'e, autre notion du programme). Pas de nouvelle id\'ee.
    \item \textbf{Phase 5 : Relecture Active \og{}Grille en main\fg{} (1h25 - 1h30 \textbf{[5 min]}) :} V\'erifier : citations exactes ? concepts d\'efinis ? plan visible ? fautes d'accord ? lisibilit\'e ?
\end{enumerate}
\end{methode}

\begin{attention}[Gestion du Stress \& Le Pi\`ege du \og{}Vide\fg{}]
\textbf{Erreur fr\'equente :} Paniquer face \`a un texte difficile (ex: Mbembe ou Foucault) et se mettre \`a \'ecrire n'importe quoi d\`es la premi\`ere minute.
\textbf{Technique \og{}Respiration 4-7-8\fg{} + Lecture \og{}Sans Stylo\fg{} :} Avant d'ouvrir le sujet, fermez les yeux. Inspirez 4 sec, bloquez 7 sec, expirez 8 sec (x3). Puis lisez le texte \textbf{sans stylo en main} pendant 5 min. Le stylo fige la pens\'ee ; la main libre lib\`ere la compr\'ehension globale. Le \textbf{brouillon est obligatoire} (20\% du temps = 18 min) : c'est votre \og{}bo\^ite noire\fg{}, le seul endroit o\`u vous avez le droit d'\^etre confus.
\end{attention}

\paragraph{Sujet Type Bac Tech (In\'edit - Simulation) :}
\textbf{Texte :} Extrait d'\emph{Achille Mbembe, \emph{Critique de la raison n\`egre}, Paris, La D\'ecouverte, 2013 (chap. \og{}La chose et le f\'etiche\fg{})} \emph{OU} \textbf{Michel Foucault, \emph{Surveiller et punir}, Paris, Gallimard, 1975 (Partie 3, chap. 3 \og{}Le carcan\fg{}, sur l'institution scolaire comme appareil disciplinaire)}.
\textit{Contexte camerounais implicite :} L'\'ecole technique comme lieu de \og{}fabrication\fg{} du sujet productif, le corps de l'\'el\`eve (tenue, posture, horaire), l'examen comme rituel de normalisation.

\textbf{Questions (Type Bac Tech - 20 pts) :}
\begin{enumerate}
    \item \textbf{Expliquez la th\`ese de l'auteur et les concepts principaux qu'il mobilise pour l'\'etayer.} (6 pts)
    \item \textbf{Montre comment l'auteur structure son argumentation pour d\'emontrer que l'institution (scolaire/carc\'erale) produit un sujet \og{}docile\fg{}.} (6 pts)
    \item \textbf{La discipline est-elle la condition n\'ecessaire de la libert\'e, ou son contraire ? Discutez en vous appuyant sur le texte et vos connaissances.} (8 pts)
\end{enumerate}

\courssub{Correction par les Pairs : La Grille Officielle MINESEC Simplifi\'ee}

\emph{Intuition :} Corriger la copie d'un autre, c'est changer de lunettes. On voit les d\'efauts de structure, les flous de langage, les citations approximatives qu'on ne voit pas dans sa propre prose. C'est l'apprentissage du \cle{regard du correcteur}.

\begin{definition}[Grille d'Auto/H\'et\'ero-\'evaluation Bac Philo Tech (Sur 20)]
La grille officielle MINESEC p\`ese 4 grands items. Pour l'entra\^inement, on la d\'ecline en sous-crit\`eres observables (Oui/Non/Partiel).
\begin{center}
\begin{tabularx}{\linewidth}{|l|X|c|}\hline
\textbf{Item Principal (Pts)} & \textbf{Sous-crit\`eres observables (Check-list)} & \textbf{Pts} \\ \hline
\textbf{1. Introduction \& Probl\'ematique (4)} & Accroche pertinente / Pr\'esentation Auteur/\OE uvre/Th\`eme / Th\`ese formul\'ee / \textbf{Probl\'ematique explicite (question philosophique)} / Plan annonc\'e clair (2-3 parties) & 4 \\ \hline
\textbf{2. D\'eveloppement : Structure \& Contenu (10)} & Plan respect\'e / Paragraphes distincts (saut de ligne, indentation) / \textbf{Citations exactes} (guillemets, ref) / \textbf{Concepts d\'efinis} / Analyse textuelle (pas paraphrase) / Apports personnels (auteurs, ex. Cameroun) / \textbf{Transitions logiques} / Discussion argument\'ee (th\`eses/antith\`eses/synth\`eses) & 10 \\ \hline
\textbf{3. Conclusion (4)} & Bilan synth\'etique (r\'eponse \`a la prob.) / Ouverture pertinente (pas \og{}pour conclure je dirai que...\fg{}) / Style soign\'e & 4 \\ \hline
\textbf{4. Expression \& Langue (2)} & Syntaxe correcte / Vocabulaire philosophique pr\'ecis / Orthographe / Lisibilit\'e (\'ecriture, ratures) & 2 \\ \hline
\end{tabularx}
\end{center}
\end{definition}

\begin{popfigure}
\begin{tikzpicture}[node distance=1.2cm, every node/.style={font=\small, align=center, inner sep=4pt}]
\node[draw, rectangle, rounded corners, fill=popblueL, thick] (grille) {\textbf{GRILLE AUTO-\'EVALUATION (Sur 20)}};
\node[draw, rectangle, rounded corners, fill=popgreenL, below left of=grille, xshift=-1.8cm, align=center] (c1){\textbf{Intro (4)} : \\ Accroche / Pr\'esent / Th\`ese / Prob / Plan};
\node[draw, rectangle, rounded corners, fill=popgoldL, below of=grille, align=center] (c2){\textbf{D\'eveloppement (10)} : \\ Structure / Textuel / Concepts / Logique / Transition};
\node[draw, rectangle, rounded corners, fill=poporangeL, below right of=grille, xshift=1.8cm, align=center] (c3){\textbf{Conclusion (4)} : \\ Bilan / R\'eponse Prob / Ouverture};
\node[draw, rectangle, rounded corners, fill=poppinkL, below of=c2, yshift=-0.6cm, align=center] (c4){\textbf{Langue (2)} : \\ Syntaxe / Vocab / Orthographe};
\draw[popblue, thick, ->] (grille.south west) -- (c1.north east);
\draw[popblue, thick, ->] (grille.south) -- (c2.north);
\draw[popblue, thick, ->] (grille.south east) -- (c3.north west);
\draw[popblue, thick, ->] (c2.south) -- (c4.north);
\end{tikzpicture}
\legende{Sch\'ema de la grille d'auto-\'evaluation simplifi\'ee (4 items). Chaque item se d\'ecompose en sous-crit\`eres binaires (Valid\'e / Non valid\'e) pour une notation rapide et objective.}
\end{popfigure}

\begin{methode}[Proc\'edure de Correction par les Pairs (30 min)]
\begin{enumerate}
    \item \textbf{Anonymat :} \'Echanger les copies (num\'eros au lieu des noms).
    \item \textbf{Lecture \og{}Correcteur\fg{} (10 min) :} Lire \emph{une fois} sans stylo pour l'impression globale. Relire \emph{avec la grille} : cocher chaque sous-crit\`ere. Mettre un \textbf{+} (acquis), \textbf{/} (partiel), \textbf{-} (manquant) dans la marge.
    \item \textbf{Commentaire \og{}Miroir\fg{} (10 min) :} Sur une fiche annexe, r\'ediger : \og{}Ce qui est fort (2 points pr\'ecis)\fg{} / \og{}Ce qui fait perdre des points (3 points pr\'ecis avec n$^\circ$ de crit\`ere)\fg{} / \og{}Conseil n$^\circ$1 pour gagner 2 pts au Bac\fg{}.
    \item \textbf{Restitution (10 min) :} Remettre la copie + fiche. L'\'el\`eve lit la fiche, note ses 3 faiblesses majeures dans son \og{}Plan d'Action Personnel\fg{} (voir 5.4).
\end{enumerate}
\end{methode}

\courssub{Analyse de Copies Types (Anonymis\'ees) : Le Diagnostic Diff\'erentiel}

\emph{Intuition :} Comprendre \emph{pourquoi} une note est ce qu'elle est, c'est internaliser les standards. Nous analysons trois profils arch\'etypaux.

\begin{exemplebox}[Copie A : 16/20 --- \og{}La Copie Mod\`ele Structurelle\fg{}]
\textbf{Profil :} \'El\`eve s\'erieux, bon niveau langue, a int\'egr\'e la m\'ethode.
\begin{itemize}
    \item \textbf{Intro (4/4) :} Accroche sur la \og{}discipline\fg{} au Cameroun (ex: lever des couleurs, tenue kaki). Th\`ese de Foucault/Mbembe pr\'ecise : \og{}L'institution fabrique la docilit\'e par l'espace, le temps, l'examen\fg{}. Probl\'ematique : \textit{\og{}Jusqu'o\`u l'assujettissement disciplinaire est-il le prix de l'autonomie ?\fg{}}. Plan : I. L'anatomie politique du corps \'el\`eve (Espace/Temps). II. L'examen comme ritualisation de la norme. III. La marge de manœuvre : r\'esistance et subjectivation.
    \item \textbf{D\'ev (9/10) :} Paragraphes calibr\'es. Citation exacte : \og{}L'examen entoure de toute sa puissance documentaire ceux qu'il soumet\fg{} (Foucault). Concept \cle{normation} d\'efini. Exemple camerounais : le Probatoire comme \og{}jugement normalisant\fg{}. Transition I$\rightarrow$II : \textit{\og{}Mais cette emprise spatiale/temporelle trouve son point d'orgue dans l'\'epreuve certificative...\fg{}}. Discussion : Apport de Kant (\emph{Qu'est-ce que les Lumi\`eres ?} : sortie de la minorit\'e) vs Foucault.
    \item \textbf{Conclu (3/4) :} Bilan solide. Ouverture sur \cle{biopolitique} (Foucault) / \cle{n\'ecropolitique} (Mbembe) : \og{}L'\'ecole technique forme-t-elle des sujets libres ou des ressources humaines ?\fg{}. Perte 1 pt : ouverture un peu rapide.
    \item \textbf{Langue (2/2) :} Impeccable.
\end{itemize}
\textbf{Le\c{c}on :} La \textbf{visibilit\'e du plan} (sauts de lignes, connecteurs \og{}En premier lieu / Par ailleurs / En d\'efinitive\fg{}) et les \textbf{citations \og{}chirurgicales\fg{}} (courtes, int\'egr\'ees, analys\'ees) font la diff\'erence.
\end{exemplebox}

\begin{exemplebox}[Copie B : 9/20 --- \og{}Le Pi\`ege de la Paraphrase Narrative\fg{}]
\textbf{Profil :} \'El\`eve qui a \og{}lu le texte\fg{} mais pas \og{}fait de philo\fg{}.
\begin{itemize}
    \item \textbf{Intro (1/4) :} Pas d'accroche. \og{}Foucault dit que l'\'ecole c'est comme la prison.\fg{} Pas de th\`ese formul\'ee. Probl\'ematique absente (remplac\'ee par : \og{}Nous allons voir comment l'\'ecole discipline\fg{}). Plan : \og{}D'abord l'espace, ensuite le temps, enfin l'examen\fg{} (descriptif, pas argumentatif).
    \item \textbf{D\'ev (5/10) :} \textbf{Paraphrase continue} : \og{}Foucault explique que les \'el\`eves sont assis en rang. Il dit que l'horaire est strict. Il raconte l'examen.\fg{} Aucune citation (ou citation de 3 lignes mal copi\'ee). Concepts absents (\cle{docilit\'e}, \cle{normalisation}, \cle{panoptisme} non nomm\'es). Exemples : \og{}Chez nous au lyc\'ee, on porte l'uniforme.\fg{} (Sociologique, pas philosophique). Pas de transitions. Blocs de texte sans paragraphes.
    \item \textbf{Conclu (1/4) :} \og{}En conclusion, l'\'ecole discipline mais c'est pour notre bien.\fg{} Pas de bilan, pas d'ouverture.
    \item \textbf{Langue (2/2) :} Correcte.
\end{itemize}
\textbf{Diagnostic :} \textbf{Confusion Explicitation / Discussion.} L'\'el\`eve raconte le texte (Q1) au lieu de l'analyser (Q2/Q3). \textbf{Absence de probl\'ematique = plan descriptif = note < 10.}
\end{exemplebox}

\begin{exemplebox}[Copie C : 11/20 --- \og{}Le Bon Fond G\^ach\'e par la Forme\fg{}]
\textbf{Profil :} \'El\`eve cultiv\'e, de bonnes id\'ees, mais d\'esorganis\'e et press\'e.
\begin{itemize}
    \item \textbf{Intro (1/4) :} \textbf{Intro r\'edig\'ee... \`a la fin de la copie} (ou absente au d\'ebut). Pas de probl\'ematique claire. Plan flou.
    \item \textbf{D\'ev (8/10) :} \textbf{Excellent fond !} Analyse fine du \cle{panoptisme} appliqu\'e aux ateliers techniques. Citation de Mbembe sur le \cle{f\'etiche} marchandise. Discussion Kant/Foucault pertinente. \textbf{MAIS :} Pas de paragraphes visibles (\'ecriture continue), pas de transitions explicites, plan invisible pour le correcteur press\'e. Q3 (Discussion) : \textbf{Hors-sujet partiel} : glisse vers \og{}L'\'education au Cameroun manque de moyens\fg{} (sociologie/politique) au lieu de \og{}Discipline vs Libert\'e\fg{} (philosophie).
    \item \textbf{Conclu (1/4) :} Bacl\'ee en 2 lignes. \og{}Pour finir, la discipline est n\'ecessaire.\fg{}
    \item \textbf{Langue (1/2) :} Quelques fautes d'inattention.
\end{itemize}
\textbf{Diagnostic :} \textbf{Gestion du temps d\'efaillante} (trop de temps sur le d\'eveloppement, zapp\'e Intro/Conclu). \textbf{Invisibilit\'e de la structure} = correcteur ne voit pas le plan = note plafonn\'ee \`a 11-12 malgr\'e le fond. \textbf{Hors-sujet Q3} = perte de 3-4 pts sur 8.
\end{exemplebox}

\begin{exoresolu}[Analyse Comparative : O\`u sont les points ?]
\textbf{Question :} Pourquoi Copie A (16) vs Copie C (11) alors que C a un meilleur fond philosophique ?
\textbf{Solution :}
\begin{align*}
\text{Copie A} &: \text{Intro (4) + Dev (9) + Conclu (3) + Langue (2) = 18} \rightarrow \text{Ajustement s\'ev\'erit\'e ouverture } \rightarrow \mathbf{16} \\
\text{Copie C} &: \text{Intro (1) + Dev (8) + Conclu (1) + Langue (1) = 11} \\
\text{\'Ecart (5 pts)} &= \underbrace{(4-1)}_{\text{Intro}} + \underbrace{(3-1)}_{\text{Conclu}} + \underbrace{(2-1)}_{\text{Langue/Pr\'esentation}} + \underbrace{\text{P\'enalit\'e Hors-sujet Q3}}_{\text{-2 pts implicites}} \\
&= 3 + 2 + 1 + 2 = 8 \text{ pts de \og{}forme/m\'ethode\fg{} perdus par C.}
\end{align*}
\textbf{Enseignement :} Au Bac Tech, \textbf{la m\'ethode EST le fond}. Une d\'emonstration philosophique \emph{exige} une architecture visible. \og{}Je sais mais je n'ai pas le temps de faire l'intro\fg{} = \og{}Je ne sais pas d\'emontrer\fg{}.
\end{exoresolu}

\courssub{\'Elaboration du \og{}Plan d'Action Personnel\fg{} (PAP) : La M\'eta-cognition Op\'eratoire}

\emph{Intuition :} \og{}Je vais r\'eviser\fg{} ne veut rien dire. \og{}Je m'entra\^ine 30 min chaque mardi sur la probl\'ematisation\fg{} est un plan d'action. La m\'eta-cognition, c'est transformer le diagnostic (faiblesses) en prescription (exercices dat\'es).

\begin{methode}[Construction de votre PAP (A faire imm\'ediatement apr\`es la correction)]
\begin{enumerate}
    \item \textbf{Identifiez vos 3 Faiblesses Majeures (Top 3) :} Relisez votre fiche \og{}Miroir\fg{} + votre grille. Choisissez les 3 items qui vous font perdre le plus de points \emph{r\'eguli\`erement}.
    \item \textbf{Nommez-les avec le vocabulaire technique :} Pas \og{}je suis nul en intro\fg{} mais \textbf{\og{}Probl\'ematique floue / Plan non annonc\'e\fg{}}. Pas \og{}j'\'ecris mal\fg{} mais \textbf{\og{}Connecteurs logiques absents / Paragraphes non visibles\fg{}}. Pas \og{}je g\`ere mal le temps\fg{} mais \textbf{\og{}D\'epassement temps Dev / Zappage Intro/Conclu/Relecture\fg{}}.
    \item \textbf{Prescrivez 1 Exercice Cibl\'e / Semaine par Faiblesse (Jusqu'au Bac) :}
\end{enumerate}
\end{methode}

\begin{center}
\begin{tabularx}{\linewidth}{|l|X|X|}\hline
\textbf{Faiblesse Majeure (Item Grille)} & \textbf{Exercice Hebdomadaire Cibl\'e (30-45 min)} & \textbf{Indicateur de R\'eussite (KPI)} \\ \hline
\textbf{1. Probl\'ematique / Plan} & \textbf{Lundi :} \og{}1 Sujet = 3 Probl\'ematiques\fg{}. Prendre 1 sujet bac (annales). \'Ecrire 3 probs diff\'erentes. Choisir la meilleure. Faire le plan d\'etaill\'e (args + citations) en 15 min. & 3 probs distinctes / Plan complet en $<$ 15 min. \\ \hline
\textbf{2. Citation / Analyse Textuelle} & \textbf{Mercredi :} \og{}Chasse au Tr\'esor\fg{}. Relire 1 texte programme. Extraire 5 citations \og{}utilisables\fg{} (courtes, conceptuelles). \'Ecrire la phrase d'analyse pour chacune : \textit{\og{}Cette citation illustre le concept de X en montrant que...\fg{}}. & 5 citations exactes + 5 analyses types pr\^etes. \\ \hline
\textbf{3. Transitions / Structure Visible} & \textbf{Vendredi :} \og{}R\'edaction Squelette\fg{}. \'Ecrire \textbf{seulement} l'intro + les 4-6 phrases d'amorce de paragraphes (avec connecteurs) + la conclu. Pas le corps. V\'erifier : Plan visible ? Logique ? & Squelette complet en 20 min. Correcteur valide la logique. \\ \hline
\textbf{4. Gestion Temps / Relecture} & \textbf{Samedi (Simu) :} 1 Sujet complet \textbf{chronom\'etr\'e strict 1h30}. Obligation : Stylo pos\'e \`a 1h25 pour relecture grille. & Copie rendue \`a 1h30 pile. Grille auto-remplie $>$ 14/20. \\ \hline
\textbf{5. Hors-sujet / Discussion} & \textbf{Dimanche :} \og{}Le Duel\fg{}. Prendre 1 sujet discussion. \'Ecrire : Th\`ese / Antith\`ese / Synth\`ese en 10 lignes max. V\'erifier : Chaque argument r\'epond-il \emph{exactement} \`a la question ? & 0 argument hors-sujet. Synth\`ese philosophique (pas politique). \\ \hline
\end{tabularx}
\end{center}

\begin{aretenir}[Le PAP est un Contrat avec Soi-m\^eme]
Il doit \^etre \textbf{\'ecrit}, \textbf{dat\'e}, \textbf{affich\'e} (dans le cahier, sur le mur). Il est \textbf{r\'evisable} chaque dimanche soir (bilan de la semaine : exercices faits ? KPI atteints ? Ajustement semaine suivante). C'est le seul outil qui transforme l'anxi\'et\'e en pilotage.
\end{aretenir}

\courssub{Cl\^oture : La \og{}Fiche M\'emo Bac Philo Tech\fg{} (Recto-Verso)}

\emph{Intuition :} Le jour J, la m\'emoire de travail est satur\'ee. Une fiche A5 (recto-verso), apprise par c\oe ur, lib\`ere le cortex pr\'efrontal pour la pens\'ee.

\begin{exemplebox}[FICHE M\'EMO BAC PHILO TECH --- RECTO : LE CHRONO \& LA STRUCTURE]
\textbf{$\blacktriangleright$ CHRONO 1h30 (Le \og{}Tempo\fg{} Gagnant)}
\begin{itemize}
    \item \textbf{0h00-0h15 (15\%)} : \textbf{LECTURE/ANALYSE} (3 lectures : Th\`ese / Structure+Concepts / Prob+Plan). \textit{Brouillon obligatoire.}
    \item \textbf{0h15-0h30 (15\%)} : \textbf{INTRO AU PROPRE} (Accroche / Pr\'esent / Th\`ese / \textbf{Prob} / Plan). \textit{Ne pas recopier le brouillon, r\'ediger.}
    \item \textbf{0h30-1h10 (40\%)} : \textbf{D\'EVELOPPEMENT} (2 parties $\times$ 20 min). \textit{1 paragraphe = 1 id\'ee + Citation + Concept + Exemple + Transition. Sauter lignes !}
    \item \textbf{1h10-1h20 (10\%)} : \textbf{CONCLUSION} (Bilan / R\'eponse Prob / Ouverture).
    \item \textbf{1h20-1h30 (10\%)} : \textbf{RELECTURE GRILLE} (Citations ? Concepts ? Plan visible ? Accords ? Lisibilit\'e ?).
\end{itemize}

\textbf{$\blacktriangleright$ STRUCTURE TYPE INTRO (4 pts faciles)}
\begin{enumerate}
    \item \textbf{Accroche} : Actualit\'e / Exemple Cameroun / Citation courte / D\'efinition paradoxe.
    \item \textbf{Pr\'esentation} : Auteur, \OE uvre (date), Th\`eme g\'en\'eral, \textbf{Th\`ese} (en 1 phrase : \og{}L'auteur soutient que...\fg{}).
    \item \textbf{Probl\'ematique} : \textbf{Question philosophique} (tension, paradoxe) issue de la th\`ese. \textit{Ex: \og{}La discipline lib\`ere-t-elle ou ali\`ene-t-elle ?\fg{}}
    \item \textbf{Annonce de Plan} : \textit{\og{}Nous verrons d'abord que... (I), puis que... (II), pour enfin montrer que... (III).\fg{}}
\end{enumerate}

\textbf{$\blacktriangleright$ STRUCTURE TYPE CONCLUSION (4 pts faciles)}
\begin{enumerate}
    \item \textbf{Bilan} : \textit{\og{}Au terme de cette analyse, il appara\^it que...\fg{}} (R\'eponse directe \`a la prob.).
    \item \textbf{Nuance} : \textit{\og{}Toutefois, cette r\'eponse doit \^etre nuanc\'ee car...\fg{}}
    \item \textbf{Ouverture} : Autre auteur prog. / Notion li\'ee / Actualit\'e Tech / Perspective \'ethique. \textit{Ex: \og{}Cette tension discipline/libert\'e rejoint la question de l'autonomie chez Kant...\fg{}}
\end{enumerate}
\end{exemplebox}

\begin{exemplebox}[FICHE M\'EMO BAC PHILO TECH --- VERSO : MOTS-CL\'ES, CONNECTEURS \& AUTO-\'EVAL]
\textbf{$\blacktriangleright$ CONNECTEURS LOGIQUES (Le \og{}Ciment\fg{} du Dev - 2 pts structure)}
\begin{center}
\begin{tabular}{|l|l|}\hline
\textbf{Fonction} & \textbf{Mots-cl\'es (Variez !)} \\ \hline
\textbf{Encha\^inement} & Premièrement, En premier lieu, D'une part... D'autre part, Par ailleurs, De surcro\^it, En outre. \\ \hline
\textbf{Opposition / Nuance} & Cependant, Toutefois, N\'eanmoins, Certes... mais, En revanche, Au contraire, Si... toutefois. \\ \hline
\textbf{Cons\'equence / Synth\`ese} & Par cons\'equent, D\`es lors, C'est pourquoi, En d\'efinitive, En somme, Au total, Ultimement. \\ \hline
\textbf{Illustration / Pr\'ecision} & En effet, Par exemple, Notamment, Autrement dit, Plus pr\'ecis\'ement, C'est le cas de... \\ \hline
\end{tabular}
\end{center}

\textbf{$\blacktriangleright$ VOCABULAIRE PHILO \og{}HAUTE VALEUR\fg{} (A placer naturellement)}
\cle{Probl\'ematiser}, \cle{Conceptualiser}, \cle{Argumenter}, \cle{Nuancer}, \cle{Pr\'esupposer}, \cle{Impliquer}, \cle{Contredire}, \cle{R\'efuter}, \cle{Distinction} (ex: fait/droit, l\'egal/l\'egitime, condition/conditionn\'e), \cle{Paradoxe}, \cle{Tension}, \cle{Aporie}.

\textbf{$\blacktriangleright$ GRILLE AUTO-\'EVAL EXPRESS (Derni\`eres 5 min --- Cochez mentalement)}
\begin{itemize}
    \item [ ] \textbf{Intro :} Accroche ? Auteur/\OE uvre/Th\`ese ? \textbf{Prob (?) } Plan (I, II, III) ?
    \item [ ] \textbf{Dev :} Paragraphes visibles (saut ligne + retrait) ? \textbf{Citations exactes (guillemets) ?} Concepts d\'efinis ? Arguments textuels + hors-texte ? \textbf{Transitions (mots-cl\'es) ?} Plan respect\'e ?
    \item [ ] \textbf{Conclu :} Bilan = R\'eponse Prob ? Ouverture pertinente ?
    \item [ ] \textbf{Langue :} Relu \`a l'envers (chasse aux accords) ? \'Ecriture lisible ? Ratures propres ?
\end{itemize}
\textbf{$\blacktriangleright$ RAPPEL STRAT\'EGIQUE :} \textit{\og{}Le correcteur a 50 copies/heure. La lisibilit\'e (sauts de lignes, paragraphes, plan apparent) vaut 1-2 pts inconscients. Soigner la forme, c'est respecter le correcteur.\fg{}} (Voir \textit{Le Saviez-Vous ?})
\end{exemplebox}

\begin{savaistu}[Le Correcteur Invisible : La Psychologie de la Notation]
Des \'etudes en psychologie cognitive de l'\'evaluation montrent que la \textbf{fluidit\'e de traitement} (processing fluency) influence la note. Une copie a\'er\'ee, structur\'ee (titres I/II/III visibles dans la marge ou en d\'ebut de paragraphe), sans ratures sales, est trait\'ee plus vite et jug\'ee \og{}meilleure\fg{} \`a contenu \'egal. \textbf{Astuce Pro :} Num\'erotez vos parties \textbf{en gras dans la marge gauche} (I, A, 1 ; II, A, 1...). Le correcteur voit le plan en 3 secondes. C'est 1 point offert pour 10 secondes d'effort.
\end{savaistu}

\begin{methode}[Auto-\'evaluation = Devenir son Propre Correcteur (Le Levier N$^\circ$1)]
\textbf{Principe :} Ne rendez \textbf{JAMAIS} une copie (devoir, bac blanc, bac) sans l'avoir not\'ee vous-m\^eme avec la grille.
\begin{enumerate}
    \item \textbf{Simulez le correcteur :} Mettez la copie de c\^ot\'e 1h (ou nuit). Reprenez-la avec la grille officielle (Recto Verso Fiche M\'emo).
    \item \textbf{Notez impitoyablement :} \og{}Est-ce que \emph{moi}, correcteur, je vois la prob ? Je vois la citation ? Je vois la transition ?\fg{} Soyez plus dur que le jury.
    \item \textbf{Corrigez avant de rendre :} Si vous voyez \og{}Pas de transition I$\rightarrow$II\fg{}, \'ecrivez-la au propre (si temps) ou au stylo rouge en marge (signe de m\'eta-cognition pour le correcteur : \og{}Je l'ai vu\fg{}).
    \item \textbf{R\'esultat :} En 2 semaines de cette pratique (1 copie/semaine), la note monte m\'ecaniquement de 3 \`a 5 points. Vous int\'egrorisez la grille.
\end{enumerate}
\end{methode}

\begin{attention}[Derniers Pi\`eges du Jour J --- Check-list \og{}Anti-Crash\fg{}]
\begin{itemize}
    \item \textbf{Stylo unique (bleu/noir) :} Pas de changement de couleur (sauf correcteur si autoris\'e). Pas de crayon \`a papier pour le propre.
    \item \textbf{Pas de \og{}Je pense que\fg{}, \og{}A mon avis\fg{} :} Langage philosophique impersonnel : \textit{\og{}Il convient de montrer\fg{}, \og{}Force est de constater\fg{}, \og{}L'analyse r\'ev\`ele\fg{}}.
    \item \textbf{Citation = Guillemets fran\c{c}ais \og{} \fg{} + R\'ef\'erence (ligne/paragraphe) :} \og{}L'\'ecole est une machine \`a fabriquer des corps dociles\fg{} (l. 12-13). Pas de citation \og{}flottante\fg{} sans analyse imm\'ediate.
    \item \textbf{Gestion Q3 (Discussion) :} Ne pas faire un \og{}pour/contre\fg{} binaire. Structure : \textbf{Th\`ese} (argument 1 + ex) $\rightarrow$ \textbf{Antith\`ese/Objection} (argument 2 + ex) $\rightarrow$ \textbf{Synth\`ese/D\'epassement} (nouveau concept, distinction, condition de validit\'e).
    \item \textbf{Hors-sujet \og{}Sociologique/Politique\fg{} :} La philo ne demande pas \og{}Que faire au Cameroun ?\fg{} mais \og{}Que penser conceptuellement de X ?\fg{}. Gardez l'ancrage camerounais comme \textbf{exemple illustratif}, pas comme \textbf{th\`ese}.
\end{itemize}
\end{attention}

\begin{exercice}[Entra\^inement Final : \og{}Le Grand Oral Blanc\fg{} (Optionnel mais Recommand\'e)]
Pr\'eparez, sur une fiche bristol, la \textbf{pr\'esentation orale de 5 min} de votre copie type (Copie A id\'eale) :
\begin{enumerate}
    \item Sujet / Th\`ese / Probl\'ematique.
    \item Plan (I, II, III) avec \textbf{1 argument phare + 1 citation par partie}.
    \item R\'eponse synth\'etique \`a la prob. + Ouverture.
\end{enumerate}
Pr\'esentez-la \`a un pair / prof en 5 min chrono. C'est l'ultime test de clart\'e : si vous ne pouvez pas l'expliquer simplement en 5 min, vous ne l'avez pas assez structur\'ee pour l'\'ecrit.
\end{exercice}

\begin{corrige}[Exercice : Le Grand Oral Blanc]
\textbf{Attendu :} L'\'el\`eve tient en 5 min. Regard droit. Pas de lecture. Vocabulaire pr\'ecis. Plan visible sur le bristol (mots-cl\'es seulement). Le pair \'evalue avec la grille orale : Clart\'e / Structure / Pertinence / Ma\^itrise du temps. Objectif : \textbf{Valid\'e} = Pr\^et pour l'\'ecrit.
\end{corrige}

\par\bigskip\noindent
\textbf{\cle{Mot de la fin :}} La philosophie au Bac Tech, ce n'est pas \og{}avoir des id\'ees\fg{}. C'est \textbf{montrer qu'on sait les ordonner, les \'etayer, les discuter} dans un cadre contraint. Vous avez la m\'ethode (Sections 1-4), vous avez l'entra\^inement (Section 5), vous avez la grille (Fiche M\'emo). \textbf{Faites confiance \`a votre architecture.} Le jour J, appliquez le tempo. Respirez. \'Ecrivez. Relisez. \textbf{Bonne chance \`a toutes et tous pour le Probatoire / Baccalaur\'eat Technique 2027 !}

\newpage

\coursec{Activité d'intégration}

\coursec{Consolidation des acquis de la méthodologie de l'exercice sur texte philosophique}

\courssub{Objectifs de l'activité}

À l'issue de cette activité d'intégration, l'élève de Terminale technique doit être capable de :

\begin{itemize}
\item mobiliser de manière simultanée l'ensemble des savoir-faire de l'exercice sur texte philosophique : lecture analytique, dégagement de la thèse, problématisation, analyse des arguments, discussion critique et conclusion ouverte ;
\item rédiger un mémoire argumenté d'une page respectant la structure canonique du développement de l'explication de texte (thèse, analyse, discussion) ;
\item adapter un raisonnement philosophique à une situation concrète de la vie publique camerounaise (débat sur la sécurité numérique) ;
\item s'exprimer oralement de façon claire et argumentée devant une assemblée, et évaluer la production d'autrui à l'aide d'une grille de critères méthodologiques ;
\item justifier une démarche, une réponse ou une conclusion, comme l'exige l'épreuve de Philosophie au Probatoire et au Baccalauréat technique.
\end{itemize}

\courssub{Situation}

\textbf{Communiqué de la Commission Nationale des Droits de l'Homme et des Libertés (CNDHL), Yaoundé.}

Dans le cadre de l'examen du projet de loi n° 2026/014 portant \emph{Sécurité Numérique et Protection des Données}, la CNDHL organise une audition citoyenne. Ce projet prévoit notamment : l'installation de caméras intelligentes dans les grandes villes (Yaoundé, Douala, Bafoussam), la conservation des données de connexion des internautes pendant douze mois, et la création d'une cellule de lutte contre les \emph{fake news} sur les réseaux sociaux. Les organisations de la société civile s'inquiètent d'une possible atteinte aux libertés. Pour éclairer les débats, la Commission a soumis aux participants le texte philosophique suivant (20 lignes) :

\begin{exemplebox}[Texte soumis à l'audition : « Surveiller pour protéger ? »]
« Toute société moderne fait face à un paradoxe : pour garantir la liberté de chacun, l'État doit exercer un contrôle qui semble la menacer. Comme l'a montré Foucault, le pouvoir contemporain ne se contente plus d'interdire : il surveille, enregistre, classe. La caméra et le fichier numérique sont les nouveaux instruments d'une discipline douce, d'autant plus efficace que l'individu, se sachant observé, finit par se surveiller lui-même. Pourtant, comme le souligne Habermas, l'espace public démocratique suppose que les citoyens puissent délibérer librement, sans crainte d'être fichés pour leurs opinions. Or la désinformation qui circule sur les réseaux sociaux — rumeurs, fausses nouvelles, manipulations — détruit précisément les conditions d'une discussion rationnelle : on ne peut délibérer sur la base de mensonges. La sécurité est donc une condition de la liberté, mais une surveillance sans limites la détruit. Tout l'art du législateur consiste alors à trouver le point d'équilibre : protéger les citoyens sans les enfermer dans une transparence totale où nul ne serait plus libre de penser autrement. La question n'est pas de choisir entre sécurité et liberté, mais de déterminer jusqu'où l'une peut aller sans abolir l'autre. »
\end{exemplebox}

Quatre groupes sont conviés à l'audition, chacun portant un point de vue déterminé :

\begin{center}
\begin{tabularx}{\linewidth}{|l|X|X|}\hline
\rowcolor{popblueL} \textbf{Rôle} & \textbf{Position attendue} & \textbf{Références suggérées} \\ \hline
Gouvernement & Défendre le projet de loi : la sécurité conditionne l'exercice des libertés. & Hobbes : l'État protège contre l'insécurité. \\ \hline
Opposition & Critiquer le projet : risque d'État policier et de censure. & Foucault : la surveillance discipline les conduites. \\ \hline
Société civile / ONG & Défendre les libertés individuelles et la vie privée. & Locke : droits naturels, limites du pouvoir. \\ \hline
Experts juristes & Poser les conditions d'un encadrement légal équilibré. & Habermas : espace public, droit positif. \\ \hline
\end{tabularx}
\end{center}

\courssub{Consignes}

\textbf{Séance 1 — Préparation en groupes de quatre.} Chaque groupe rédige un \emph{Mémoire} d'une page destiné à l'audition publique. Travail guidé :

\begin{enumerate}
\item \textbf{Lecture analytique.} Soulignez dans le texte la thèse principale et relevez les deux arguments qui la soutiennent. Formulez la thèse en une phrase, avec vos propres mots.
\item \textbf{Problématisation.} Transformez la question du texte en une question directrice claire, de forme interrogative, qui pourrait servir de problématique à l'audition (par exemple : « Jusqu'où l'État peut-il surveiller les citoyens sans détruire leur liberté ? »).
\item \textbf{Analyse.} Expliquez en quelques lignes le rôle des références à Foucault et à Habermas dans l'économie du texte : que démontre chacune ?
\item \textbf{Discussion.} Selon le rôle attribué à votre groupe, discutez la thèse du texte : approuvez-la, nuancez-la ou réfutez-la, en mobilisant au moins \emph{un} auteur de la fiche « Rôles et arguments » et \emph{un} exemple camerounais concret (usage des réseaux sociaux, rumeurs en période électorale, vidéosurveillance urbaine).
\item \textbf{Rédaction.} Composez le mémoire selon la structure : (a) thèse du texte, (b) analyse des arguments, (c) discussion critique argumentée, (d) conclusion personnelle justifiée.
\end{enumerate}

\textbf{Séance 2 — Audition publique.} Chaque groupe présente son mémoire en cinq minutes. Les autres groupes, réunis en jury, notent la prestation sur la grille simplifiée (chaque critère sur 5 points) : \cle{fidélité au texte}, \cle{pertinence de la discussion}, \cle{qualité argumentative}, \cle{références philosophiques}. Le professeur, président de séance, synthétise les débats et évalue la méthodologie.

\courssub{Ressources à mobiliser}

\begin{aretenir}[Rappels méthodologiques utiles]
\begin{itemize}
\item \textbf{Structure de l'explication de texte} : introduction (présentation de l'auteur et du texte, dégagement de la thèse, problématique, annonce du plan) ; développement (analyse : expliquer la thèse et ses arguments ; discussion : évaluer la thèse avec des arguments et des références) ; conclusion (bilan et réponse personnelle justifiée).
\item \textbf{Problématique} : une question directrice, de forme interrogative, qui exprime le problème philosophique du texte — ni trop large, ni hors sujet.
\item \textbf{Discussion} : on ne se contente pas de paraphraser ; on confronte la thèse à d'autres doctrines (Hobbes, Locke, Foucault, Habermas) et à l'expérience concrète.
\item \textbf{Exigences de forme} : langue correcte, connecteurs logiques (\emph{d'abord, en effet, cependant, par conséquent}), citations brèves et exactes, aucune copie du texte.
\end{itemize}
\end{aretenir}

\begin{attention}[Pièges fréquents à éviter]
Ne confondez pas \emph{résumé} et \emph{explication} : recopier le texte ne vaut aucun point d'analyse. Ne plaidez pas pour votre rôle sans passer par le texte : la discussion doit d'abord s'appuyer sur la thèse dégagée. Enfin, une opinion personnelle non justifiée par un argument ou une référence philosophique n'est pas recevable à l'audition comme au Baccalauréat.
\end{attention}

\courssub{Démarche et corrigé commenté}

\begin{exoresolu}[Exemple de mémoire réussi : groupe « Société civile / ONG »]
Énoncé : rédiger le mémoire du groupe Société civile, chargé de discuter la thèse du texte à partir de Locke et d'un exemple camerounais.
\tcblower
\textbf{Étape 1 — Thèse du texte (fidélité).} Le texte soutient que la sécurité est une condition de la liberté, mais qu'une surveillance sans limites la détruit : le législateur doit donc fixer un équilibre entre les deux. \emph{Commentaire : la thèse est reformulée en une phrase, sans paraphrase.}

\textbf{Étape 2 — Problématique.} « Jusqu'où la sécurité numérique peut-elle aller sans abolir la liberté des citoyens ? » \emph{Commentaire : question directrice, de forme interrogative, fidèle au problème du texte.}

\textbf{Étape 3 — Analyse des arguments.} Premier argument : avec Foucault, la surveillance moderne est une « discipline douce » — l'internaute qui se sait observé s'autocensure, ce qui menace la liberté de penser. Second argument : avec Habermas, la délibération démocratique exige un espace public libre, que les \emph{fake news} corrompent ; une régulation est donc nécessaire. \emph{Commentaire : chaque référence est expliquée dans son rôle, pas simplement citée.}

\textbf{Étape 4 — Discussion selon le rôle (ONG, référence : Locke).} Nous approuvons le diagnostic du texte, mais nous contestons que le projet de loi respecte l'équilibre annoncé. Selon Locke, les individus possèdent des droits naturels — dont la liberté et la propriété de soi — que le pouvoir reçoit mission de protéger, non de violer ; un pouvoir qui s'en affranchit devient arbitraire. Or conserver douze mois les données de connexion de tous les internautes, sans contrôle d'un juge, institue une surveillance généralisée de citoyens présumés innocents : c'est précisément la « transparence totale » que le texte lui-même dénonce. Exemple camerounais : lors des campagnes électorales, des rumeurs virales ont certes semé le trouble ; mais couper ou surveiller massivement les réseaux pénalise aussi les étudiants, les commerçants en ligne de Douala et les familles qui communiquent par messagerie. La lutte contre la désinformation ne justifie pas de traiter chaque citoyen en suspect. \emph{Commentaire : la discussion confronte la thèse à Locke et à un fait national, avec un raisonnement enchaîné par des connecteurs.}

\textbf{Étape 5 — Conclusion justifiée.} Nous demandons que la conservation des données soit soumise à l'autorisation préalable d'une autorité judiciaire indépendante et limitée aux enquêtes ciblées. Ainsi la sécurité restera la \emph{condition} de la liberté sans devenir son \emph{tombeau} : l'équilibre que le texte appelle de ses vœux sera juridiquement garanti, et non simplement proclamé. \emph{Commentaire : réponse personnelle, argumentée, qui répond exactement à la problématique.}
\end{exoresolu}

\courssub{Bilan de l'activité}

Cette simulation du « Tribunal du Texte » a permis de mettre en évidence trois acquis essentiels. D'abord, la méthodologie de l'exercice sur texte n'est pas un rituel scolaire abstrait : lire fidèlement, problématiser, analyser puis discuter sont les gestes mêmes du citoyen qui prend part à un débat public éclairé. Ensuite, l'activité a montré qu'une même thèse peut être discutée de manière contradictoire mais rigoureuse : gouvernement, opposition, ONG et juristes mobilisent Hobbes, Foucault, Locke ou Habermas pour construire des arguments différents à partir du même texte — preuve que la qualité d'une copie tient à la solidité de l'argumentation, non à l'opinion défendue. Enfin, l'évaluation croisée par grille a développé la méta-cognition des élèves : en notant la fidélité au texte, la pertinence de la discussion et la qualité argumentative chez les autres, chacun a intériorisé les critères du correcteur au Probatoire et au Baccalauréat technique. L'élève est désormais prêt à affronter seul, en temps limité, l'exercice sur texte philosophique.

\newpage

\coursec{Méthodes et exercices résolus}

\coursec{Consolidation des acquis de la méthodologie de l'exercice sur texte philosophique}

\courssub{Diagnostic \& Rappel des Exigences Officielles}

\begin{methode}[Analyser la consigne et identifier les attendus de l'épreuve]
\textbf{Objectif :} Ne pas confondre vitesse et précipitation. Les 10 premières minutes décident de la note finale.
\begin{enumerate}
    \item \textbf{Lire le texte deux fois :} 1\textsuperscript{re} lecture \og découverte \fg (sans stylo), 2\textsuperscript{e} lecture \og active \fg (surlignage des concepts-clés, connecteurs logiques, thèse principale).
    \item \textbf{Décortiquer le paratexte :} Auteur, œuvre, date, collection. Cela situe le \cle{contexte historico-philosophique} (ex : Descartes \emph{Discours de la méthode} 1637 $\rightarrow$ rupture scolastique / naissance science moderne).
    \item \textbf{Identifier la structure du texte :} Repérer les \cle{étapes du raisonnement} (thèse, arguments, objections, réponses, conclusion, distinctions conceptuelles). Numéroter les paragraphes.
    \item \textbf{Lister les concepts opératoires :} Mots techniques (ex : \cle{cogito}, \cle{liberté}, \cle{aliénation}, \cle{progrès}, \cle{technique}). Définir chacun \emph{in situ} (dans le texte) et \emph{ex situ} (dans la doctrine de l'auteur).
    \item \textbf{Formuler la thèse centrale en une phrase :} \og L'auteur soutient que... parce que... \fg.
    \item \textbf{Analyser les questions :} Distinguer \og Expliquez \fg (restitution interne), \og Montrez comment \fg (analyse de la démonstration), \og Discutez \fg (évaluation critique externe).
\end{enumerate}
\end{methode}

\begin{exoresolu}[Diagnostic sur un extrait type Bac Tech — Kant, \emph{Fondement de la métaphysique des mœurs} (Préface)]
\textbf{Énoncé :} \og La morale ne peut jamais être tirée de l'expérience... C'est une philosophie pure... \fg (Extrait : distinction empirique / pur, nécessité de la métaphysique des mœurs).
\textbf{Questions :}
1. Expliquez la distinction entre \og philosophie empirique \fg et \og philosophie pure \fg.
2. Montrez comment l'auteur justifie la nécessité d'une \og métaphysique des mœurs \fg.
3. \og La morale est-elle indépendante de l'expérience ? \fg Discutez.

\tcblower
\textbf{Solution détaillée :}

\textbf{Diagnostic (Brouillon structuré) :}
\begin{itemize}
    \item \textbf{Auteur/Contexte :} Emmanuel Kant (1724-1804), Critique de la raison pure (1781) $\rightarrow$ Fondement (1785). Projet : fonder la morale \emph{a priori}, contre l'utilitarisme (Hume) et l'eudémonisme.
    \item \textbf{Thèse centrale :} La loi morale doit être \cle{formelle}, \cle{universelle} et \cle{nécessaire} $\rightarrow$ donc indépendante de l'expérience (contingente).
    \item \textbf{Structure :}
    \begin{enumerate}
        \item §1 : Définition philosophie empirique (fondée sur l'expérience) vs pure (fondée sur la raison \emph{a priori}).
        \item §2 : Application à la morale : \og métaphysique des mœurs \fg = partie pure de la morale. L'expérience ne donne que des cas, pas la loi.
        \item §3 : Danger de l'empirisme moral (hétérogénéité des mobiles, absence de nécessité).
        \item §4 : Nécessité pratique : éviter le \cle{circulus vitiosus} (cercle vicieux) où la raison cherche le mobile dans l'expérience.
    \end{enumerate}
    \item \textbf{Concepts-clés :} \emph{A priori / A posteriori}, \cle{Impératif catégorique}, \cle{Autonomie}, \cle{Hétéronomie}, \cle{Bonheur} (concept empirique) vs \cle{Devoir} (concept pur).
\end{itemize}

\textbf{Réponse Q1 (Expliquez) :}
La \cle{philosophie empirique} (ex : physique, psychologie, anthropologie pragmatique) tire ses principes de l'expérience (\emph{a posteriori}). Ses jugements sont contingents (peuvent être autrement). La \cle{philosophie pure} (ex : logique, métaphysique, mathématiques) tire ses principes de la raison pure (\emph{a priori}). Ses jugements sont nécessaires et universels. En morale, l'empirique donne l'\cle{anthropologie} (ce que font les hommes), le pur donne la \cle{morale} (ce que les hommes \emph{doivent} faire).

\textbf{Réponse Q2 (Montrez comment) :}
Kant procède par \cle{nécessité logique} :
\begin{enumerate}
    \item Prémisse majeure : La loi morale s'adresse à \emph{tous} les êtres raisonnables (universalité) et \emph{nécessairement} (obligation inconditionnelle).
    \item Prémisse mineure : L'expérience ne fournit que des jugements contingents (particuliers, variables selon les inclinations, le climat, l'éducation).
    \item Conclusion : L'expérience est \emph{inapte} à fonder la loi morale. Il faut une \cle{métaphysique des mœurs} (philosophie pure appliquée au pratique) pour isoler le principe \emph{a priori} de la moralité (la forme de la loi : universalisabilité).
\end{enumerate}
Il montre aussi le \textbf{risque pratique} : sans métaphysique, la morale devient \cle{hétéronome} (soumise à l'inclination), détruisant la \cle{dignité} de l'homme.

\textbf{Réponse Q3 (Discutez) :}
\textbf{Thèse (Kant) :} Oui, indépendance radicale. Le mobile moral pur est le \cle{respect de la loi}. Tout mobile empirique (bonheur, pitié, intérêt) corrompt la moralité (hétéronomie). Exemple : le marchand honnête par intérêt vs par devoir.
\textbf{Antithèse (Aristote / Hume / Nietzsche / Contextuel Cameroun) :} Non, la morale s'enracine dans l'expérience.
\begin{itemize}
    \item \emph{Aristote :} Vertu = habitude (\emph{ethos}) acquise par la pratique, l'éducation, le milieu (polis). Pas de morale sans \cle{phronèsis} (prudence) situationnelle.
    \item \emph{Hume :} La raison est \og l'esclave des passions \fg. Le fondement est le \cle{sentiment moral} (sympathie), empirique.
    \item \emph{Réalisme africain (Eboussi Boulaga) :} La morale est \cle{communautaire}, vécue dans le \og nous \fg ancestral, transmise par l'initiation et l'expérience des aînés, non déduite d'un \emph{a priori} abstrait.
\end{itemize}
\textbf{Synthèse :} Kant a raison sur la \cle{forme} (exigence d'universalité = idéal régulateur), mais tort sur la \cle{matière} (application concrète). Une morale \emph{sans} expérience est \og vide \fg (Kant lui-même : \og Concepts sans intuitions sont vides \fg). Une morale \emph{par} l'expérience seule est \og aveugle \fg (relativisme). La \cle{méthode} rigoureuse : distinguer \emph{fondation} (idéal \emph{a priori}) et \emph{application} (jugement \emph{a posteriori} / prudence).

\textbf{Conclusion de l'exercice :} Le diagnostic structure la copie : Q1 = §1, Q2 = §2-4, Q3 = Dépassement critique structuré (Thèse/Antithèse/Synthèse).
\end{exoresolu}

\courssub{L'Atelier Problématisation : Du Sujet à la Question Directrice}

\begin{methode}[Construire la problématique en 4 temps (Méthode \og Tension Dialectique \fg)]
\textbf{Principe :} Une problématique n'est pas un résumé, c'est la mise en lumière d'un \cle{conflit conceptuel} que le texte tente de résoudre.
\begin{enumerate}
    \item \textbf{Extraire les concepts antagonistes :} Identifier les paires opposées dans le texte (ex : \cle{Nature} vs \cle{Culture}, \cle{Liberté} vs \cle{Déterminisme}, \cle{Individu} vs \cle{État}, \cle{Technique} vs \cle{Humanisme}, \cle{Théorie} vs \cle{Pratique}).
    \item \textbf{Formuler la thèse de l'auteur :} \og Pour [Auteur], [Concept A] prime sur [Concept B] car [Raison centrale du texte] \fg.
    \item \textbf{Formuler l'antithèse plausible (l'objection interne ou externe) :} \og Mais on peut objecter que [Concept B] est tout aussi essentiel car [Argument contraire fort] \fg.
    \item \textbf{Rédiger la Question Directrice (Problématique) :} Syntaxe type : \og \textbf{Dans quelle mesure} [Thèse de l'auteur] \textbf{ne risque-t-elle pas de nier / d'oublier / de réduire} [Antithèse] \textbf{?} \fg OU \og \textbf{Comment} [Concept A] \textbf{peut-il être concilié avec} [Concept B] \textbf{dans la perspective de} [Auteur] \textbf{?} \fg
    \item \textbf{Valider :} La question appelle-t-elle un développement en 2 ou 3 parties ? (Oui = Bonne problématique).
\end{enumerate}
\end{methode}

\begin{exoresolu}[Problématisation : Texte sur \og La Technique et l'Humanité \fg (Type Bac Tech — inspiration Heidegger / Ellul / Simondon)]
\textbf{Énoncé :} Extrait affirmant : \og La technique n'est pas un simple outil neutre ; elle structure notre rapport au monde et risque d'enfermer l'homme dans un \emph{enframing} (\textit{Gestell}) où tout devient \og fonds disponible \fg (ressource). \fg
\textbf{Consigne :} \og Problématisez le texte et proposez un plan détaillé. \fg

\tcblower
\textbf{Application de la méthode :}

\textbf{1. Concepts antagonistes :}
\begin{itemize}
    \item \cle{Technique} (Moyen, Instrument, Puissance, \textit{Gestell}, Détermination) \textbf{vs} \cle{Humanité} (Fin, Liberté, Autonomie, Sens, Réceptivité).
    \item \cle{Disponibilité} (Monde comme stock/ressource) \textbf{vs} \cle{Vérité/Être} (Monde comme don/sens).
\end{itemize}

\textbf{2. Thèse de l'auteur (Heideggerien) :}
\og La technique moderne, par son essence (\textit{Gestell}), \textbf{menace} l'humanité en la réduisant à n'être que le gestionnaire des ressources, oubliant la \cle{question du sens} et la \cle{libération} (gelassenheit). \fg

\textbf{3. Antithèse (Simondon / Technophile / Contexte Cameroun) :}
\og La technique est \textbf{libératrice} : elle soulage la peine (machines agricoles, santé), connecte les hommes (numérique), permet l'expression culturelle. L'aliénation vient du \cle{système socio-économique} (capitalisme), pas de la technique en soi. Au Cameroun, le téléphone mobile a plus \og humanisé \fg (banque mobile, éducation) qu'aliéné. \fg

\textbf{4. Question Directrice (Problématique) :}
\og \textbf{Dans quelle mesure l'essence de la technique moderne comme mise en réserve (Gestell) compromet-elle irréversiblement l'autonomie humaine, ou bien une appropriation critique et éthique (Simondon, éthique de la responsabilité) peut-elle réorienter la technique vers l'émancipation ?} \fg

\textbf{5. Plan détaillé validé (3 parties dialectiques) :}

\begin{center}
\begin{tabularx}{\linewidth}{|l|X|}\hline
\textbf{Partie} & \textbf{Contenu / Arguments / Auteurs / Exemples} \\ \hline
\textbf{I. La technique comme menace : l'enfermement dans le \og Gestell \fg (Thèse du texte)} &
\begin{itemize}
    \item Essence de la technique $\neq$ instrumentalité (Heidegger). \textit{Gestell} = mise en réserve généralisée.
    \item Conséquence : \cle{Réification} de l'homme (ressource humaine), \cle{calculabilité} totale (Big Data, IA), perte du \cle{sens} (Nihilisme technique).
    \item Exemple : Algorithmes de notation sociale / Gestion RH par IA / Déforestation industrielle (Sud-Cameroun) : la forêt devient \og bois d'œuvre \fg.
\end{itemize} \\ \hline
\textbf{II. La technique comme médiation libératrice : l'aliénation n'est pas une fatalité (Antithèse)} &
\begin{itemize}
    \item \cle{Simondon} : Technique = \og individuation \fg. Objet technique porte du sens, de la culture. \cle{Concret} vs \cle{Abstrait}.
    \item \cle{Arendt} : \emph{Homo faber} vs \emph{Animal laborans}. Technique libère de la nécessité biologique pour la \cle{politique} (action libre).
    \item Exemples camerounais : Mobile Money (inclusion financière), Télémedecine (accès soins zones rurales), Machines agricoles (pénibilité).
\end{itemize} \\ \hline
\textbf{III. Vers une éthique de la technique : responsabilité et \og Gelassenheit \fg (Synthèse)} &
\begin{itemize}
    \item Dépasser le dualisme : Ni technophobie, ni technophilie béate. \cle{Jonas} : \og Principe responsabilité \fg (heuristique de la peur).
    \item \cle{Heidegger} : \textit{Gelassenheit} (détachement / disponibilité ouverte) = attitude poétique face à la technique.
    \item \cle{Éthique ingénieur / Régulation} : \og Privacy by design \fg, IA explicable, Transition écologique (énergies renouvelables vs extractivisme).
    \item Conclusion : L'humain reste \cle{berger de l'être} (Heidegger) / \cle{sujet de responsabilité} (Jonas) seulement s'il pense la technique, ne la subit pas.
\end{itemize} \\ \hline
\end{tabularx}
\end{center}
\end{exoresolu}

\courssub{Rédaction Chirurgicale : Introduction, Développement, Conclusion}

\begin{methode}[Le \og Standard Bac \fg : Introduction en 4 temps, Développement en \og Argument-Exemple-Analyse \fg, Conclusion en 2 temps]
\textbf{Introduction (15-20 lignes max, 1 paragraphe unique) :}
\begin{enumerate}
    \item \textbf{Amorce contextuelle (1 phrase) :} Citation courte, fait de société camerounais, actualité tech, ou définition étymologique. \emph{Pas de \og Depuis la nuit des temps \fg.}
    \item \textbf{Sujet \& Thèse du texte (2 phrases) :} \og Dans cet extrait de [Œuvre], [Auteur] soutient que [Thèse centrale résumée]... \fg
    \item \textbf{Problématique (1 phrase claire, interrogative) :} \og Dès lors, [Question Directrice issue de l'atelier précédent] ? \fg
    \item \textbf{Annonce du Plan (1 phrase structurée) :} \og Nous verrons d'abord que [I]... ; puis que [II]... ; enfin que [III]... \fg (Mots-clés : \textbf{D'abord, Ensuite, Enfin}).
\end{enumerate}

\textbf{Développement (Corps de la dissertation / Réponse Q3 "Discutez") :}
\begin{enumerate}
    \item \textbf{1 Paragraphe = 1 Idée Force = 1 Argument majeur.}
    \item \textbf{Structure \og A.E.A. \fg (Argument - Exemple - Analyse) :}
    \begin{itemize}
        \item \textbf{Argument :} Phrase thèse du paragraphe (lien avec problématique).
        \item \textbf{Exemple / Référence :} Auteur précis, concept technique, fait précis (Cameroun, Histoire, Science).
        \item \textbf{Analyse / Lien :} \og Cela montre que... \fg, \og Ce qui prouve que... \fg, \og Cependant... \fg (transition vers paragraphe suivant).
    \end{itemize}
    \item \textbf{Connecteurs logiques obligatoires :} \textbf{En effet, Certes, Cependant, Par conséquent, En outre, De surcroît, En définitive.}
\end{enumerate}

\textbf{Conclusion (10-15 lignes) :}
\begin{enumerate}
    \item \textbf{Bilan synthétique (2-3 phrases) :} Reprendre le fil des 3 parties sans liste à puces. \og Au terme de cette analyse, il apparaît que [I]... mais que [II]... conduisant à [III]... \fg
    \item \textbf{Ouverture (1 phrase percutante) :} Question nouvelle, autre auteur, actualité brûlante, perspective éthique. \og Cette réflexion invite à repenser [Concept] à l'ère de [IA / Crise écologique / Gouvernance numérique]... \fg
\end{enumerate}
\end{methode}

\begin{exoresolu}[Rédaction Intro + Conclusion sur le sujet : \og L'État est-il le garant de la liberté ? \fg (Programme : Société/État)]
\textbf{Contexte :} Dissertation type Bac Tech. Thèse : Hobbes/Locke (État protecteur). Antithèse : Rousseau/Anarchistes/Thoreau (État oppresseur). Synthèse : État de droit / Démocratie participative / Société civile.

\tcblower
\textbf{Introduction rédigée :}
\og \textbf{[Amorce]} « L'homme naît libre, et partout il est dans les fers » (Rousseau, \emph{Du contrat social}). Cette parole prophétique résonne singulièrement dans nos États africains post-coloniaux où la promesse de l'indépendance a parfois cédé la place à l'autoritarisme ou à la faillite de l'État de droit. \textbf{[Sujet/Thèse]} Dans la tradition contractuelle (Hobbes, Locke), l'État naît précisément pour \emph{garantir} la liberté par la sécurité : \og Le but de la loi n'est pas d'abolir ou de restreindre, mais de préserver et d'élargir la liberté \fg (Locke). \textbf{[Problématique]} \textbf{Pourtant, l'histoire montre que l'institution étatique peut basculer dans la tyrannie : dans quelle mesure l'État, né pour protéger la liberté, ne porte-t-il pas en lui le germe de sa propre négation ?} \textbf{[Plan]} Nous analyserons d'abord comment l'État fonde la liberté par le droit (I) ; nous montrerons ensuite les dérives totalitaires ou prédatrices qui en font le premier prédateur de la liberté (II) ; enfin, nous verrons que seule une démocratie substantielle, contrôlée par une société civile vigilante, permet de résoudre cette aporie (III). \fg

\textbf{Conclusion rédigée :}
\og \textbf{[Bilan]} Au terme de cette démarche, il apparaît que l'État est bien le \textbf{condition de possibilité} de la liberté publique (I) : sans monopole de la violence légitime et sans cadre juridique, c'est la loi du plus fort (Hobbes). Cependant, l'État \textbf{concentre un pouvoir} qui, sans contre-pouvoirs effectifs, se retourne contre la liberté (II) : surveillance de masse, lois liberticides, corruption prédatoire (ex : détournements de fonds publics au détriment des services sociaux). \textbf{[Synthèse]} La solution ne réside pas dans la suppression de l'État (utopie anarchiste) mais dans sa \textbf{constitutionnalisation rigoureuse} (III) : séparation des pouvoirs, justice indépendante, presse libre, participation citoyenne (budgets participatifs à Yaoundé ou Douala). \textbf{[Ouverture]} Dès lors, la liberté ne se \og donne \fg pas par l'État ; elle se \og conquiert \fg et se \og veille \fg quotidiennement par le citoyen, rejoignant l'exigence kantienne d'une \og sortie de la minorité \fg : \og \textbf{Sapere aude !} \fg (Ose savoir / Ose être libre). \fg
\end{exoresolu}

\courssub{Maîtrise des Questions Types \& Pièges Spécifiques Bac Tech}

\begin{methode}[Répondre aux 3 questions canoniques de l'exercice sur texte (Q1, Q2, Q3)]
\textbf{Règle d'or :} \textbf{Ne jamais mélanger les registres.} Q1 = \emph{Exégèse} (dedans le texte). Q2 = \emph{Analyse logique} (structure du texte). Q3 = \emph{Philosophie personnelle} (dehors le texte, mais guidé par lui).
\begin{enumerate}
    \item \textbf{Q1 : \og Expliquez / Précisez le sens de... \fg (ou \og Que veut dire l'auteur quand il écrit... ? \fg)}
    \begin{itemize}
        \item \textbf{Consigne :} Reformuler \emph{avec ses mots} le passage visé, en explicitant les concepts techniques \emph{dans le contexte du texte}.
        \item \textbf{Structure :} \og L'auteur affirme que [Concept A] signifie [Définition contextuelle]. Il le distingue de [Concept B] car [Raison textuelle]. Par exemple, dans le texte, il écrit : « [Citation courte] » ce qui montre que... \fg
        \item \textbf{Piège :} Faire du hors-texte, donner son avis, faire un cours sur l'auteur.
    \end{itemize}
    \item \textbf{Q2 : \og Montrez comment l'auteur démontre que... \fg / \og Analysez le raisonnement... \fg}
    \begin{itemize}
        \item \textbf{Consigne :} Reconstruire la \cle{chaîne logique} (Syllogisme, Arguments d'autorité, Analogie, Réduction à l'absurde, Distinction conceptuelle).
        \item \textbf{Structure :} \og L'auteur procède en [X] étapes : 1. Il pose que [Prémisse 1]... 2. Il en déduit que [Inférence]... 3. Il illustre par [Exemple/Analogie]... 4. Il conclut que [Thèse]. La force du raisonnement repose sur [Distinction clé / Concept opératoire]. \fg
        \item \textbf{Piège :} Refaire l'explication de texte (Q1) ou donner son avis (Q3).
    \end{itemize}
    \item \textbf{Q3 : \og Discutez / Que pensez-vous de cette thèse ? \fg / \og Cette thèse est-elle convaincante ? \fg}
    \begin{itemize}
        \item \textbf{Consigne :} Mini-dissertation structurée (Thèse/Antithèse/Synthèse) \textbf{en 1 page max}.
        \item \textbf{Structure imposée :}
        \begin{enumerate}
            \item \textbf{Accord partiel (Thèse) :} \og La thèse est solide car [Argument fort du texte + Exemple précis]... \fg
            \item \textbf{Objection rigoureuse (Antithèse) :} \og Cependant, elle pèche par [Limite : oubli, excès, contre-exemple, auteur adverse]... \fg
            \item \textbf{Dépassement (Synthèse) :} \og Dès lors, il faut nuancer : [Distinction / Condition / Nouveau concept]... \fg
        \end{enumerate}
        \item \textbf{Piège majeur Bac Tech :} \og Je pense que... \fg $\rightarrow$ \textbf{Interdit.} Remplacer par \og On peut soutenir que... \fg, \og Il est pertinent de noter que... \fg, \og La thèse résiste à l'épreuve de... \fg.
    \end{itemize}
\end{enumerate}
\end{methode}

\begin{exoresolu}[Application Q1, Q2, Q3 sur : Descartes, \emph{Discours de la méthode}, 3\textsuperscript{e} partie (Morale par provision)]
\textbf{Texte :} \og ...obéir aux lois et coutumes de mon pays... / ...suivre les opinions les plus modérées... / ...me conduire selon les maximes de la religion... \fg

\tcblower
\textbf{Q1 : Expliquez ce que Descartes entend par \og morale par provision \fg.}
\textbf{Solution :}
\og La \cle{morale par provision} (ou \og morale provisoire \fg) est l'ensemble des règles pratiques que Descartes s'impose \textbf{le temps de la reconstruction métaphysique} (le doute hyperbolique). Comme il a détruit ses anciennes opinions (doute) mais n'a pas encore établi la vérité (le \emph{Cogito}, Dieu, le monde), il a besoin de \cle{règles d'action} pour ne pas rester paralysé. Il choisit la \cle{conformité} (lois du pays, religion, opinions modérées) non par conviction absolue, mais par \cle{prudence} et \cle{sagesse} : cela évite l'excentricité, le danger politique et l'erreur morale grossière. C'est une \cle{morale de l'attente}, subordonnée à la future \cle{morale par principe} (fondée sur la raison claire et distincte). \fg

\textbf{Q2 : Montrez comment Descartes justifie le choix de la \og conformité \fg (lois, coutumes, religion) comme règle d'action.}
\textbf{Solution :}
\og Descartes justifie ce choix par un \textbf{raisonnement prudentialiste (coût/bénéfice) en 3 temps} :
\begin{enumerate}
    \item \textbf{Constat d'incertitude :} Ma raison doute de tout (coutumes, religion comprises). Je ne sais pas si elles sont \emph{vraies} ou \emph{bonnes}.
    \item \textbf{Analyse des risques (Asymétrie) :}
    \begin{itemize}
        \item Si je \textbf{suis} les coutumes/lois : Risque minime (erreur sur la vérité, mais sécurité sociale, paix, salut religieux \emph{au cas où}).
        \item Si je \textbf{m'écarte de la conformité} / innove : Risque maximal (prison, mort, damnation, chaos).
    \end{itemize}
    \item \textbf{Critère de la \cle{modération} :} Choisir les opinions \og les plus modérées \fg (les plus reçues par les plus sages) minimise l'erreur probable. C'est le \og pari le plus sûr \fg (anticipation du \og Pari \fg de Pascal).
\end{enumerate}
La \cle{distinction conceptuelle} clé est : \og Vérité théorique \fg (suspendue par le doute) $\neq$ \og Utilité pratique \fg (maintien de l'ordre). \fg

\textbf{Q3 : \og Suivre les coutumes de son pays est-il une marque de sagesse ou de lâcheté ? \fg Discutez.}
\textbf{Solution :}
\textbf{Thèse (Sagesse / Prudence - Descartes / Burke / Conservatisme) :} Oui, marque de \cle{sagesse pratique}. La coutume est le \og dépôt de la raison collective \fg (Burke), fruit de l'expérience séculaire. La respecter assure la \cle{stabilité sociale}, condition de toute vie commune. Au Cameroun, le respect des \cle{chefferies traditionnelles} et des \cle{coutumes successorales} apaise les conflits fonciers. Descartes a raison : le temps de la réflexion philosophique ne doit pas être celui de la subversion sociale.
\textbf{Antithèse (Lâcheté / Aliénation - Kant / Antigone / Décolonisation) :} Non, c'est de la \cle{lâcheté intellectuelle} (Kant : \og Paresse et couardise \fg causes de la minorité). Obéir \og par provision \fg ou par habitude, c'est refuser l'\cle{autonomie} (s'auto-légiférer). \cle{Antigone} désobéit à Créon (loi de la cité) au nom de la \cle{loi divine non écrite} (supérieure). Au Cameroun, la \cle{résistance nationaliste} (UPC, Ruben Um Nyobè) a \emph{brisé} la légalité coloniale (coutume/loi française) pour fonder une \cle{légitimité} nouvelle. La conformité fige les injustices (ex : excision, exclusion des femmes de l'héritage coutumier).
\textbf{Synthèse :} La distinction \cle{Descartes} est cruciale : \textbf{Provisoire vs Définitive}. Suivre la coutume \emph{provisoirement} (le temps d'élaborer une critique rationnelle) = Sagesse. Suivre la coutume \emph{définitivement} (par peur, paresse, intérêt) = Lâcheté / Hétéronomie. L'idéal : \og Obéir aux lois, mais travailler à les améliorer par la raison \fg (Kant, \emph{Qu'est-ce que les Lumières ?}). \fg
\end{exoresolu}

\courssub{Simulation Bac Blanc \& Méta-cognition (Auto-évaluation)}

\begin{methode}[Grille d'Auto-évaluation \og Objectif 16+/20 \fg — À utiliser en fin de copie (5 min)]
\textbf{Principe :} Devenir son propre correcteur. Cocher \textbf{OUI / NON / PARTIEL} pour chaque item. Si $<$ 80\% OUI $\rightarrow$ Revoir la copie (si temps) ou noter les points de progression pour la prochaine fois.
\begin{enumerate}
    \item \textbf{INTRODUCTION :} [ ] Amorce accrocheuse (pas cliché) ? [ ] Thèse du texte fidèlement résumée ? [ ] Problématique \textbf{explicite, interrogative, dialectique} ? [ ] Plan \textbf{annoncé clairement} (I, II, III) ?
    \item \textbf{DÉVELOPPEMENT (Structure) :} [ ] 3 parties équilibrées ? [ ] Titres de parties \textbf{thématiques} (pas \og I. Les arguments pour \fg) ? [ ] Transitions \textbf{logiques} entre I-II et II-III (pas \og Maintenant passons à... \fg) ?
    \item \textbf{DÉVELOPPEMENT (Contenu - Par paragraphe) :} [ ] 1 Idée force par § ? [ ] Argument \textbf{philosophique} (pas simple opinion) ? [ ] \textbf{Référence explicite} (Auteur, Concept, Œuvre) ? [ ] \textbf{Exemple concret} (Cameroun, Histoire, Science, Actualité) ? [ ] \textbf{Analyse} du lien Argument/Exemple/Problématique ?
    \item \textbf{QUESTIONS SUR TEXTE (si applicable) :} [ ] Q1 : Explication \emph{intratextuelle} pure ? [ ] Q2 : Analyse \emph{logique} (connecteurs, structure) ? [ ] Q3 : Mini-dissertation \textbf{structurée T/A/S} (pas liste d'avis) ?
    \item \textbf{CONCLUSION :} [ ] Bilan \textbf{synthétique} (réponse à la problématique) ? [ ] Ouverture \textbf{pertinente} (nouvelle question, horizon) ? [ ] Pas de \og En conclusion, j'ai montré que... \fg (redite) ?
    \item \textbf{FORME / LANGUE :} [ ] 0 faute d'orthographe/grammaire visible (relecture à l'envers) ? [ ] Vocabulaire \textbf{philosophique précis} (\cle{aliénation} $\neq$ \cle{exploitation}, \cle{liberté} $\neq$ \cle{licence}) ? [ ] Phrases \textbf{complètes}, syntaxe soignée ? [ ] Lisibilité (sauts de ligne, paragraphes indentés) ?
\end{enumerate}
\end{methode}

\begin{exoresolu}[Simulation Bac Blanc : Sujet Zéro — \og La Technique libère-t-elle l'homme ? \fg (4h) + Auto-correction guidée]
\textbf{Consigne :} Traiter le sujet en dissertation. Temps imparti : 4h. Utiliser la grille ci-dessus en fin d'épreuve.

\tcblower
\textbf{Simulation Chronométrée (Modèle de gestion du temps) :}
\begin{center}
\begin{tabularx}{\linewidth}{|c|X|c|}\hline
\textbf{Étape} & \textbf{Action} & \textbf{Durée} \\ \hline
0-30 min & \textbf{Brouillon intensif :} Analyse concepts, Thèses/Antithèses, Auteurs, Exemples, Problématique, Plan détaillé (niveau §). & 30 min \\ \hline
30-50 min & \textbf{Rédaction Introduction} (soignée, au net direct ou sur brouillon propre). & 20 min \\ \hline
50min-2h50 & \textbf{Rédaction Développement} (3 parties $\approx$ 40 min chacune). \textbf{Relire chaque partie} avant de passer à la suivante. & 2h \\ \hline
2h50-3h10 & \textbf{Rédaction Conclusion} (Bilan + Ouverture). & 20 min \\ \hline
3h10-3h40 & \textbf{Relecture GLOBALE :} Cohérence, Orthographe, Noms propres, Dates, Respect consignes. & 30 min \\ \hline
3h40-4h00 & \textbf{Auto-évaluation Grille} + Corrections mineures (ratures propres, flèches d'insertion). & 20 min \\ \hline
\end{tabularx}
\end{center}

\textbf{Exemple de Problématique validée :}
\og \textbf{Dans quelle mesure la technique, en maîtrisant la nature pour satisfaire nos besoins, ne finit-elle pas par asservir l'homme à ses propres créations, exigeant une réappropriation éthique et politique de la puissance technique ?} \fg

\textbf{Plan Type (3 Parties) :}
I. \textbf{La technique comme puissance libératrice : maîtrise de la nécessité, émancipation du travail.} (Prométhée, Marx \emph{travail}, Simondon, Ex: agriculture mécanisée Nord-Cameroun, santé).
II. \textbf{L'envers de la libération : l'aliénation technique, l'homme devenu rouage/ressource.} (Heidegger \textit{Gestell}, Ellul \emph{système technicien}, Marx \emph{fétichisme}, Ex: précarité ubérisation, surveillance algorithmique, déchets électroniques (Agbogbloshie au Ghana, décharges de Douala au Cameroun)).
III. \textbf{Vers une technique humanisée : responsabilité, subsidiarité, low-tech, gouvernance démocratique.} (Jonas \emph{responsabilité}, Illich \emph{outils conviviaux}, Droit du numérique, Ex: FabLabs Yaoundé, Open Source, Agroécologie).

\textbf{Auto-correction Type (Méta-cognition) :}
\og \textit{Auto-évaluation post-copie :} \textbf{Point fort :} Plan dialectique solide, exemples camerounais pertinents (Agroécologie, Mobile Money), vocabulaire précis (\textit{Gestell}, \emph{fétichisme}, \emph{convivialité}). \textbf{Point faible :} Transition I$\rightarrow$II un peu brutale (\og Cependant \fg sec). Ouverture trop générale (\og Il faut une éthique \fg). \textbf{Action corrective :} Travailler les \textbf{transitions dialectiques} (reprise du dernier mot de la partie précédente + annonce du problème nouveau). Préparer 3 \og ouvertures types \fg prêtes à l'emploi (Éthique, Politique, Esthétique/Anthropologique). \fg
\end{exoresolu}

\courssub{Gestion du Temps \& Hygiène de l'Épreuve (Bonus Stratégique)}

\begin{methode}[Le \og Chrono-Philosophe \fg : 4h = 240 min = Capital à investir]
\textbf{Stratégie \og Zéro Stress \fg :}
\begin{enumerate}
    \item \textbf{Choix du sujet (5 min max) :} Lire les 2 sujets (Dissertation / Texte). Choisir celui où \textbf{vous maîtrisez les concepts ET les auteurs}. Pas celui qui \og inspire \fg. Cochez : \og J'ai 3 auteurs + 2 exemples concrets + une problématique claire pour ce sujet \fg.
    \item \textbf{Brouillon = Investissement (30-40 min) :} Ne pas écrire la copie au brouillon. Faire : \textbf{Arbre conceptuel} $\rightarrow$ \textbf{Tableau T/A/S} $\rightarrow$ \textbf{Problématique} $\rightarrow$ \textbf{Plan détaillé (niveau § avec mots-clés auteurs/exemples)}.
    \item \textbf{Rédaction \og Flux \fg (2h30) :} Écrire \textbf{au net} directement si brouillon est propre. Sinon recopier vite. \textbf{Ne pas s'arrêter} sur un mot. Mettre \og XXX \fg et continuer. Revenir à la relecture.
    \item \textbf{Relecture \og 3 Passes \fg (30 min) :}
    \begin{enumerate}
        \item \textbf{Sens/Structure :} La problématique est-elle répondue ? Le plan suivi ? Les transitions logiques ?
        \item \textbf{Contenu/Précision :} Noms d'auteurs exacts ? Concepts définis ? Exemples datés/localisés ? Pas de contresens ?
        \item \textbf{Forme/Langue :} Orthographe (relecture \emph{à l'envers}, mot par mot), Ponctuation, Paragraphes, Lisibilité.
    \end{enumerate}
    \item \textbf{Kit de survie :} Stylos (bleu/noir) x3, Correcteur (blanc/roller), Montre (pas téléphone), Eau, Chocolat/Noix (glycémie), Mouchoirs.
\end{enumerate}
\end{methode}

\begin{exoresolu}[Cas pratique : \og Panne sèche \fg au milieu de la Partie II (Il reste 1h30, page blanche)]
\textbf{Situation :} Vous avez fait Intro + Partie I. Vous bloquez sur les arguments de la Partie II (Antithèse). Stress montant.

\tcblower
\textbf{Protocole \og Réanimation \fg (10 min max) :}
\begin{enumerate}
    \item \textbf{Respiration carrée (1 min) :} Inspire 4s, Retiens 4s, Expire 4s, Retiens 4s $\times$ 4. Cortisol $\downarrow$, Cortex préfrontal $\uparrow$.
    \item \textbf{Relire la Problématique \& la Thèse I (2 min) :} \og Ma thèse I : Technique libère. Problématique : Comment éviter l'asservissement ? \fg $\rightarrow$ La Partie II \textbf{DOIT} montrer l'asservissement.
    \item \textbf{Activer la \cle{Mémoire Sémantique} par mots-clés (3 min) :} Écrire en vrac sur brouillon : \og Aliénation, Hétéronomie, Déterminisme, Ellul, Heidegger, Jonas, Marx, Ubérisation, Écran, Algorithme, Déchet, Nucléaire, Armes, Perte sens, Bureaucratie \fg.
    \item \textbf{Construire 2 arguments "Bouée de sauvetage" (4 min) :}
    \begin{itemize}
        \item \textbf{Arg 1 (Classique) :} \og L'efficacité technique devient fin en soi (Ellul). L'homme s'adapte à la machine (cadences, horaires, notation). Ex: Ouvrier chaîne / Livreur Uber. \fg
        \item \textbf{Arg 2 (Contextuel Fort) :} \og La technique détruit le vivant (Jonas). Responsabilité sur l'avenir. Ex: Changement climatique (CO2 technique), Déchets plastiques océans, Nucléaire (Fukushima). Au Cameroun : Pollution plastique Douala, Déforestation bois d'œuvre. \fg
    \end{itemize}
    \item \textbf{Rédiger le § de transition + Arg 1 + Arg 2 (30 min) :} \og \textbf{Cependant}, cette libération par la technique porte en elle son ombre... \fg $\rightarrow$ Écrire sans viser la perfection, viser la \textbf{cohérence}.
    \item \textbf{Enchaîner direct sur Partie III (Synthèse) :} \og \textbf{Dès lors}, comment habiter la technique autrement ? \fg (Responsabilité, Low-tech, Démocratie technique).
\end{enumerate}
\textbf{Résultat :} Copie sauvée, structure respectée, note préservée (12-14/20 au lieu de 6/20 abandon).
\end{exoresolu}

\courssub{Boîte à Outils : Références \og Clés en main \fg pour le Bac Tech (Auteurs / Concepts / Exemples Cameroun)}

\begin{aretenir}[Répertoire Minimaliste \og High Yield \fg (Haut Rendement) — À mémoriser / Avoir en fiche bristol]
\textbf{Par Thème Programme (Terminale Tech) :}
\begin{center}
\begin{tabularx}{\linewidth}{|l|X|X|X|}\hline
\textbf{Thème} & \textbf{Auteurs \og Incontournables \fg (3 max)} & \textbf{Concepts \og Opératoires \fg (Définissables)} & \textbf{Exemples \og Cameroun / Concrets \fg (Cités nommément)} \\ \hline
\cle{Condition Humaine} & \textbf{Descartes} (Sujet/Cogito), \textbf{Sartre} (Liberté/Existence précède essence), \textbf{Levinas} (Visage/Autrui) & Cogito, \cle{Liberté radicale}, \cle{Mauvaise foi}, \cle{Responsabilité}, \cle{Condition}, \cle{Nature/Culture} & Rites d'initiation (Bamiléké, Bassa), \cle{Mort} / Deuil, Handicap \& inclusion, \cle{Identité} numérique. \\ \hline
\cle{Société / État / Droit} & \textbf{Hobbes/Locke/Rousseau} (Contrat), \textbf{Kant} (Droit/Cosmopolitique), \textbf{Arendt} (Politique/Totalitarisme), \textbf{Rawls} (Justice) & \cle{Contrat social}, \cle{Souveraineté}, \cle{Légitimité}, \cle{État de droit}, \cle{Séparation pouvoirs}, \cle{Justice distributive}, \cle{Bien commun} & \cle{Constitution 1996/2008}, \cle{Décentralisation} (Communes, Régions), \cle{Crise anglophone} (Justice/Éducation), \cle{Corruption} (CONAC), \cle{Élections}. \\ \hline
\cle{Vérité / Connaissance / Science} & \textbf{Descartes} (Doute/Preuve), \textbf{Kant} (Limites raison), \textbf{Popper} (Falsifiabilité), \textbf{Bachelard} (Obstacle épistémologique) & \cle{Doute methodique}, \cle{Critérium de vérité}, \cle{Paradigme}, \cle{Falsification}, \cle{Modèle}, \cle{Science/Opinion}, \cle{Technoscience} & \cle{Paludisme/Vaccin RTS,S}, \cle{Agriculture} (Semences améliorées vs OGM), \cle{Tradition orale} vs \cle{Science écrite}, \cle{IA} (ChatGPT / DeepSeek) fiabilité. \\ \hline
\cle{Morale / Politique / Éthique} & \textbf{Kant} (Devoir/Impératif catégorique), \textbf{Aristote} (Vertu/Bonheur/Phronèsis), \textbf{Jonas} (Responsabilité/Avenir), \textbf{Ricoeur} (Éthique/Estime de soi) & \cle{Devoir/Inclination}, \cle{Autonomie/Hétéronomie}, \cle{Vertu}, \cle{Fin/ Moyens}, \cle{Responsabilité}, \cle{Dignité}, \cle{Droits de l'homme} & \cle{Bioéthique} (Avortement, Fin de vie, Don d'organes), \cle{Éthique IA} (Biais algorithmiques), \cle{Corruption} (Devoir fonctionnaire), \cle{Écologie} (Devoir envers générations futures). \\ \hline
\cle{Esthétique / Art / Technique} & \textbf{Platon} (Mimésis/Interdit), \textbf{Aristote} (Catharsis), \textbf{Kant} (Jugement de goût/Désintéressement), \textbf{Heidegger} (Œuvre/Vérité), \textbf{Simondon} (Technique/Individuation) & \cle{Beau}, \cle{Art}, \cle{Technique}, \cle{Œuvre}, \cle{Goût}, \cle{Gestell}, \cle{Objet technique}, \cle{Esthétique du quotidien} & \cle{Art urbain} (Graffiti Yaoundé/Douala), \cle{Musique} (Makossa, Bikutsi, Afrobeats - technique studio), \cle{Architecture} (Cases traditionnelles vs Béton), \cle{Mode} (Tissus Ndop, Kaba ngondo). \\ \hline
\end{tabularx}
\end{center}
\textbf{Astuce :} Pour chaque concept, savoir donner : 1. Définition philosophique précise. 2. Auteur de référence. 3. Exemple concret (Cameroun si possible). 4. Contre-exemple / Limite.
\end{aretenir}

\begin{attention}[Les 5 Erreurs qui coûtent des points au Bac (Philosophie Terminale Tech)]
\begin{enumerate}
    \item \textbf{\og Le résumé de texte déguisé en dissertation \fg (Q3 / Dissertation) :} Raconter ce que dit l'auteur au lieu de \textbf{penser avec lui}. \textit{Correctif :} Pour chaque paragraphe, demander : \og Quel argument \textbf{MOI} j'avance pour répondre à la problématique ? \fg (Même si l'arg vient d'un auteur, c'est \textbf{votre} choix de l'utiliser).
    \item \textbf{\og La \og Pense-bête \fg d'auteurs sans articulation \fg :} Lister : \og Platon dit..., Aristote dit..., Kant dit... \fg sans lien logique. \textit{Correctif :} \og \textbf{Contre} la thèse de Platon (pour qui l'art est imitation dangereuse), \textbf{Aristote oppose} la catharsis purificatrice... \fg. Les auteurs \textbf{dialoguent} entre eux et avec votre problématique.
    \item \textbf{\og L'exemple \og décoratif \fg ou \og Wikipédia \fg \fg :} \og Par exemple, au Cameroun, il y a la corruption \fg. \textit{Correctif :} \og L'affaire \textbf{Opération Épervier} (2006-...) illustre la tension entre \cle{légalité} (procédures) et \cle{légitimité} (justice sociale perçue), montrant que l'État de droit peine à s'incarner sans \cle{vertu civique} (Rousseau). \fg $\rightarrow$ Exemple \textbf{daté, précis, analysé}.
    \item \textbf{\og La confusion \og Liberté / Licence \fg , \og Égalité / Identique \fg , \og Technique / Technologie \fg , \og Morale / Moralité \fg \fg :} Imprécision vocabulaire = perte de points \textbf{garantie}. \textit{Correctif :} \textbf{Définir} les concepts clés en intro ou au premier § où ils apparaissent. Utiliser le \textbf{vocabulaire technique} (\cle{Autonomie}, \cle{Hétéronomie}, \cle{Finalité}, \cle{Efficience}, \cle{Axiologie}).
    \item \textbf{\og La conclusion \og bac à sable \fg ou \og copier-coller intro \fg \fg :} \og En conclusion, j'ai montré que I, II, III. \fg Ou ouverture hors-sujet : \og Cela nous amène à la question de Dieu. \fg \textit{Correctif :} \textbf{Bilan dialectique} (La thèse I est vraie MAIS insuffisante car II, D'OÙ la nécessité de III). \textbf{Ouverture ciblée} : \og Cette analyse du rapport Technique/Liberté invite à repenser l'\textbf{éducation technique} au Cameroun : former des \textbf{ingénieurs-citoyens} (Simondon) capables de \textbf{désobéissance épistémique} face aux algorithmes opaques. \fg
\end{enumerate}
\end{attention}

\newpage

\coursec{Exercices}

\coursec{Consolidation des acquis de la méthodologie de l'exercice sur texte philosophique}

\courssub{Diagnostic \& Rappel des Exigences Officielles}

\begin{aretenir}[Attendus officiels de l'exercice sur texte (Probatoire / Bac Tech)]
L'exercice sur texte philosophique évalue quatre compétences majeures :
\begin{enumerate}
\item \textbf{Identifier} la thèse (l'idée directrice que l'auteur \emph{démontre}, pas seulement énoncée) et le thème.
\item \textbf{Problématiser} : transformer le thème en une question philosophique ouverte qui révèle la tension interne du texte.
\item \textbf{Expliquer} (et non paraphraser) : analyser les concepts, articuler les arguments, illustrer par des exemples pertinents (dont le contexte camerounais), citer le texte avec précision.
\item \textbf{Discuter critique} : analyser la cohérence interne (analyse interne), puis confronter la thèse à des objections argumentées (contre-exemples, distinctions conceptuelles, limites), pour aboutir à une position nuancée (accepter, nuancer, rejeter).
\item \textbf{Conclure} : bilan synthétique de la discussion + ouverture philosophique (autre auteur, autre contexte, question nouvelle).
\end{enumerate}
\end{aretenir}

\begin{methode}[Les 4 étapes canoniques de l'introduction]
\begin{enumerate}
\item \textbf{Accroche} : citation d'auteur, fait d'actualité, proverbe, définition du thème \emph{en lien direct avec le texte}.
\item \textbf{Présentation du texte} : Auteur (si connu), œuvre, date, thème général, \textbf{thèse} formulée en une phrase claire, structure logique du texte (repérage des grandes parties).
\item \textbf{Problématique} : Question ouverte (commence par \emph{Comment}, \emph{Pourquoi}, \emph{Dans quelle mesure}, \emph{Peut-on}) qui exprime l'enjeu philosophique du texte.
\item \textbf{Annonce du plan} : Deux parties équilibrées (ex: \og I. Analyse de la thèse... II. Discussion critique...\fg ou \og I. Explication... II. Discussion...\fg). Utiliser des connecteurs logiques (\og Dans un premier temps... Dans un second temps...\fg).
\end{enumerate}
\end{methode}

\begin{attention}[Pièges « tueurs » à éviter absolument]
\begin{itemize}
\item \textbf{Confondre thèse et thème :} Le thème est le sujet (ex: \og la liberté \fg) ; la thèse est la position défendue (ex: \og La liberté ne saurait être illimitée \fg).
\item \textbf{Paraphraser au lieu d'expliquer :} Répéter le texte avec d'autres mots $\neq$ en dégager le sens, la portée, les présupposés.
\item \textbf{Oublier les citations :} Une explication sans ancrage textuel (guillemets + numéro de ligne ou \og l. 3\fg) est une dissertation hors sujet.
\item \textbf{Problématique fermée :} \og L'auteur a-t-il raison ? \fg (Oui/Non) $\rightarrow$ Mauvais. \og Dans quelle mesure la thèse de l'auteur résiste-t-elle à l'épreuve des faits ? \fg $\rightarrow$ Bon.
\item \textbf{Plan invisible :} Pas d'annonce $\rightarrow$ note plafonnée.
\end{itemize}
\end{attention}

\begin{exercice}[QCM méthodologique — Diagnostic rapide]
Pour chaque question, coche la (ou les) bonne(s) réponse(s).
\begin{enumerate}
\item \textbf{La \cle{thèse} d'un texte philosophique est :}
\begin{enumerate}
\item le résumé de tout ce que dit l'auteur ;
\item \textbf{l'idée principale que l'auteur défend et cherche à démontrer} ;
\item l'opinion du candidat sur le sujet ;
\item la biographie de l'auteur.
\end{enumerate}

\item \textbf{Parmi les propositions suivantes, laquelle est une \cle{paraphrase} ?}
\begin{enumerate}
\item \textbf{Dire avec d'autres mots ce que l'auteur a déjà dit, sans l'expliquer} ;
\item Dégager le sens et la portée d'une formule de l'auteur ;
\item Discuter la pertinence de la thèse ;
\item Citer un autre philosophe.
\end{enumerate}

\item \textbf{L'argument d'autorité consiste à :}
\begin{enumerate}
\item frapper l'adversaire ;
\item \textbf{justifier une idée en invoquant le prestige d'un penseur ou d'une institution} ;
\item donner une définition précise ;
\item raconter un exemple personnel.
\end{enumerate}

\item \textbf{La meilleure problématique d'un commentaire de texte est :}
\begin{enumerate}
\item une question fermée appelant seulement oui ou non ;
\item \textbf{une question ouverte qui engage l'enjeu philosophique du texte} ;
\item une affirmation sans question ;
\item une question sur la vie de l'auteur.
\end{enumerate}

\item \textbf{Remets dans l'ordre les étapes de l'introduction :} (1) Accroche et présentation du texte ; (2) Formulation de la thèse ; (3) Problématique ; (4) Annonce du plan.
\end{enumerate}
\end{exercice}

\begin{exercice}[Vrai ou Faux justifié]
Réponds par \textbf{Vrai} ou \textbf{Faux} et justifie en deux ou trois lignes.
\begin{enumerate}
\item Dans un commentaire de texte, il est interdit de discuter (critiquer) la thèse de l'auteur.
\item Une copie sans aucune citation du texte étudié ne peut pas obtenir la moyenne.
\item La conclusion d'un commentaire doit rappeler la thèse, le bilan de la discussion et peut s'ouvrir sur une question nouvelle.
\item Expliquer un texte, c'est d'abord le recopier en changeant quelques mots.
\item Le plan du développement doit être visible et annoncé dès l'introduction.
\end{enumerate}
\end{exercice}

\begin{exercice}[Définitions précises — Mots-clés de la méthode]
Définis en trois lignes maximum chacune des notions suivantes, telles qu'elles s'emploient dans l'exercice sur texte philosophique :
\begin{enumerate}
\item la \cle{thèse} ;
\item la \cle{problématique} ;
\item l'\cle{explication} (par opposition à la paraphrase) ;
\item la \cle{discussion critique} ;
\item l'\cle{ouverture} de la conclusion.
\end{enumerate}
\end{exercice}

\begin{exercice}[Repérage rapide dans un texte]
On donne l'extrait suivant : \og L'école coloniale n'a pas seulement instruit : elle a appris à l'Africain à douter de lui-même. Or un peuple qui doute de ses propres lumières finit par importer jusqu'à ses problèmes. \fg
\begin{enumerate}
\item Relève dans cet extrait la phrase qui exprime la thèse.
\item Identifie le type d'argument utilisé (argument de fait, d'autorité, par analogie, \emph{ad hominem}).
\item Formule en une phrase le sens de l'expression \og importer jusqu'à ses problèmes \fg.
\end{enumerate}
\end{exercice}

\courssub{L'Atelier Problématisation : Du Sujet à la Question Directrice}

\begin{methode}[Construire une problématique en 3 temps]
\begin{enumerate}
\item \textbf{Repérer la tension} : Quel paradoxe, quelle opposition, quelle limite le texte met-il en lumière ? (ex: Universalité de la science vs Dépendance africaine).
\item \textbf{Formuler la question} : Transformer la tension en question ouverte (\og Comment... ? \fg, \og Dans quelle mesure... ? \fg, \og Pourquoi... ? \fg). La question doit \emph{contenir} les termes clés du texte.
\item \textbf{Vérifier la pertinence} : La réponse à cette question oblige-t-elle à parcourir \emph{tout} le texte ? Engage-t-elle une réflexion philosophique (pas seulement factuelle) ?
\end{enumerate}
\end{methode}

\begin{exemplebox}[Problématiques : Bonnes vs Mauvaises]
\begin{center}
\begin{tabularx}{\linewidth}{|l|X|X|}
\hline
\textbf{Thèse courte} & \textbf{Mauvaise problématique (fermée / hors-sujet)} & \textbf{Bonne problématique (ouverte / philosophique)} \\
\hline
L'école technique aliène l'ouvrier. & L'école technique aliène-t-elle l'ouvrier ? (Fermée) & \textbf{Dans quelle mesure} la spécialisation technique constitue-t-elle une aliénation pour le travailleur ? \\
\hline
Le mobile libère le paysan. & Le téléphone est-il utile au paysan ? (Factuelle) & \textbf{Comment} l'appropriation des technologies numériques transforme-t-elle l'autonomie du paysan camerounais ? \\
\hline
La tradition freine le progrès. & La tradition est-elle un frein ? (Fermée) & \textbf{Peut-on} penser l'articulation entre tradition et progrès sans les opposer radicalement ? \\
\hline
\end{tabularx}
\end{center}
\end{exemplebox}

\begin{exercice}[Plan flash — 5 minutes chrono par sujet]
Pour chacune des cinq thèses courtes suivantes, formule en cinq minutes : (a) une problématique sous forme de question ouverte ; (b) l'annonce d'un plan en deux parties équilibrées (titres des parties I et II).
\begin{enumerate}
\item \og L'école technique aliène l'ouvrier. \fg
\item \og Le téléphone mobile libère le paysan camerounais. \fg
\item \og La tradition est un frein au progrès. \fg
\item \og Obéir à l'État, c'est renoncer à sa liberté. \fg
\item \og La science rend la religion inutile. \fg
\end{enumerate}
\end{exercice}

\courssub{Rédaction Chirurgicale : Introduction, Développement, Conclusion}

\begin{methode}[L'Explication : La méthode « D.A.I.C. » (Définir, Analyser, Illustrer, Citer)]
Pour chaque concept ou argument clé du texte :
\begin{enumerate}
\item \textbf{Définir} le sens précis du terme dans le contexte de l'auteur (pas le dictionnaire courant).
\item \textbf{Analyser} la logique de l'argument : comment l'auteur passe-t-il d'une prémisse à la conclusion ? Quels connecteurs logiques (donc, car, or, mais) ?
\item \textbf{Illustrer} par un exemple concret, si possible camerounais (réalité vécue, fait d'actualité, référence culturelle), pour montrer la portée réelle du concept.
\item \textbf{Citer} le texte (guillemets + référence ligne) pour ancrer l'explication dans l'extrait imposé.
\end{enumerate}
\end{methode}

\begin{methode}[La Discussion Critique : Structure en 3 temps (A.O.S.)]
\begin{enumerate}
\item \textbf{Analyse interne (A)} : \og L'auteur a raison car...\fg. Montrer la cohérence, la force des arguments, la pertinence des exemples \emph{du texte}. C'est l'explication approfondie.
\item \textbf{Objection argumentée (O)} : \og Cependant / Toutefois...\fg. Apporter une limite, un contre-exemple, une distinction conceptuelle non vue par l'auteur. Mobiliser un autre auteur ou un fait réel (ex: médecine traditionnelle, tontine, chefferie au Cameroun).
\item \textbf{Synthèse / Nuance (S)} : \og Au total...\fg. La thèse est-elle à accepter, nuancer (valable sous conditions) ou rejeter ? Justifier la position finale.
\end{enumerate}
\end{methode}

\begin{exemplebox}[Paraphrase vs Explication : Le test du « Prix de la dignité »]
\textbf{Formule :} \og L'autonomie théorique est le prix de la dignité. \fg (Hountondji)

\textbf{Paraphrase (Zéro point en explication) :} \og Avoir sa propre théorie, c'est ce qui donne de la valeur à l'homme. \fg $\rightarrow$ Simple reformulation, aucun concept défini, aucune portée, aucun exemple.

\textbf{Explication (Niveau Bac) :}
\og L'\textbf{autonomie théorique} désigne ici la capacité du chercheur africain à poser lui-même ses objets d'étude, à construire ses propres concepts et méthodes, au lieu d'appliquer des protocoles venus du Nord. Le terme \textbf{prix} implique un coût, un effort : il faut investir dans la formation, les laboratoires, la publication locale. La \textbf{dignité} n'est pas un don, c'est une conquête : être \og interlocuteur \fg (pair) dans la communauté scientifique universelle, et non \og élève \fg (subalterne). \textbf{Exemple :} L'IRAD (Institut de Recherche Agricole pour le Développement) au Cameroun qui développe des variétés de manioc résistantes à la mosaïque adaptées aux sols locaux, au lieu d'attendre des semences importées. \fg
\end{exemplebox}

\begin{exercice}[De la paraphrase à l'explication]
Voici une formule d'un texte : \og L'autonomie théorique est le prix de la dignité. \fg
\begin{enumerate}
\item Écris une phrase qui ne fait que \emph{paraphraser} cette formule.
\item Écris ensuite un paragraphe de cinq à huit lignes qui l'\emph{explique} véritablement : sens des termes (\og autonomie théorique \fg, \og prix \fg, \og dignité \fg), enjeu, exemple concret tiré de la recherche scientifique au Cameroun (ex: IRAD, université, startup tech).
\item Indique ce qui distingue ton paragraphe de la paraphrase (utilise le vocabulaire : définition, articulation, illustration, citation).
\end{enumerate}
\end{exercice}

\begin{exercice}[Chasse aux erreurs — Copie notée 8/20]
Voici le début de la copie d'un élève sur un texte traitant de la liberté de penser :
\og Dans ce texte, l'auteur parle de la liberté. La liberté est très importante dans la vie. Moi je pense que la liberté c'est bien. L'auteur dit que la liberté c'est important et que les hommes doivent être libres. Depuis toujours les hommes aiment la liberté. En Afrique aussi il y a la liberté. \fg
\begin{enumerate}
\item À l'aide de la grille méthodologique (thèse, problématique, explication, citations, plan, conclusion), relève au moins cinq erreurs commises.
\item Annote chaque erreur en indiquant la règle méthodologique violée.
\item Réécris entièrement cette introduction en respectant les quatre étapes (accroche, présentation du texte et thèse, problématique, annonce du plan).
\end{enumerate}
\end{exercice}

\begin{exercice}[Construire une discussion critique — Cas camerounais]
Un texte soutient la thèse suivante : \og Le développement technique du Cameroun exige l'abandon des pratiques traditionnelles. \fg
\begin{enumerate}
\item Rédige un paragraphe d'\emph{analyse interne} : explique la thèse et montre sa cohérence (arguments possibles : modernité, efficacité, standardisation).
\item Rédige un paragraphe de \emph{discussion critique} : formule une objection argumentée, en mobilisant un exemple camerounais précis (médecine traditionnelle et pharmacopée, tontine et inclusion financière, chefferie traditionnelle et médiation sociale, agriculture paysanne et agroécologie).
\item Conclus en cinq lignes : la thèse doit-elle être acceptée, nuancée ou rejetée ? Justifie.
\end{enumerate}
\end{exercice}

\courssub{Maîtrise des Questions Types \& Pièges Spécifiques Bac Tech}

\begin{aretenir}[Typologie des questions au Bac Tech]
\begin{itemize}
\item \textbf{Présentez le texte / Dégagez la thèse} : Exigence de concision (1 phrase pour la thèse) + contextualisation auteur/œuvre.
\item \textbf{Expliquez telle formule / tel passage} : Exigence de la méthode D.A.I.C. (Définir, Analyser, Illustrer, Citer). Ne pas déborder sur le reste du texte.
\item \textbf{Montrez comment l'auteur articule... / Analysez les arguments} : Exigence de repérage structurel (connecteurs, enchainement logique) + évaluation de la force probante.
\item \textbf{Discutez / Commentez} : Exigence de la structure A.O.S. (Analyse interne, Objection, Synthèse). La discussion \emph{porte sur la thèse du texte}, pas sur le thème en général.
\item \textbf{Concluez} : Bilan + Ouverture.
\end{itemize}
\end{aretenir}

\begin{attention}[Pièges spécifiques « Bac Tech »]
\begin{itemize}
\item \textbf{Le hors-sujet technique :} Parler de son stage, de sa filière (Génie civil, Électrotechnique...) sans lien explicite avec la thèse du texte.
\item \textbf{L'exemple « décoratif » :} Citer \og le Cameroun \fg ou \og Yaoundé \fg sans que l'exemple ne fasse \emph{avancer} l'argumentation (ne prouve ni n'illustre le concept).
\item \textbf{La confusion « Commentaire / Dissertation » :} Si le sujet est \og Commentez ce texte \fg, on ne fait \textbf{pas} une dissertation sur le thème. Le texte est la \emph{seule} boussole.
\item \textbf{Gestion du temps :} 2h = 30 min lecture/analyse/brouillon + 1h15 rédaction + 15 min relecture. Ne pas rédiger l'introduction au brouillon en entier (mots-clés seulement).
\end{itemize}
\end{attention}

\begin{exercice}[Sujet type Bac Tech : Paulin Hountondji, \emph{La connaissance de soi}, 2002 (extrait adapté)]
\textbf{Texte :} \og La recherche scientifique en Afrique reste largement dépendante : les problèmes étudiés, les instruments utilisés, les revues où l'on publie viennent d'ailleurs. Le chercheur africain travaille souvent pour des laboratoires du Nord, sur des questions qui ne sont pas celles de son peuple. Pourtant, la science est universelle : une vérité mathématique ou agronomique ne porte pas de passeport. Mais pour participer à cette universalité, encore faut-il se poser soi-même ses propres questions. L'autonomie théorique est le prix de la dignité : tant que nous consommons les théories des autres sans produire les nôtres, nous restons des élèves et non des interlocuteurs. Cela ne signifie ni rejeter la science venue d'ailleurs, ni enfermer la pensée dans des traditions figées : il s'agit de faire de nos réalités un point de départ de la recherche, et non seulement un terrain d'expérimentation pour autrui. \fg

\textbf{Travail à faire :}
\begin{enumerate}
\item Présente le texte (auteur, thème) et dégage sa thèse en une phrase.
\item Explique la formule : \og L'autonomie théorique est le prix de la dignité \fg (sens des termes, enjeu, exemple camerounais).
\item Montre comment l'auteur articule l'universalité de la science et la particularité des contextes africains.
\item Discussion : L'Afrique peut-elle produire une science universelle sans renier ses traditions ? Appuie ta réponse sur au moins deux exemples camerounais (ex: pharmacopée traditionnelle / recherche pharmaceutique ; savoirs endogènes de gestion de l'eau / génie hydraulique moderne).
\item Rédige une conclusion complète avec ouverture.
\end{enumerate}
\end{exercice}

\begin{exercice}[Probatoire blanc : John Stuart Mill, \emph{De la liberté}, 1859 (extrait)]
\textbf{Texte :} \og Il y a une limite à l'interférence légitime de l'opinion collective dans l'indépendance individuelle ; trouver cette limite et la maintenir contre l'empiètement est aussi indispensable à une bonne condition des affaires humaines que la protection contre le despotisme politique. Car la tyrannie de la majorité, quand elle s'exerce par la pression de l'opinion et non par la loi, peut être plus oppressive que bien des gouvernements : elle enchaîne non seulement les actes, mais la pensée elle-même, et pénètre plus profondément dans les détails de la vie. L'individu qui pense autrement que la foule n'a pas seulement besoin d'être toléré : sa différence est un bienfait pour tous, car c'est en se confrontant aux opinions contraires que la vérité se fortifie. Une société qui étouffe l'excentricité se prive de ses propres progrès. \fg

\textbf{Question unique :} Commentez ce texte.

\textbf{Exigences :}
\begin{enumerate}
\item Une introduction complète (accroche, présentation du texte, thèse, problématique, annonce du plan).
\item Un développement en deux parties :
\begin{enumerate}
\item Analyse interne du texte (thèse, arguments, explication de \og la tyrannie de la majorité \fg).
\item Discussion critique (la majorité a-t-elle toujours tort ? Limites de la thèse de Mill : harmonie sociale, protection des vulnérables, relativisme).
\end{enumerate}
\item Une conclusion avec ouverture sur le contexte électoral et démocratique camerounais (pluralisme, liberté d'expression, rôle de l'opposition, loi sur la décentralisation).
\end{enumerate}
\end{exercice}

\courssub{Simulation Bac Blanc \& Méta-cognition (Auto-évaluation)}

\begin{experience}[Simulation Bac Blanc — Conditions réelles (2 heures)]
\textbf{Consignes :} Chronomètre enclenché. Aucun document, aucun téléphone. Une seule feuille de brouillon. Rédaction finale sur copie double.
\end{experience}

\begin{exercice}[Sujet complet — Simulation Bac Blanc]
\textbf{Durée : 2 heures.}

\textbf{Texte :} \og La technique moderne promet à l'homme la maîtrise du monde. Pourtant, celui qui croit dominer la nature par ses machines découvre souvent qu'il est devenu dépendant de ses propres outils : le paysan ne sait plus semer sans engrais importés, l'artisan ne sait plus réparer sans pièces venues d'ailleurs, l'élève ne sait plus chercher sans l'écran. La technique n'est pas mauvaise en elle-même : elle est un pouvoir, et tout pouvoir exige d'être gouverné. La question n'est donc pas d'accepter ou de refuser la technique, mais de savoir qui la maîtrise, et pour quelles fins. Un peuple qui consomme des techniques sans les comprendre ni les produire reste mineur ; un peuple qui s'en approprie les principes devient adulte. L'éducation technique, bien comprise, est ainsi une école de liberté : elle transforme l'usager en acteur. \fg

\textbf{Travail à faire :}
\begin{enumerate}
\item Dégage la thèse du texte et formule la problématique.
\item Explique la distinction fondamentale du texte entre \og usager \fg et \og acteur \fg de la technique.
\item Analyse les arguments de l'auteur et relève les exemples qu'il mobilise.
\item Discussion : L'enseignement technique au Cameroun réalise-t-il effectivement cette \og école de liberté \fg ? Argumente à partir de ton expérience d'élève de Terminale technique (stages, ateliers, projets, programmes, insertion).
\item Rédige une introduction et une conclusion conformes aux exigences méthodologiques.
\end{enumerate}
\end{exercice}

\begin{aretenir}[Grille d'Auto-évaluation — Méta-cognition (À remplir après correction)]
Note-toi sur \textbf{4 points} pour chacun des critères suivants (4 = Maîtrise parfaite / 0 = Absent) :
\begin{center}
\begin{tabularx}{\linewidth}{|X|c|c|c|c|c|}
\hline
\textbf{Critère} & \textbf{4} & \textbf{3} & \textbf{2} & \textbf{1} & \textbf{0} \\
\hline
Thèse dégagée (précise, en 1 phrase) & & & & & \\
\hline
Problématique pertinente (ouverte, enjeu philosophique) & & & & & \\
\hline
Explication (D.A.I.C. : concepts, arguments, exemples, citations) & & & & & \\
\hline
Citations du texte (précises, intégrées, analysées) & & & & & \\
\hline
Plan visible (annoncé, respecté, 2 parties équilibrées) & & & & & \\
\hline
Discussion argumentée (A.O.S. : interne, objection, synthèse) & & & & & \\
\hline
Conclusion avec ouverture (bilan + horizon nouveau) & & & & & \\
\hline
\end{tabularx}
\end{center}

\textbf{Calcul :} Total sur 28 points $\rightarrow$ Note sur 20 = (Total $\times$ 20) / 28.
\begin{itemize}
\item \textbf{16--20 / 20} : Prêt pour l'examen (Autonome).
\item \textbf{12--15 / 20} : Solide, mais des lacunes ciblées (Revoir fiches méthode).
\item \textbf{8--11 / 20} : Fragile (Revoir les bases : thèse, citation, plan).
\item \textbf{$<$ 8 / 20} : Insuffisant (Remédiation urgente : ateliers encadrés).
\end{itemize}
\textbf{Action :} Identifie tes 2 points les plus faibles et programme une séance de révision ciblée cette semaine.
\end{aretenir}

\newpage

\coursec{Corrigés des exercices}

\coursec{Consolidation des acquis de la méthodologie de l'exercice sur texte philosophique}

\courssub{Diagnostic \& Rappel des Exigences Officielles}

\begin{aretenir}[Les 4 étapes canoniques de l'introduction (exigibles au Bac)]
\begin{enumerate}
\item \textbf{Accroche et présentation du texte :} Auteur, œuvre, date, thème général.
\item \textbf{Formulation de la thèse :} Position précise de l'auteur (une phrase complète : sujet + verbe + attribut). Distincte du thème.
\item \textbf{Problématique :} Question ouverte (Comment ? Pourquoi ? Dans quelle mesure ?) issue de la thèse, qui commande le plan.
\item \textbf{Annonce du plan :} Les deux (ou trois) grandes parties du développement, équilibrées (explication puis discussion).
\end{enumerate}
\end{aretenir}

\begin{definition}[Thèse vs Thème]
Le \textbf{thème} est le sujet général abordé (ex. : « la liberté », « la technique »). La \textbf{thèse} est la position philosophique précise que l'auteur défend et démontre dans le texte (ex. : « La liberté de penser est une exigence fondamentale de la dignité humaine »).
\end{definition}

\begin{definition}[Paraphrase vs Explication (Méthode D.A.I.C.)]
La \textbf{paraphrase} reformule le texte avec d'autres mots sans en dégager le sens. L'\textbf{explication} suit la méthode \textbf{D.A.I.C.} : \textbf{D}éfinir les concepts clés, \textbf{A}rticuler les arguments (logique interne), \textbf{I}llustrer par des exemples concrets (camerounais si possible), \textbf{C}iter le texte (guillemets, références de lignes).
\end{definition}

\begin{exoresolu}[Exercice 1 — QCM méthodologique — Diagnostic rapide]
\begin{enumerate}
\item \textbf{Réponse : (b).} La thèse est l'idée principale que l'auteur défend et cherche à démontrer. (a) est faux : la thèse n'est pas un résumé exhaustif, c'est une position tranchée. (c) est faux : l'opinion du candidat n'intervient que dans la discussion. (d) relève de la présentation, non de la thèse.
\item \textbf{Réponse : (a).} La paraphrase consiste à redire avec d'autres mots ce que l'auteur a déjà dit, sans dégager le sens ni la portée. (b) définit l'explication, (c) la discussion, (d) une mobilisation de référence.
\item \textbf{Réponse : (b).} L'argument d'autorité justifie une thèse en invoquant le prestige d'un penseur, d'une institution ou d'une tradition (\og Aristote l'a dit, donc c'est vrai \fg). C'est un argument recevable comme appui, mais insuffisant seul en philosophie, où seule la démonstration fait foi.
\item \textbf{Réponse : (b).} Une bonne problématique est une question \emph{ouverte} (Comment ? Pourquoi ? Dans quelle mesure ?) qui engage l'enjeu philosophique du texte et dont la réponse exige de parcourir tout l'extrait. Une question fermée (oui/non) ferme la réflexion au lieu de l'ouvrir.
\item \textbf{Ordre correct : (1) $\rightarrow$ (2) $\rightarrow$ (3) $\rightarrow$ (4).} Accroche et présentation du texte, puis formulation de la thèse, puis problématique, puis annonce du plan. Inverser la problématique et la thèse est une faute logique : on ne peut problématiser que ce qu'on a d'abord identifié.
\end{enumerate}
\textbf{Points de vigilance :} au Bac, le correcteur vérifie dès l'introduction que la thèse est formulée \emph{avant} la problématique et que celle-ci est une véritable question.
\end{exoresolu}

\begin{exoresolu}[Exercice 2 — Vrai ou Faux justifié]
\begin{enumerate}
\item \textbf{Faux.} La discussion critique est une exigence officielle de l'exercice sur texte : après l'explication (analyse interne), le candidat doit confronter la thèse à des objections argumentées et prendre position de manière nuancée. Un commentaire sans discussion est plafonné.
\item \textbf{Vrai.} L'ancrage textuel est la condition de validité du commentaire : sans citations précises (guillemets, références de lignes), la copie devient une dissertation hors sujet sur le thème, sanctionnée en dessous de la moyenne.
\item \textbf{Vrai.} La conclusion canonique comporte trois temps : rappel de la thèse, bilan de la discussion (position justifiée), ouverture sur une question nouvelle, un autre auteur ou un autre contexte.
\item \textbf{Faux.} Recopier en changeant quelques mots, c'est la paraphrase, faute méthodologique majeure. Expliquer, c'est définir les concepts, articuler les arguments, dégager la portée et illustrer (méthode D.A.I.C.).
\item \textbf{Vrai.} L'annonce du plan est la quatrième étape obligatoire de l'introduction. Un plan non annoncé est un \og plan invisible \fg, qui plafonne la note même si le développement est bon.
\end{enumerate}
\end{exoresolu}

\begin{exoresolu}[Exercice 3 — Définitions précises — Mots-clés de la méthode]
\begin{enumerate}
\item \textbf{La thèse :} position philosophique précise que l'auteur défend et s'efforce de démontrer dans le texte. Elle se formule en une phrase complète (sujet + verbe + attribut) et se distingue du thème, qui n'est que le sujet général abordé.
\item \textbf{La problématique :} question ouverte (Comment ? Pourquoi ? Dans quelle mesure ?) qui traduit la tension interne du texte et dont la réponse commande l'organisation de tout le devoir. Elle doit contenir les termes clés du texte.
\item \textbf{L'explication :} analyse qui dégage le sens, la portée et les présupposés d'un passage : définition des concepts, articulation des arguments, illustration par des exemples, citation du texte. Elle s'oppose à la paraphrase, simple reformulation sans valeur ajoutée.
\item \textbf{La discussion critique :} confrontation argumentée de la thèse : analyse interne (cohérence, force des arguments), objection (contre-exemple, distinction, limite), synthèse (accepter, nuancer ou rejeter la thèse, position justifiée).
\item \textbf{L'ouverture :} dernier mouvement de la conclusion : après le bilan, on élargit la réflexion vers une question nouvelle, un autre philosophe, un autre domaine ou un autre contexte, sans ouvrir un nouveau développement.
\end{enumerate}
\end{exoresolu}

\courssub{L'Atelier Problématisation : Du Sujet à la Question Directrice}

\begin{methode}[Construire une problématique en 3 temps]
\begin{enumerate}
\item \textbf{Identifier la thèse} (position de l'auteur) et les concepts clés.
\item \textbf{Repérer la tension} : paradoxe, limite, opposition interne au texte (ex. : universalité vs particularité, liberté vs contrainte, tradition vs modernité).
\item \textbf{Formuler la question} : ouverte, contenant les termes de la thèse, dont la réponse exige d'expliquer \emph{et} de discuter.
\end{enumerate}
\end{methode}

\begin{exoresolu}[Exercice 4 — Repérage rapide dans un texte (Extrait type : École coloniale)]
\begin{enumerate}
\item \textbf{Phrase-thèse :} \og L'école coloniale n'a pas seulement instruit : elle a appris à l'Africain à douter de lui-même. \fg C'est la position que l'auteur défend ; la seconde phrase en tire la conséquence.
\item \textbf{Type d'argument :} argument \emph{de fait} (ou par les conséquences) : l'auteur s'appuie sur un constat historique (l'effet réel de l'école coloniale) et en déduit une conséquence observable (la dépendance intellectuelle). Ce n'est ni un argument d'autorité (aucun penseur invoqué), ni une analogie, ni un \emph{ad hominem} (aucune personne attaquée).
\item \textbf{Sens de \og importer jusqu'à ses problèmes \fg :} un peuple qui a perdu confiance en sa propre raison ne se contente pas d'importer des solutions étrangères : il finit par adopter aussi les questions, les catégories et les problèmes des autres, au lieu de poser les siens à partir de ses réalités. L'expression radicalise l'idée de dépendance : elle devient intellectuelle et non seulement matérielle.
\end{enumerate}
\textbf{Point de vigilance :} bien distinguer la phrase-thèse (position défendue) de la phrase-conséquence (ce que la thèse entraîne).
\end{exoresolu}

\begin{exoresolu}[Exercice 5 — Plan flash — 5 minutes chrono par sujet]
\begin{enumerate}
\item \textbf{Thèse :} \og L'école technique aliène l'ouvrier. \fg\\
(a) Problématique : Dans quelle mesure la spécialisation technique constitue-t-elle une aliénation pour le travailleur ?\\
(b) Plan : I. L'école technique comme facteur d'aliénation (explication : spécialisation, fragmentation des tâches). II. L'école technique comme voie d'émancipation (discussion : qualification, autonomie, insertion — exemple des lycées techniques camerounais).
\item \textbf{Thèse :} \og Le téléphone mobile libère le paysan camerounais. \fg\\
(a) Problématique : Comment l'appropriation des technologies numériques transforme-t-elle l'autonomie du paysan camerounais ?\\
(b) Plan : I. Le mobile comme instrument de libération (information sur les prix, mobile money, accès aux marchés). II. Les limites de cette libération (dépendance au réseau, coûts, fracture numérique, nouvelle forme de dépendance).
\item \textbf{Thèse :} \og La tradition est un frein au progrès. \fg\\
(a) Problématique : Peut-on penser l'articulation entre tradition et progrès sans les opposer radicalement ?\\
(b) Plan : I. La tradition comme obstacle (conservatisme, résistance au changement). II. La tradition comme ressource du progrès (pharmacopée, tontine, savoirs endogènes intégrés au développement).
\item \textbf{Thèse :} \og Obéir à l'État, c'est renoncer à sa liberté. \fg\\
(a) Problématique : L'obéissance aux lois de l'État est-elle compatible avec la liberté du citoyen ?\\
(b) Plan : I. L'obéissance comme renonciation (contrainte, obéissance passive). II. L'obéissance comme condition de la liberté (la loi protège la liberté de tous ; distinction liberté/libertinage).
\item \textbf{Thèse :} \og La science rend la religion inutile. \fg\\
(a) Problématique : Le progrès de la connaissance scientifique supprime-t-il la fonction propre de la religion ?\\
(b) Plan : I. La science comme substitut à la religion (explication rationnelle du monde). II. L'irréductibilité du religieux (questions de sens, de morale, de finitude, que la science ne tranche pas).
\end{enumerate}
\textbf{Points de vigilance :} la problématique doit reprendre les termes de la thèse ; les deux parties doivent être équilibrées et couvrir ensemble l'explication et la discussion.
\end{exoresolu}

\courssub{Rédaction Chirurgicale : Introduction, Développement, Conclusion}

\begin{aretenir}[Structure type du développement (Bac Tech)]
\begin{itemize}
\item \textbf{Première partie (Explication / Analyse interne) :} Suivre le fil du texte, paragraphe par paragraphe ou argument par argument. Appliquer D.A.I.C. à chaque étape. Citer systématiquement.
\item \textbf{Deuxième partie (Discussion critique / A.O.S.) :} \textbf{A}nalyse interne (cohérence), \textbf{O}bjection(s) argumentée(s) (contre-exemple, distinction, limite), \textbf{S}ynthèse (position nuancée et justifiée). Mobiliser des exemples camerounais précis.
\end{itemize}
\end{aretenir}

\begin{exoresolu}[Exercice 6 — De la paraphrase à l'explication (Extrait : Hountondji, \emph{Autonomie théorique})]
\begin{enumerate}
\item \textbf{Paraphrase (faute) :} \og Avoir sa propre théorie, c'est ce qui donne de la valeur à l'homme. \fg (Simple reformulation : mêmes idées, autres mots, aucun concept défini, aucun exemple.)
\item \textbf{Explication (corrigé type - D.A.I.C.) :}
\textbf{Définir :} L'\textbf{autonomie théorique} désigne la capacité du chercheur africain à poser lui-même ses objets d'étude, à construire ses propres concepts et méthodes, au lieu d'appliquer des protocoles élaborés ailleurs. Le terme \textbf{prix} implique un coût : cette autonomie s'achète par l'effort — formation, laboratoires, publications locales. La \textbf{dignité} n'est donc pas un don mais une conquête : être \og interlocuteur \fg, c'est-à-dire pair reconnu dans la communauté scientifique universelle, et non \og élève \fg subalterne qui consomme les théories d'autrui.
\textbf{Analyser :} La formule établit un rapport de condition : pas de dignité scientifique sans production endogène de savoir — \og tant que nous consommons les théories des autres sans produire les nôtres, nous restons des élèves \fg. L'enjeu est politique autant que scientifique : produire ses propres savoirs, c'est cesser d'être un simple terrain d'expérimentation.
\textbf{Illustrer (Cameroun) :} L'IRAD (Institut de Recherche Agricole pour le Développement) développe des variétés de manioc résistantes à la mosaïque, adaptées aux sols locaux, au lieu d'attendre des semences importées ; de même, des startups technologiques de Yaoundé et Douala conçoivent des applications de mobile money adaptées aux réalités du pays.
\textbf{Citer :} Ancrer chaque temps par des guillemets sur l'extrait (\og autonomie théorique \fg, \og prix \fg, \og dignité \fg, \og interlocuteur \fg).
\item \textbf{Ce qui distingue l'explication de la paraphrase :} la paraphrase se contente de reformuler ; l'explication \emph{définit} les trois termes clés, \emph{articule} la logique de la formule (coût $\rightarrow$ conquête $\rightarrow$ reconnaissance), \emph{illustre} par des exemples concrets camerounais, et s'ancre par des \emph{citations}. Elle ajoute du sens au lieu de répéter.
\end{enumerate}
\end{exoresolu}

\begin{exoresolu}[Exercice 7 — Chasse aux erreurs — Copie notée 8/20]
\begin{enumerate}
\item \textbf{Erreurs relevées (au moins cinq) :}
\begin{enumerate}
\item \og l'auteur parle de la liberté \fg : confusion entre \textbf{thème} et \textbf{thèse} — la position défendue n'est jamais formulée.
\item Aucune \textbf{présentation du texte} : ni auteur, ni œuvre, ni contexte.
\item \og Moi je pense que la liberté c'est bien \fg : opinion personnelle \textbf{non argumentée}, hors de propos dans l'explication.
\item \og L'auteur dit que la liberté c'est important \fg : \textbf{paraphrase} vague, sans citation précise ni analyse.
\item Aucune \textbf{problématique} : pas de question directrice.
\item Aucune \textbf{annonce de plan} : le correcteur ne sait pas où va la copie.
\item \og En Afrique aussi il y a la liberté \fg : exemple \textbf{décoratif}, sans lien argumentatif avec le texte.
\end{enumerate}
\item \textbf{Règles violées :} (a) distinction thème/thèse ; (b) étape 2 de l'introduction ; (c) interdiction de l'opinion non fondée ; (d) exigence D.A.I.C. et citations ; (e) étape 3 de l'introduction ; (f) étape 4 de l'introduction ; (g) pertinence de l'illustration.
\item \textbf{Introduction réécrite (corrigé type) :}\\
\og La liberté de penser, disait déjà les Lumières, est le premier des biens, car sans elle aucun autre ne peut être défendu. Le texte soumis à notre étude, extrait d'une réflexion sur la liberté, soutient la thèse suivante : la liberté de penser est une exigence fondamentale de la dignité humaine, que nulle autorité ne saurait légitimement étouffer. Mais cette thèse soulève une difficulté : la liberté de penser et de s'exprimer peut-elle être absolue, ou rencontre-t-elle des limites légitimes ? Dans quelle mesure la défense de la liberté de penser est-elle compatible avec les exigences de la vie en société ? Pour répondre à cette question, nous analyserons d'abord la thèse de l'auteur et ses arguments, puis nous la soumettrons à la discussion critique. \fg
\end{enumerate}
\textbf{Barème indicatif :} 5 erreurs relevées et annotées : 5 pts ; introduction réécrite conforme aux 4 étapes : 5 pts.
\end{exoresolu}

\begin{exoresolu}[Exercice 8 — Construire une discussion critique — Cas camerounais (Thèse : \og Développement technique = abandon traditions \fg)]
\begin{enumerate}
\item \textbf{Analyse interne (corrigé type) :} La thèse affirme que le développement technique du Cameroun exige l'abandon des pratiques traditionnelles. Sa cohérence repose sur trois arguments. D'abord, la \emph{modernité} technique suppose des procédures standardisées, mesurables et reproductibles, que les pratiques traditionnelles, transmises oralement et variables, ne garantiraient pas. Ensuite, l'\emph{efficacité} : la mécanisation agricole, par exemple, produit davantage que la houe traditionnelle. Enfin, la \emph{standardisation} : pour intégrer les marchés et les normes internationales, un pays doit adopter des techniques uniformes. La tradition apparaît alors comme un conservatisme qui retarde l'industrialisation.
\item \textbf{Discussion critique (corrigé type - Structure A.O.S.) :}
\textbf{Objection 1 (Ressource médicale) :} La \emph{médecine traditionnelle} camerounaise n'est pas un obstacle mais une ressource — la recherche pharmaceutique s'appuie sur la pharmacopée endogène pour identifier des principes actifs, et l'OMS reconnaît la médecine traditionnelle comme complément des systèmes de santé.
\textbf{Objection 2 (Ressource financière) :} La \emph{tontine}, pratique d'épargne rotative traditionnelle, assure l'inclusion financière de millions de Camerounais exclus des banques ; loin d'être abandonnée, elle inspire aujourd'hui des applications de mobile money.
\textbf{Objection 3 (Ressource sociale) :} La \emph{chefferie traditionnelle} joue un rôle de médiation sociale et de gestion des conflits fonciers que l'administration moderne ne remplace pas.
\textbf{Synthèse :} L'opposition tradition/modernité est conceptuellement contestable : ce qui freine le développement n'est pas la tradition en tant que telle, mais le refus de la transformer et de l'articuler aux techniques nouvelles.
\item \textbf{Conclusion (corrigé type) :} Au total, la thèse doit être \emph{nuancée} plutôt qu'acceptée. Elle a raison de rappeler que le développement exige des techniques efficaces et standardisées ; mais elle a tort de poser l'\emph{abandon} des traditions comme condition nécessaire. Les exemples de la pharmacopée, de la tontine et de la chefferie montrent que les pratiques traditionnelles peuvent être des ressources du développement, à condition d'être critiquées, sélectionnées et hybridées avec les techniques modernes. Le vrai enjeu n'est donc pas d'abandonner la tradition, mais de la penser.
\end{enumerate}
\textbf{Points de vigilance :} l'objection doit viser la \emph{thèse} (et non le thème en général) ; chaque exemple camerounais doit faire avancer l'argumentation ; la position finale doit être justifiée, non simplement affirmée.
\end{exoresolu}

\courssub{Maîtrise des Questions Types \& Pièges Spécifiques Bac Tech}

\begin{attention}[Pièges fréquents à éviter]
\begin{itemize}
\item \textbf{Confondre thème et thèse} (0 pt à la thèse).
\item \textbf{Problématique fermée} (Oui/Non) ou hors sujet (ne reprenant pas les termes du texte).
\item \textbf{Paraphrase} à la place de l'explication (D.A.I.C. absent).
\item \textbf{Citations manquantes} ou approximatives (pas de guillemets, pas de lignes).
\item \textbf{Discussion générale} sur le thème au lieu de discussion \emph{de la thèse du texte}.
\item \textbf{Exemples décoratifs} (cités pour faire joli, sans lien argumentatif).
\item \textbf{Conclusion sans ouverture} ou ouverture qui lance un nouveau développement.
\item \textbf{Gestion du temps} : brouillon trop long, rédaction bâclée, pas de relecture.
\end{itemize}
\end{attention}

\begin{exoresolu}[Exercice 9 — Sujet type Bac Tech : Paulin Hountondji, \emph{La connaissance de soi} (2002)]
\textbf{Corrigé détaillé avec barème indicatif (sur 20).}

\begin{enumerate}
\item \textbf{Présentation et thèse (3 pts).} Le texte est un extrait adapté de \emph{La connaissance de soi} (2002) du philosophe béninois Paulin Hountondji ; son thème est la recherche scientifique en Afrique. \textbf{Thèse (en une phrase) :} l'Afrique ne deviendra un interlocuteur digne dans la science universelle qu'en conquérant son autonomie théorique, c'est-à-dire en se posant elle-même ses propres questions à partir de ses réalités. \emph{Point de vigilance :} la thèse doit être formulée en une seule phrase complète, distincte du thème.
\item \textbf{Explication de la formule (5 pts — méthode D.A.I.C.).} \textbf{Définir :} l'\og autonomie théorique \fg est la capacité à produire ses propres problèmes, concepts et méthodes ; le \og prix \fg signifie le coût de cette conquête (formation, institutions, moyens) ; la \og dignité \fg désigne le statut d'interlocuteur (pair) et non d'élève (subalterne). \textbf{Analyser :} la formule établit un rapport de condition : pas de dignité scientifique sans production endogène de savoir — \og tant que nous consommons les théories des autres sans produire les nôtres, nous restons des élèves \fg. \textbf{Illustrer :} l'IRAD qui développe des variétés de manioc résistantes adaptées aux sols camerounais ; les laboratoires universitaires camerounais travaillant sur les plantes de la pharmacopée locale. \textbf{Citer :} ancrer chaque temps de l'explication par des guillemets sur l'extrait.
\item \textbf{Articulation universalité/particularité (4 pts).} L'auteur procède en trois temps : (i) il affirme l'universalité de la science — \og une vérité mathématique ou agronomique ne porte pas de passeport \fg ; (ii) il montre que l'universalité n'exclut pas la particularité des points de départ : \og pour participer à cette universalité, encore faut-il se poser soi-même ses propres questions \fg ; (iii) il écarte les deux excès symétriques : ni rejet de la science venue d'ailleurs, ni enfermement dans des traditions figées. La vérité scientifique est universelle dans sa \emph{validité}, mais particulière dans ses \emph{questions de départ}. \emph{Point de vigilance :} repérer les connecteurs (\og Pourtant \fg, \og Mais \fg, \og Cela ne signifie ni... ni \fg) qui structurent l'argumentation.
\item \textbf{Discussion (5 pts — structure A.O.S.).} \textbf{Analyse interne :} Hountondji a raison : la dépendance des problèmes, instruments et revues est documentée ; produire ses questions est bien la condition de la reconnaissance. \textbf{Objection :} toutefois, l'autonomie théorique peut-elle être totale dans une science mondialisée où la collaboration est la norme ? Le risque est l'isolement. De plus, la mobilisation des traditions exige un tri critique : tout savoir endogène n'est pas scientifiquement valide. \textbf{Exemples camerounais attendus (au moins deux) :} la pharmacopée traditionnelle soumise à la recherche pharmaceutique (validation des principes actifs) ; les savoirs endogènes de gestion de l'eau (puits, calendriers hydrologiques villageois) articulés au génie hydraulique moderne. \textbf{Synthèse :} l'Afrique peut produire une science universelle sans renier ses traditions, à condition de les traiter comme des \emph{points de départ} soumis à l'épreuve de la méthode, et non comme des vérités intouchables ni comme des obstacles à abolir.
\item \textbf{Conclusion avec ouverture (3 pts).} Bilan : la thèse de Hountondji, nuancée par la discussion, demeure valable — la dignité scientifique de l'Afrique passe par la production endogène de problèmes, sans repli identitaire ni soumission. Ouverture possible : cette exigence d'autonomie théorique vaut-elle aussi pour d'autres domaines que la science, par exemple l'économie ou l'éducation technique au Cameroun ?
\end{enumerate}
\textbf{Barème récapitulatif :} 3 + 5 + 4 + 5 + 3 = 20. \textbf{Points de vigilance généraux :} citations obligatoires dans les questions 2 et 3 ; la discussion doit porter sur la thèse du texte ; les exemples doivent être précis et argumentés, non décoratifs.
\end{exoresolu}

\begin{exoresolu}[Exercice 10 — Probatoire blanc : John Stuart Mill, \emph{De la liberté} (1859)]
\textbf{Corrigé type complet (barème indicatif sur 20 : introduction 4 pts ; analyse interne 6 pts ; discussion 6 pts ; conclusion 4 pts).}

\textbf{1. Introduction (corrigé type).}\\
\og La démocratie, écrivait Tocqueville, porte en elle un danger paradoxal : celui où la volonté du plus grand nombre devient une servitude pour chacun. C'est ce danger qu'analyse John Stuart Mill dans un extrait de \emph{De la liberté} (1859), ouvrage fondamental de la pensée libérale. Le philosophe anglais y soutient la thèse suivante : il existe une limite à l'interférence légitime de l'opinion collective dans l'indépendance individuelle, et la défense de la différence individuelle est une condition du progrès de tous. Mais cette défense de l'excentricité soulève une difficulté : la majorité a-t-elle toujours tort contre l'individu, et la liberté d'opinion peut-elle être sans limites ? Dans quelle mesure la protection de l'indépendance individuelle est-elle compatible avec les exigences de la vie collective ? Pour répondre à cette question, nous analyserons d'abord la thèse de Mill et ses arguments, puis nous la soumettrons à la discussion critique. \fg

\textbf{2. Développement — Première partie : analyse interne (corrigé type).}\\
La thèse de Mill s'articule en trois arguments. \textbf{Premier argument :} l'existence d'une limite à l'interférence collective — \og trouver cette limite et la maintenir contre l'empiètement \fg est aussi indispensable que la protection contre le despotisme politique : Mill élargit la notion de tyrannie au-delà du seul pouvoir d'État. \textbf{Deuxième argument :} la \og tyrannie de la majorité \fg. Il faut définir cette formule : lorsque la majorité impose ses opinions non par la loi mais par la pression sociale (moquerie, exclusion, conformisme), elle devient \og plus oppressive que bien des gouvernements \fg, car le despotisme politique n'enchaîne que les actes, tandis que la tyrannie de l'opinion \og pénètre plus profondément dans les détails de la vie \fg et atteint \og la pensée elle-même \fg. L'oppression sociale est plus redoutable que l'oppression juridique parce qu'elle est intériorisée. \textbf{Troisième argument :} l'utilité de la différence. L'individu qui pense autrement \og n'a pas seulement besoin d'être toléré : sa différence est un bienfait pour tous \fg, car \og c'est en se confrontant aux opinions contraires que la vérité se fortifie \fg. La logique est celle du libre examen : une vérité jamais contestée devient un préjugé mort ; l'excentricité est le moteur du progrès — \og une société qui étouffe l'excentricité se prive de ses propres progrès \fg. La cohérence interne du texte est donc forte : de la limite (quoi protéger) à la tyrannie (contre quoi) puis à l'utilité (pourquoi).

\textbf{3. Développement — Deuxième partie : discussion critique (corrigé type).}\\
\textbf{Objection 1 — La majorité n'a pas toujours tort.} Mill semble supposer que l'opinion commune est suspecte par nature ; or le consensus majoritaire peut exprimer des valeurs protectrices : l'interdiction de l'incitation à la haine, par exemple, limite la liberté d'expression pour protéger les plus vulnérables. Toute \og excentricité \fg n'est pas un bienfait : la différence n'est féconde que si elle s'engage dans la confrontation rationnelle, non si elle propage le mensonge ou la violence.\\
\textbf{Objection 2 — Le risque du relativisme.} Si toutes les opinions doivent s'affronter librement, comment éviter que les erreurs les plus séduisantes ne l'emportent dans l'opinion ? Mill répond que la vérité \og se fortifie \fg dans la confrontation, mais cela suppose des citoyens éduqués et un espace public rationnel — conditions qui ne sont pas toujours réunies.\\
\textbf{Objection 3 — L'harmonie sociale.} Une société a aussi besoin d'un minimum de convictions partagées pour exister ; la protection absolue de l'excentricité peut fragiliser le lien commun.\\
\textbf{Synthèse :} la thèse de Mill doit être \emph{nuancée} : elle a raison de défendre l'indépendance individuelle contre la tyrannie de l'opinion, car sans dissidence aucun progrès n'est possible ; mais la liberté d'expression rencontre des limites légitimes lorsqu'elle nuit à autrui — ce que Mill lui-même reconnaît d'ailleurs par son principe de non-nuisance.

\textbf{4. Conclusion avec ouverture (corrigé type).}\\
Au terme de notre analyse, Mill défend l'existence d'une limite infranchissable à l'interférence de l'opinion collective dans l'indépendance individuelle, et fait de la différence une condition du progrès commun. Notre discussion a montré que cette thèse, fondatrice, doit être complétée par la protection des vulnérables et l'exigence d'un débat public rationnel. Cette réflexion éclaire directement le contexte démocratique camerounais : le pluralisme politique, la liberté d'expression de la presse et le rôle de l'opposition dans le jeu électoral posent concrètement la question millienne de la majorité et de la dissidence ; de même, la décentralisation, en rapprochant la décision des citoyens, fait de la confrontation des opinions locales un enjeu quotidien. Reste ouverte une question : à l'ère des réseaux sociaux, où la pression de l'opinion devient instantanée et mondiale, la tyrannie de la majorité n'a-t-elle pas pris une forme nouvelle que Mill n'avait pas prévue ?
\end{exoresolu}

\courssub{Simulation Bac Blanc \& Méta-cognition (Auto-évaluation)}

\begin{methode}[Gestion du temps type (Épreuve de 2h ou 3h selon série)]
\begin{enumerate}
\item \textbf{Lecture active \& Analyse (30 min) :} 3 lectures (découverte, repérage thèse/arguments, sélection citations). Remplir la fiche de brouillon structurée (Thèse, Problématique, Plan, Citations clés, Exemples camerounais).
\item \textbf{Rédaction soignée (1h15 à 1h30) :} Introduction (15 min), Développement Partie 1 (25 min), Développement Partie 2 (25 min), Conclusion (10 min). Écrire lisiblement, sauter des lignes, numéroter les pages.
\item \textbf{Relecture active (15 min) :} Vérifier : thèse formulée ? Problématique ouverte ? Plan annoncé et respecté ? Citations présentes (guillemets + lignes) ? D.A.I.C. respecté ? Discussion A.O.S. ? Exemples camerounais argumentés ? Conclusion en 3 temps ? Orthographe / syntaxe.
\end{enumerate}
\end{methode}

\begin{exoresolu}[Exercice 11 — Sujet complet — Simulation Bac Blanc (Texte anonyme : Technique, Usager, Acteur, Éducation technique)]
\textbf{Corrigé détaillé avec barème indicatif (sur 20 : thèse/problématique 3 pts ; explication 4 pts ; analyse des arguments 4 pts ; discussion 5 pts ; introduction/conclusion 4 pts).}

\begin{enumerate}
\item \textbf{Thèse et problématique.}\\
\textbf{Thèse (une phrase) :} la technique est un pouvoir qui, loin d'assurer automatiquement la maîtrise du monde, rend dépendant celui qui la consomme sans la comprendre, si bien que seule l'appropriation de ses principes — notamment par l'éducation technique — transforme l'usager en acteur et constitue une véritable école de liberté.\\
\textbf{Problématique :} Dans quelle mesure l'éducation technique peut-elle transformer l'usager dépendant de la technique en acteur libre et maître de ses outils ?

\item \textbf{Explication de la distinction \og usager \fg / \og acteur \fg (méthode D.A.I.C.).}\\
\textbf{Définir :} l'\og usager \fg est celui qui consomme la technique sans en comprendre les principes ni savoir la produire : il utilise, mais dépend. L'\og acteur \fg est celui qui s'est approprié les principes de la technique : il comprend, répare, adapte et produit. \textbf{Analyser :} la distinction repose sur le rapport au savoir : l'usager possède l'objet mais pas la connaissance ; l'acteur possède la connaissance, donc le pouvoir. L'auteur en tire une opposition morale et politique : \og un peuple qui consomme des techniques sans les comprendre ni les produire reste mineur ; un peuple qui s'en approprie les principes devient adulte \fg — la métaphore mineur/adulte renvoie à l'émancipation. \textbf{Illustrer :} l'élève qui ne sait plus chercher sans l'écran est un usager ; l'élève de Terminale technique qui diagnostique et répare une panne en atelier devient acteur. \textbf{Citer :} \og l'éducation technique, bien comprise, est ainsi une école de liberté : elle transforme l'usager en acteur \fg.

\item \textbf{Analyse des arguments et des exemples.}\\
\textbf{Argument 1 (par les conséquences) :} la promesse de maîtrise se retourne en dépendance. \textbf{Exemples mobilisés par l'auteur :} le paysan qui \og ne sait plus semer sans engrais importés \fg ; l'artisan qui \og ne sait plus réparer sans pièces venues d'ailleurs \fg ; l'élève qui \og ne sait plus chercher sans l'écran \fg. Ces trois exemples couvrent trois secteurs (agriculture, artisanat, éducation) pour montrer l'universalité du phénomène.\\
\textbf{Argument 2 (par distinction conceptuelle) :} \og la technique n'est pas mauvaise en elle-même : elle est un pouvoir, et tout pouvoir exige d'être gouverné \fg. L'auteur refuse l'alternative naïve accepter/refuser : la vraie question est politique — \og qui la maîtrise, et pour quelles fins \fg.\\
\textbf{Argument 3 (par analogie) :} le couple mineur/adulte transpose à la technique la notion d'émancipation : être majeur, c'est s'approprier les principes. La conclusion (\og l'éducation technique [...] est une école de liberté \fg) découle logiquement des prémisses. \emph{Point de vigilance :} relever les connecteurs (\og Pourtant \fg, \og donc \fg, \og ainsi \fg) qui marquent les étapes du raisonnement.

\item \textbf{Discussion : l'enseignement technique au Cameroun, \og école de liberté \fg ?}\\
\textbf{Analyse interne :} la thèse est forte : la dépendance technique est une réalité camerounaise (importation des machines, des pièces détachées, des intrants agricoles), et seule la formation peut la rompre.\\
\textbf{Épreuve des faits (expérience d'élève de Terminale technique attendue) :} d'un côté, l'enseignement technique camerounais réalise partiellement cette promesse : les ateliers de travaux pratiques, les stages en entreprise, les projets de fin d'année et l'insertion professionnelle des lauréats (maintenance, génie civil, électrotechnique) forment de véritables acteurs capables de diagnostiquer, réparer et produire localement. De l'autre, des limites subsistent : équipements insuffisants ou importés et non réparables localement, programmes parfois éloignés des réalités du marché local, précarité des stages. L'élève peut alors devenir un usager formé sur des machines qu'il ne pourra ni entretenir ni produire.\\
\textbf{Synthèse :} la thèse doit être \emph{nuancée} : l'enseignement technique n'est une école de liberté qu'\og bien compris \fg, c'est-à-dire s'il transmet les principes et non seulement les gestes, et s'il est adossé à un tissu industriel local. La liberté technique est un projet à réaliser, non un acquis.

\item \textbf{Introduction et conclusion (corrigés types).}\\
\textbf{Introduction :} \og L'homme moderne, disait-on au siècle des machines, est enfin devenu maître et possesseur de la nature. Le texte soumis à notre étude met cette promesse à l'épreuve : il soutient que la technique, pouvoir ambivalent, asservit celui qui la consomme sans la comprendre, et que seule l'éducation technique, en transformant l'usager en acteur, constitue une véritable école de liberté. Mais cette thèse appelle une question : dans quelle mesure l'appropriation des principes techniques suffit-elle à émanciper l'individu et le peuple ? Nous analyserons d'abord la thèse de l'auteur et ses arguments, puis nous la confronterons à la réalité de l'enseignement technique au Cameroun. \fg\\
\textbf{Conclusion :} Au total, le texte établit que la technique n'émancipe que celui qui en maîtrise les principes : la dépendance n'est pas dans l'outil mais dans l'ignorance de l'outil. Notre discussion a montré que l'enseignement technique camerounais réalise partiellement cette école de liberté, à condition de former des acteurs capables de comprendre, réparer et produire, et non de simples opérateurs. Une question demeure ouverte : à l'heure de l'intelligence artificielle et des objets connectés, l'appropriation des principes techniques reste-t-elle possible pour le seul utilisateur, ou la dépendance devient-elle irréversible ?
\end{enumerate}
\textbf{Gestion du temps rappelée :} 30 min de lecture, analyse et brouillon ; 1 h 15 de rédaction ; 15 min de relecture (orthographe, citations, annonce du plan respectée).
\end{exoresolu}

\begin{savaistu}[Repères pour l'auto-évaluation (Méta-cognition)]
Avant de rendre ta copie, coche chaque item :
\begin{itemize}
\item [ ] Thèse formulée en 1 phrase complète, distincte du thème.
\item [ ] Problématique : question ouverte, termes du texte, engage explication + discussion.
\item [ ] Plan annoncé (2 parties équilibrées) et respecté dans le développement.
\item [ ] Explication : concepts définis, arguments articulés, exemples camerounais, citations exactes.
\item [ ] Discussion : Analyse interne + Objection(s) argumentée(s) + Synthèse nuancée.
\item [ ] Conclusion : Rappel thèse + Bilan discussion + Ouverture pertinente.
\item [ ] Langage : vocabulaire philosophique précis, connecteurs logiques, orthographe soignée.
\end{itemize}
\end{savaistu}

\newpage

\coursec{Fiche bilan}

\begin{popfigure}
\begin{tikzpicture}[mindmap, grow cyclic, every node/.style={concept, text=white, font=\small}, concept color=popblue, level 1/.append style={level distance=3.2cm, sibling angle=72}, level 2/.append style={level distance=1.8cm}]
\node[concept, font=\small\bfseries] {Exercice sur texte philosophique}
  child[concept color=poporange] { node {Diagnostic \& exigences} }
  child[concept color=popgreen] { node {Problématisation} }
  child[concept color=poppurple] { node {Rédaction : intro, développement, conclusion} }
  child[concept color=poppink] { node {Questions types \& pièges Bac Tech} }
  child[concept color=popgold] { node {Simulation \& auto-évaluation} };
\end{tikzpicture}
\legende{Carte mentale de la leçon : les cinq ateliers de consolidation de l'exercice sur texte philosophique au Baccalauréat technique.}
\end{popfigure}

\begin{definition}[L'essentiel de la leçon : définitions clés]
\begin{itemize}
\item \cle{Exercice sur texte} : épreuve consistant à expliquer et à discuter un texte philosophique à partir de questions précises, en mobilisant les notions du programme.
\item \cle{Problématique} : question directrice qui dégage l'enjeu philosophique du texte, c'est-à-dire le problème que l'auteur cherche à résoudre.
\item \cle{Thèse} : position défendue par l'auteur, formulable en une phrase claire.
\item \cle{Argument} : raison avancée pour justifier la thèse ; il peut être logique, factuel ou fondé sur un exemple.
\item \cle{Concept} : notion générale et abstraite (liberté, devoir, État, technique) dont le texte précise le sens.
\end{itemize}
\end{definition}

\begin{aretenir}[Repères à connaître par c\oe ur]
\begin{center}
\begin{tabularx}{\linewidth}{|l|X|}\hline
\rowcolor{popblueL} \textbf{Élément} & \textbf{Exigence officielle} \\ \hline
Introduction & Accroche brève, présentation du texte (auteur, \oe uvre), thèse, problématique, annonce du plan \\ \hline
Développement & Réponses ordonnées aux questions : idée, explication, exemple, référence à un auteur du programme \\ \hline
Conclusion & Bilan des réponses, réponse explicite à la problématique, ouverture éventuelle \\ \hline
Méthode de lecture & 1\textsuperscript{re} lecture globale, 2\textsuperscript{e} lecture analytique (surligner thèse, arguments, exemples) \\ \hline
\end{tabularx}
\end{center}
\end{aretenir}

\begin{methode}[Méthodes express]
\begin{enumerate}
\item \textbf{Repérer la thèse} : chercher les connecteurs (donc, ainsi, par conséquent) et les verbes d'affirmation.
\item \textbf{Dégager la problématique} : transformer le sujet en une question ouverte commençant par « Comment... », « Pourquoi... », « Faut-il... ».
\item \textbf{Répondre à une question} : une idée par paragraphe ; annoncer, expliquer, illustrer (exemple camerounais ou de la vie courante), conclure le paragraphe.
\item \textbf{Justifier} : toujours relier sa réponse au texte (citation courte entre guillemets) et au cours.
\end{enumerate}
\end{methode}

\begin{attention}[Pièges fréquents à l'exercice sur texte]
\begin{itemize}
\item \textbf{La paraphrase} : recopier le texte avec d'autres mots ne vaut aucun point ; il faut expliquer le sens et l'enjeu des idées.
\item \textbf{La citation trop longue} : une citation de plus de deux lignes est sanctionnée ; citer court et commenter.
\item \textbf{La confusion expliquer / opiner} : l'exercice demande d'abord de rendre compte fidèlement de la pensée de l'auteur ; l'avis personnel, s'il est demandé, vient en discussion argumentée, jamais à la place de l'explication.
\item \textbf{L'oubli d'une question} : répondre à toutes les questions posées, dans l'ordre, en annonçant chacune explicitement.
\item \textbf{Le hors-sujet} : chaque référence au cours doit éclairer le texte, et non servir d'exposé général sur la notion.
\end{itemize}
\end{attention}

\begin{exemplebox}[De la question à la problématique]
Texte étudié : un extrait où l'auteur affirme que « la technique libère l'homme du travail pénible ».
\begin{itemize}
\item \textbf{Thèse repérée} : la technique est un instrument de libération humaine.
\item \textbf{Problématique possible} : « La technique libère-t-elle toujours l'homme, ou peut-elle aussi l'asservir ? »
\item \textbf{Exemple de la vie courante} : le téléphone mobile facilite le commerce du vendeur de Douala (libération), mais peut rendre dépendant aux réseaux sociaux (asservissement).
\end{itemize}
\end{exemplebox}

\begin{aretenir}[Checklist avant le Bac]
Avant de rendre ma copie d'exercice sur texte, je vérifie :
\begin{itemize}
\item[$\square$] J'ai lu le texte deux fois et repéré la thèse de l'auteur.
\item[$\square$] J'ai formulé une problématique sous forme de question.
\item[$\square$] Mon introduction présente le texte, la thèse et annonce le plan.
\item[$\square$] Je réponds à toutes les questions posées, dans l'ordre.
\item[$\square$] Chaque paragraphe contient une idée, une explication et un exemple.
\item[$\square$] J'ai cité brièvement le texte entre guillemets pour justifier mes réponses.
\item[$\square$] J'ai mobilisé au moins une notion ou un auteur du programme.
\item[$\square$] Je distingue bien expliquer le texte et donner mon avis personnel.
\item[$\square$] Ma conclusion répond clairement à la problématique.
\item[$\square$] Je n'ai pas paraphrasé le texte ni recopié de longs passages.
\item[$\square$] Mon écriture est lisible, sans abréviations, avec une orthographe soignée.
\item[$\square$] J'ai géré mon temps : relecture finale d'au moins cinq minutes.
\end{itemize}
\end{aretenir}

\findecours

\end{document}
